% THE GEOLOGY OF CAMEROON — Geology Upper Sixth — Cours signé OnBuch+
% Official MINESEC syllabus. Compiler avec : tectonic GEO09-the-geology-of-cameroon.tex
\newcommand{\DOCMATIERE}{Geology}
\newcommand{\DOCNIVEAU}{Upper Sixth}
\newcommand{\DOCMODULE}{Upper Sixth — Geology}
\newcommand{\DOCLECON}{9}
\newcommand{\DOCTITRE}{THE GEOLOGY OF CAMEROON}
% OnBuch+ — préambule « cours signés » ENRICHI (compilé avec tectonic / XeLaTeX)
% Système visuel « Pop » d'OnBuch+ : fond crème (jamais blanc), bordures encre
% épaisses, ombres « dures » décalées, titres Archivo Black en capitales.
% Version MATHÉMATIQUES : graphes (pgfplots/tikz), boîtes pédagogiques
% (définition, propriété, méthode, attention, le savais-tu, exercice résolu,
% exercice, TP, bilan), page de garde et signature OnBuch+ ancrée en bas.
%
% Macros attendues dans main.tex AVANT \input{preamble} :
%   \DOCMATIERE  \DOCNIVEAU  \DOCMODULE  \DOCLECON (numéro)  \DOCTITRE
\documentclass[11pt]{article}

\usepackage{fontspec}
\usepackage[english]{babel}
\usepackage[a4paper,margin=2cm,top=3.4cm,bottom=2.8cm,headheight=1.4cm,footskip=1.3cm]{geometry}
\usepackage{amsmath,amssymb,amsfonts}
\usepackage{mathtools} % \xmapsto, \coloneqq... souvent utilisés par les modèles
\usepackage{cancel} % simplifications barrées (bilans de forces)
% \abs{z} (module, valeur absolue) et \norm{u} (norme), utilisés par les modèles
\providecommand{\abs}{}\let\abs\relax\DeclarePairedDelimiter{\abs}{\lvert}{\rvert}
\providecommand{\norm}{}\let\norm\relax\DeclarePairedDelimiter{\norm}{\lVert}{\rVert}
\usepackage[table]{xcolor} % [table] : fournit aussi \rowcolors (alternance de lignes)
\usepackage{pagecolor}
\usepackage{tikz}
\usetikzlibrary{arrows.meta,positioning,calc,shapes.geometric,shapes.misc,shapes.arrows,decorations.pathmorphing,decorations.pathreplacing,decorations.markings,patterns,shadows,mindmap,trees,fit,backgrounds,matrix,intersections,angles,quotes}
\usepackage{pgfplots}
\usepackage[version=4]{mhchem} % équations nucléaires \ce{^{235}_{92}U}
\usepackage[european,siunitx]{circuitikz} % schémas électriques
\pgfplotsset{compat=1.18}
\usepgfplotslibrary{fillbetween,groupplots} % aires entre courbes (calcul intégral)
\usepackage{enumitem}
\usepackage{fancyhdr}
% [most] charge toutes les bibliothèques tcolorbox (listings, minted, raster,
% documentation...) et gonfle inutilement la mémoire TeX sur un document long
% (~300 boîtes) : on ne charge que ce qui est réellement utilisé.
\usepackage[skins,breakable]{tcolorbox}
\usepackage{graphicx}
\usepackage{colortbl}
\usepackage{booktabs}
\usepackage{tabularx}
\usepackage{diagbox} % case d'en-tête diagonale des tableaux à double entrée
% Colonnes extensibles centrée / gauche / droite (C, L, R), souvent utilisées par les modèles
\newcolumntype{C}{>{\centering\arraybackslash}X}
\newcolumntype{L}{>{\raggedright\arraybackslash}X}
\newcolumntype{R}{>{\raggedleft\arraybackslash}X}
\usepackage{array}
\usepackage{multicol}
\usepackage{siunitx}
% parse-numbers=false : accepte aussi \SI{2,5 \times 10^{-3}}{...}
\sisetup{locale=UK,output-decimal-marker={.},group-minimum-digits=4,parse-numbers=false}
\DeclareSIUnit\ohm{\ensuremath{\symup{\Omega}}} % U+2126 absent de la police
\DeclareSIUnit\division{div} % réglages d'oscilloscope (ms/div, V/div)
\DeclareSIUnit\tr{tr} % tours (vitesse de rotation, ex. \SI{1500}{\tr\per\minute})
\usepackage{qrcode}
% \degree : commande fréquemment utilisée par les modèles pour un angle ou une
% température en dehors de \SI{}{} (non fournie par les paquets chargés).
\providecommand{\degree}{^{\circ}}
% \qed (fin de démonstration), utilisé par les modèles sans amsthm
\providecommand{\qed}{\hfill$\square$}
% \barre : double barre des valeurs interdites dans un tableau de variations (tkz-tab)
\providecommand{\barre}{\ensuremath{\|}}
% environnement proof (amsthm non chargé) utilisé par les modèles
\ifdefined\proof\else\newenvironment{proof}[1][Proof]{\par\noindent\textit{#1.}\ }{\qed\par}\fi
\usepackage{lastpage}
\usepackage{hyperref}
\hypersetup{hidelinks}

% ============================================================
% Palette « Pop » OnBuch+ — mêmes hex que la classe P (plus_ui.dart)
% ============================================================
\definecolor{poporange}{HTML}{F59321}
\definecolor{popdark}{HTML}{E07A0C}
\definecolor{popcream}{HTML}{FFF6E8}
\definecolor{popink}{HTML}{1C1714}
\definecolor{poppurple}{HTML}{7A5AE0}
\definecolor{popgreen}{HTML}{2E9E6A}
\definecolor{poppink}{HTML}{E0457B}
\definecolor{popblue}{HTML}{2F7DE1}
\definecolor{popgold}{HTML}{C98A12}
\definecolor{popmuted}{HTML}{8A7F76}
\definecolor{popline}{HTML}{EADFD2}
\definecolor{popgray}{HTML}{8A7F76}
\colorlet{popgrey}{popgray}
% Alias tolérants (le modèle nomme parfois les couleurs différemment)
\colorlet{popwhite}{white}
\colorlet{popblack}{popink}
\colorlet{popred}{poppink}
\colorlet{popyellow}{popgold}
% Teintes claires pour fonds de tableaux / zones
\colorlet{popblueL}{popblue!10!white}
\colorlet{poporangeL}{poporange!14!white}
\colorlet{popgreenL}{popgreen!12!white}
\colorlet{poppurpleL}{poppurple!12!white}
\colorlet{poppinkL}{poppink!12!white}
\colorlet{popgoldL}{popgold!14!white}

\pagecolor{popcream}

% ============================================================
% Typographie
% ============================================================
\defaultfontfeatures{Ligatures=TeX}
\setmainfont{PlusJakartaSans-Medium.ttf}[Path=../../../fonts/,BoldFont=PlusJakartaSans-Bold.ttf]
% Maths en sans-serif (Fira Math) avec chiffres et lettres droites de Plus
% Jakarta, pour que les formules se fondent dans le texte.
\usepackage[math-style=ISO,bold-style=ISO]{unicode-math}
\setmathfont{FiraMath-Regular.otf}[Path=../../../fonts/]
\setmathfont{PlusJakartaSans-Medium.ttf}[Path=../../../fonts/,range={up/num,up/{Latin,latin},\mathup}]
% (icomma retiré : décimales avec point en anglais)
% Caractères mathématiques Unicode absents des polices de texte (Plus Jakarta,
% Archivo Black) : rendus par leur équivalent mathématique, en texte comme en
% formule — sinon ils s'affichent en carré vide (ex. « Arithmétique dans ℤ »).
\usepackage{newunicodechar}
\newunicodechar{ℝ}{\ensuremath{\mathbb{R}}}
\newunicodechar{ℤ}{\ensuremath{\mathbb{Z}}}
\newunicodechar{ℂ}{\ensuremath{\mathbb{C}}}
\newunicodechar{ℕ}{\ensuremath{\mathbb{N}}}
\newunicodechar{ℚ}{\ensuremath{\mathbb{Q}}}
\newunicodechar{𝔻}{\ensuremath{\mathbb{D}}}
\newunicodechar{α}{\ensuremath{\alpha}}
\newunicodechar{β}{\ensuremath{\beta}}
\newunicodechar{γ}{\ensuremath{\gamma}}
\newunicodechar{δ}{\ensuremath{\delta}}
\newunicodechar{ε}{\ensuremath{\varepsilon}}
\newunicodechar{θ}{\ensuremath{\theta}}
\newunicodechar{λ}{\ensuremath{\lambda}}
\newunicodechar{μ}{\ensuremath{\mu}}
\newunicodechar{σ}{\ensuremath{\sigma}}
\newunicodechar{τ}{\ensuremath{\tau}}
\newunicodechar{φ}{\ensuremath{\varphi}}
\newunicodechar{ψ}{\ensuremath{\psi}}
\newunicodechar{ω}{\ensuremath{\omega}}
\newunicodechar{Σ}{\ensuremath{\Sigma}}
% majuscules produites par \MakeUppercase dans les titres (α-aminés -> Α)
\newunicodechar{Α}{\ensuremath{\alpha}}
\newunicodechar{Β}{\ensuremath{\beta}}
\newunicodechar{Γ}{\ensuremath{\Gamma}}
\newunicodechar{Δ}{\ensuremath{\Delta}}
\newunicodechar{Ε}{\ensuremath{\varepsilon}}
\newunicodechar{Θ}{\ensuremath{\Theta}}
\newunicodechar{Λ}{\ensuremath{\Lambda}}
\newunicodechar{Μ}{\ensuremath{\mu}}
\newunicodechar{Τ}{\ensuremath{\tau}}
\newunicodechar{Φ}{\ensuremath{\Phi}}
\newunicodechar{Ψ}{\ensuremath{\Psi}}
\newunicodechar{Ω}{\ensuremath{\Omega}}
\newunicodechar{↦}{\ensuremath{\mapsto}}
\newunicodechar{∈}{\ensuremath{\in}}
\newunicodechar{∉}{\ensuremath{\notin}}
\newunicodechar{∧}{\ensuremath{\wedge}}
\newunicodechar{⇔}{\ensuremath{\Leftrightarrow}}
\newunicodechar{⟺}{\ensuremath{\Longleftrightarrow}}
\newunicodechar{⇒}{\ensuremath{\Rightarrow}}
\newunicodechar{⟹}{\ensuremath{\Longrightarrow}}
\newunicodechar{∘}{\ensuremath{\circ}}
\newunicodechar{∪}{\ensuremath{\cup}}
\newunicodechar{∩}{\ensuremath{\cap}}
\newunicodechar{≡}{\ensuremath{\equiv}}
\newunicodechar{⊂}{\ensuremath{\subset}}
\newunicodechar{⊥}{\ensuremath{\perp}}
\newunicodechar{∃}{\ensuremath{\exists}}
\newunicodechar{∀}{\ensuremath{\forall}}
\newunicodechar{∛}{\ensuremath{\sqrt[3]{\,}}}
\newunicodechar{∜}{\ensuremath{\sqrt[4]{\,}}}
\newunicodechar{⃗}{\ensuremath{{}^{\to}}}
\newunicodechar{ⁿ}{\ifmmode^{n}\else\textsuperscript{\ensuremath{n}}\fi}
\newunicodechar{ˣ}{\ifmmode^{x}\else\textsuperscript{\ensuremath{x}}\fi}
\newunicodechar{ᵇ}{\ifmmode^{b}\else\textsuperscript{\ensuremath{b}}\fi}
\newunicodechar{ʳ}{\ifmmode^{r}\else\textsuperscript{\ensuremath{r}}\fi}
\newunicodechar{ᵘ}{\ifmmode^{u}\else\textsuperscript{\ensuremath{u}}\fi}
\newunicodechar{ʸ}{\ifmmode^{y}\else\textsuperscript{\ensuremath{y}}\fi}
\newunicodechar{ᵃ}{\ifmmode^{a}\else\textsuperscript{\ensuremath{a}}\fi}
\newunicodechar{ᵉ}{\ifmmode^{e}\else\textsuperscript{\ensuremath{e}}\fi}
\newunicodechar{ᵏ}{\ifmmode^{k}\else\textsuperscript{\ensuremath{k}}\fi}
\newunicodechar{ᵖ}{\ifmmode^{p}\else\textsuperscript{\ensuremath{p}}\fi}
\newunicodechar{ᴺ}{\ifmmode^{N}\else\textsuperscript{\ensuremath{N}}\fi}
\newunicodechar{ᶜ}{\ifmmode^{c}\else\textsuperscript{\ensuremath{c}}\fi}
\newunicodechar{ʰ}{\ifmmode^{h}\else\textsuperscript{\ensuremath{h}}\fi}
\newunicodechar{ᵠ}{\ifmmode^{q}\else\textsuperscript{\ensuremath{q}}\fi}
\newunicodechar{ₙ}{\ifmmode_{n}\else\textsubscript{\ensuremath{n}}\fi}
\newunicodechar{ₐ}{\ifmmode_{a}\else\textsubscript{\ensuremath{a}}\fi}
\newunicodechar{ᵢ}{\ifmmode_{i}\else\textsubscript{\ensuremath{i}}\fi}
\newunicodechar{ᵣ}{\ifmmode_{r}\else\textsubscript{\ensuremath{r}}\fi}
\newunicodechar{ₖ}{\ifmmode_{k}\else\textsubscript{\ensuremath{k}}\fi}
\newunicodechar{ᵦ}{\ifmmode_{\beta}\else\textsubscript{\ensuremath{\beta}}\fi}
\newunicodechar{ₓ}{\ifmmode_{x}\else\textsubscript{\ensuremath{x}}\fi}
\newfontfamily\popdisplay{ArchivoBlack-Regular.ttf}[Path=../../../fonts/]
\newfontfamily\poplogo{SpaceGrotesk-Bold.ttf}[Path=../../../fonts/]
\newfontfamily\popsemi{PlusJakartaSans-SemiBold.ttf}[Path=../../../fonts/]
\newfontfamily\popxbold{PlusJakartaSans-ExtraBold.ttf}[Path=../../../fonts/]
\newfontfamily\popxboldls{PlusJakartaSans-ExtraBold.ttf}[Path=../../../fonts/,LetterSpace=3.0]

\setlist[enumerate]{leftmargin=1.7em,itemsep=5pt,topsep=4pt}
\setlist[itemize]{leftmargin=1.7em,itemsep=4pt,topsep=4pt,label={\color{poporange}\textbullet}}
\setlength{\parindent}{0pt}
\setlength{\parskip}{5pt}
\renewcommand{\arraystretch}{1.25}

% pgfplots : style Pop par défaut
\pgfplotsset{
  every axis/.append style={axis line style={popink,line width=1pt},tick style={popink},
    grid style={popline,line width=0.5pt},label style={font=\small},tick label style={font=\footnotesize}},
  popcurve/.style={poporange,line width=1.8pt,no markers,smooth},
}
% Même style disponible dans un \draw TikZ simple (hors pgfplots)
\tikzset{popcurve/.style={poporange,line width=1.8pt}}
\tikzset{popgrid/.style={popline,very thin}}
\colorlet{popcurve}{poporange} % le nom du style est parfois utilisé comme couleur

% ============================================================
% Pill / squiggle
% ============================================================
\newcommand{\poppill}[3]{%
  \tikz[baseline=-0.55ex]\node[draw=popink,line width=1.2pt,rounded corners=2.6pt,%
    fill=#2,inner xsep=6.5pt,inner ysep=3pt]{%
    \popxboldls\fontsize{7}{7}\selectfont\color{#3}\MakeUppercase{#1}};%
}
\newcommand{\popsquiggle}[1]{%
  \tikz[baseline=-0.4ex]{
    \draw[#1,line width=2.6pt,line cap=round]
      plot [smooth] coordinates {(0,0.09) (0.34,0.01) (0.68,0.21) (1.02,0.01) (1.36,0.10)};
  }%
}
% Mot-clé surligné (marqueur orange sous le texte)
\newcommand{\cle}[1]{{\popxbold\color{popdark}#1}}

% ============================================================
% En-tête / pied
% ============================================================
\pagestyle{fancy}
\fancyhf{}
\renewcommand{\headrulewidth}{0pt}
\renewcommand{\footrulewidth}{0pt}
\lhead{\raisebox{-0.15ex}{\poplogo\fontsize{13}{13}\selectfont\color{popink}OnBuch\color{poporange}+}}
\rhead{\poppill{\DOCMATIERE\ \textbullet\ \DOCNIVEAU\ \textbullet\ Chapter \DOCLECON}{white}{popink}}
\lfoot{\popxbold\fontsize{7.6}{7.6}\selectfont\color{popmuted}\MakeUppercase{Signed course OnBuch+}\ \textbullet\ {\popsemi\DOCTITRE}}
\rfoot{\tikz[baseline=-0.6ex]\node[rounded corners=8.5pt,fill=popink,minimum height=17pt,minimum width=17pt,inner xsep=5pt,inner sep=0]{\popxbold\fontsize{8}{8}\selectfont\color{popcream}\thepage\,/\,\pageref*{LastPage}};}

% ============================================================
% Titres
% ============================================================
\newcommand{\chaptitre}[2]{%
  \begin{tcolorbox}[enhanced,colback=poporange,colframe=popink,boxrule=2.6pt,arc=11pt,
      drop shadow=popink,left=15pt,right=15pt,top=12pt,bottom=11pt,before skip=3pt,after skip=18pt]
    \poppill{#2}{popink}{popcream}\par\vspace{9pt}
    {\raggedright\hyphenpenalty=10000\exhyphenpenalty=10000\popdisplay\fontsize{22}{24}\selectfont\color{popcream}\MakeUppercase{#1}\par}
  \end{tcolorbox}
}
% Section numérotée : pastille orange avec le numéro + titre Archivo
\newcounter{popsec}
\newcommand{\coursec}[1]{%
  \stepcounter{popsec}\setcounter{popsub}{0}%
  \par\vspace{16pt}\noindent
  \begin{minipage}{\linewidth}
  \tikz[baseline=-0.7ex]\node[draw=popink,line width=1.6pt,fill=poporange,rounded corners=4pt,minimum size=22pt,inner sep=0]{\popdisplay\fontsize{12}{12}\selectfont\color{popcream}\thepopsec};%
  \hspace{8pt}{\popdisplay\fontsize{15}{17}\selectfont\color{popink}\MakeUppercase{#1}}\par
  \vspace{4pt}\noindent{\color{poporange}\rule{36pt}{3.6pt}}
  \end{minipage}\par\nopagebreak\vspace{8pt}
}
\newcounter{popsub}
\newcommand{\courssub}[1]{%
  \stepcounter{popsub}%
  \par\vspace{9pt}\noindent{\popxbold\fontsize{11.5}{13}\selectfont\color{popink}{\color{poporange}\thepopsec.\thepopsub}\hspace{6pt}#1}\par\nopagebreak\vspace{4pt}
}

% ============================================================
% Boîtes pédagogiques — même recette : fond blanc, bordure épaisse colorée,
% ombre dure encre, badge plein. Titre optionnel [#1].
% ============================================================
\tcbset{popbase/.style={breakable,enhanced,colback=white,boxrule=2.3pt,arc=9pt,
  shadow={3pt}{-3pt}{0pt}{popink},left=14pt,right=14pt,top=11pt,bottom=11pt,before skip=11pt,after skip=15pt,
  fontupper=\popsemi\fontsize{10.6}{14.5}\selectfont\color{popink},
  fontlower=\popsemi\fontsize{10.6}{14.5}\selectfont\color{popink}}}
% Coche dessinée (le glyphe ✓ n'existe pas dans les polices du document)
\newcommand{\popcheck}[1]{\tikz[baseline=-0.5ex]\draw[#1,line width=1.6pt,line cap=round,line join=round](-0.13,0.0)--(-0.04,-0.09)--(0.13,0.1);}
\providecommand{\coche}{\popcheck{popgreen}}
\providecommand{\resultat}[1]{\par\smallskip\noindent\fcolorbox{popgreen}{popgreenL}{\parbox{\dimexpr\linewidth-2\fboxsep-2\fboxrule\relax}{\textbf{Result.} #1}}\par\smallskip}
\newcommand{\popboxhead}[3]{\poppill{#1}{#2}{white}\ifx\relax#3\relax\else\hspace{6pt}{\popxbold\fontsize{10.6}{12}\selectfont\color{popink}#3}\fi\par\vspace{7pt}}

\newtcolorbox{aretenir}[1][]{popbase,colframe=popgreen,before upper={\popboxhead{Key points}{popgreen}{#1}}}
\newtcolorbox{exemplebox}[1][]{popbase,colframe=poporange,before upper={\popboxhead{Exemple}{poporange}{#1}}}
\newtcolorbox{definition}[1][]{popbase,colframe=popblue,before upper={\popboxhead{Definition}{popblue}{#1}}}
\newtcolorbox{propriete}[1][]{popbase,colframe=poppurple,before upper={\popboxhead{Property}{poppurple}{#1}}}
\newtcolorbox{methode}[1][]{popbase,colframe=popgold,before upper={\popboxhead{Method}{popgold}{#1}}}
\newtcolorbox{attention}[1][]{popbase,colframe=poppink,before upper={\popboxhead{Warning}{poppink}{#1}}}
\newtcolorbox{experience}[1][]{popbase,colframe=popblue,colback=popblueL,before upper={\popboxhead{Experiment}{popblue}{#1}}}
% « Le savais-tu ? » : carte violette pleine (contexte, culture, Cameroun)
\newtcolorbox{savaistu}[1][]{popbase,colback=poppurple,colframe=popink,
  fontupper=\popsemi\fontsize{10.4}{14.2}\selectfont\color{white},
  before upper={\poppill{Did you know?}{white}{poppurple}\ifx\relax#1\relax\else\hspace{6pt}{\popxbold\color{white}#1}\fi\par\vspace{7pt}}}
% Exercice résolu : énoncé en haut, \tcblower puis la solution sur fond crème
\newcounter{popexo}\newcounter{popexr}
\newtcolorbox{exoresolu}[1][]{popbase,colframe=poporange,colbacklower=popcream,
  segmentation style={popink,line width=1.2pt,dashed},
  before upper={\stepcounter{popexr}\popboxhead{Worked example \thepopexr}{poporange}{#1}},
  before lower={\poppill{Solution}{popink}{popcream}\par\vspace{6pt}}}
% Exercice (non résolu en place) — la correction est dans la partie Corrigés
\newtcolorbox{exercice}[1][]{popbase,colframe=popink,
  before upper={\stepcounter{popexo}\popboxhead{Exercise \thepopexo}{popink}{#1}}}
% Corrigé
\newtcolorbox{corrige}[1][]{popbase,colframe=popgreen,colback=popgreenL,
  before upper={\popboxhead{Solution}{popgreen}{#1}}}
% Objectifs / prérequis / bilan
\newtcolorbox{objectifs}{popbase,colframe=popink,colback=poporangeL,
  before upper={\popboxhead{Chapter objectives}{popink}{}}}
\newtcolorbox{prerequis}{popbase,colframe=popink,colback=popblueL,
  before upper={\popboxhead{Prerequisites}{popblue}{}}}
\newtcolorbox{bilan}{popbase,colframe=popink,colback=white,boxrule=2.8pt,
  before upper={\popboxhead{Fiche bilan}{poporange}{L'essentiel à connaître pour l'examen}}}
% Figure encadrée avec légende
\newtcolorbox{popfigure}[1][]{enhanced,breakable=false,colback=white,colframe=popline,boxrule=1.4pt,arc=8pt,
  left=10pt,right=10pt,top=10pt,bottom=8pt,before skip=10pt,after skip=12pt,halign=center,
  lower separated=false,halign lower=center,
  fontlower=\popsemi\fontsize{9}{11.5}\selectfont\color{popmuted},
  #1}
\newcommand{\legende}[1]{\par\vspace{6pt}{\popsemi\fontsize{9}{11.5}\selectfont\color{popmuted}#1}}

% ============================================================
% Signature OnBuch+
% ============================================================
\newcommand{\poplogoinline}[1]{{\poplogo\fontsize{#1}{#1}\selectfont\color{popink}OnBuch\color{poporange}+}}
\newcommand{\signatureonbuch}{%
  \begin{tcolorbox}[enhanced,colback=popink,colframe=popink,boxrule=2.2pt,arc=10pt,
      drop shadow=poporange,left=14pt,right=14pt,top=10pt,bottom=10pt,before skip=0pt,after skip=0pt]
    \tikz[baseline=-0.7ex]\node[circle,fill=popgreen,draw=popcream,line width=1.3pt,minimum size=22pt,inner sep=0]{\popcheck{white}};%
    \hspace{9pt}\begin{minipage}[c]{0.62\linewidth}
      {\popxbold\fontsize{12}{13}\selectfont\color{popcream}Signed\ \textbullet\ {\poplogo OnBuch\color{poporange}+}}\par\vspace{2pt}
      {\raggedright\popsemi\fontsize{8}{10.5}\selectfont\color{popline}\MakeUppercase{OnBuch+ teaching team}\\\MakeUppercase{Follows the official MINESEC syllabus}\par}
    \end{minipage}\hfill
    {\poplogo\fontsize{22}{22}\selectfont\color{popcream}OnBuch\color{poporange}+}
  \end{tcolorbox}%
}

% ============================================================
% Page de garde
% #1 titre · #2 matière · #3 niveau · #4 module · #5 n° leçon · #6 durée ·
% #7 description / objectifs (texte ou itemize)
% ============================================================
\newcommand{\pagegarde}[7]{%
  \thispagestyle{empty}
  \newgeometry{a4paper,margin=1.7cm,top=1.6cm,bottom=1.4cm}%
  \begin{tikzpicture}[remember picture,overlay]
    \fill[poporange!18] ([xshift=-3.2cm,yshift=-3.6cm]current page.north east) circle (4.6cm);
    \fill[poppurple!14] ([xshift=2.4cm,yshift=7.6cm]current page.south west) circle (3.8cm);
  \end{tikzpicture}%
  {\poplogo\fontsize{30}{30}\selectfont\color{popink}OnBuch\color{poporange}+}\hfill
  \raisebox{0.6ex}{\poppill{Signed course \textbullet\ Premium}{popink}{popcream}}\par
  \vspace{1pt}\popsquiggle{poppurple}\par
  \vspace{8pt}
  \begin{tcolorbox}[enhanced,colback=poporange,colframe=popink,boxrule=2.8pt,arc=13pt,
      drop shadow=popink,left=18pt,right=18pt,top=12pt,bottom=13pt,after skip=10pt]
    \poppill{#2 \textbullet\ #3}{popink}{popcream}\hspace{5pt}\poppill{#4}{white}{popink}\par\vspace{8pt}
    {\popdisplay\fontsize{14}{14}\selectfont\color{popink}CHAPTER #5}\par\vspace{4pt}
    {\raggedright\hyphenpenalty=10000\exhyphenpenalty=10000\popdisplay\fontsize{25}{27}\selectfont\color{popcream}\MakeUppercase{#1}\par}
    \vspace{8pt}
    {\popxbold\fontsize{9.5}{9.5}\selectfont\color{popink}\MakeUppercase{Suggested time: #6}}
  \end{tcolorbox}
  \begin{tcolorbox}[enhanced,colback=white,colframe=popink,boxrule=2.2pt,arc=10pt,
      drop shadow=popink,left=16pt,right=16pt,top=10pt,bottom=10pt,after skip=8pt]
    \poppill{In this chapter}{poppurple}{white}\par\vspace{5pt}
    \popsemi\fontsize{9.3}{12.6}\selectfont\color{popink}#7
  \end{tcolorbox}
  \noindent\begin{minipage}[t]{0.487\linewidth}
    \begin{tcolorbox}[enhanced,colback=popgreen,colframe=popink,boxrule=2.2pt,arc=10pt,
        drop shadow=popink,left=11pt,right=11pt,top=10pt,bottom=10pt,height=3.1cm,valign=center]
      \tcbox[colback=white,colframe=popink,boxrule=1.3pt,arc=5pt,boxsep=0pt,nobeforeafter,
        left=3pt,right=3pt,top=3pt,bottom=3pt,valign=center]{\qrcode[height=1.9cm]{https://play.google.com/store/apps/details?id=com.luvvix.onbuch&hl=en}}%
      \hspace{8pt}\begin{minipage}[c]{0.5\linewidth}
        {\popdisplay\fontsize{11}{12.5}\selectfont\color{white}\MakeUppercase{Download OnBuch}}\par\vspace{4pt}
        {\raggedright\popsemi\fontsize{8.4}{11}\selectfont\color{white}Free \textbullet\ Google Play\\AI tutor, past papers, results\par}
      \end{minipage}
    \end{tcolorbox}
  \end{minipage}\hfill
  \begin{minipage}[t]{0.487\linewidth}
    \begin{tcolorbox}[enhanced,colback=poppurple,colframe=popink,boxrule=2.2pt,arc=10pt,
        drop shadow=popink,left=11pt,right=11pt,top=10pt,bottom=10pt,height=3.1cm,valign=center]
      \tikz[baseline=-0.7ex]\node[circle,fill=white,draw=popink,line width=1.3pt,minimum size=1.6cm,inner sep=0]{\popdisplay\fontsize{24}{24}\selectfont\color{poppurple}+};%
      \hspace{8pt}\begin{minipage}[c]{0.6\linewidth}
        {\popdisplay\fontsize{11}{12.5}\selectfont\color{white}\MakeUppercase{Go OnBuch+}}\par\vspace{4pt}
        {\raggedright\popsemi\fontsize{8.4}{11}\selectfont\color{white}All signed courses, unlimited AI tutor, PDF sheets — OnBuch+ tab in the app\par}
      \end{minipage}
    \end{tcolorbox}
  \end{minipage}
  \vspace{8pt}
  \popcontenu
  \vfill
  \signatureonbuch
  \restoregeometry
  \newpage
}

% Rangée « Ce cours contient » (page de garde)
\newcommand{\poptile}[3]{% #1 couleur #2 gros texte #3 légende
  \begin{tcolorbox}[enhanced,colback=#1,colframe=popink,boxrule=2pt,arc=9pt,shadow={2.5pt}{-2.5pt}{0pt}{popink},
      left=6pt,right=6pt,top=9pt,bottom=9pt,halign=center,valign=center,height=1.75cm,width=\linewidth,nobeforeafter]
    {\popdisplay\fontsize{12.5}{13}\selectfont\color{white}\MakeUppercase{#2}}\par\vspace{3pt}
    {\centering\popsemi\fontsize{7.8}{9.5}\selectfont\color{white}#3\par}
  \end{tcolorbox}}
\newcommand{\popcontenu}{%
  \par\noindent{\popxboldls\fontsize{7.5}{7.5}\selectfont\color{popmuted}THIS SIGNED COURSE CONTAINS}\par\vspace{7pt}
  \noindent\begin{minipage}[t]{0.235\linewidth}\poptile{popblue}{Course}{complete, illustrated, step by step}\end{minipage}\hfill
  \begin{minipage}[t]{0.235\linewidth}\poptile{poporange}{Methods}{and worked examples}\end{minipage}\hfill
  \begin{minipage}[t]{0.235\linewidth}\poptile{poppink}{Exercises}{exam style + solutions}\end{minipage}\hfill
  \begin{minipage}[t]{0.235\linewidth}\poptile{popgreen}{Summary sheet}{the essentials to remember}\end{minipage}\par}

% Sommaire
\newtcolorbox{sommaire}{enhanced,colback=white,colframe=popink,boxrule=2.2pt,arc=10pt,
  drop shadow=popink,left=18pt,right=18pt,top=13pt,bottom=13pt,before skip=6pt,after skip=20pt,
  before upper={\poppill{Contents}{poppurple}{white}\par\vspace{9pt}\popsemi\fontsize{10.6}{14.5}\selectfont\color{popink}}}

% Signature de fin de cours — poussée tout en bas de la dernière page
\newcommand{\findecours}{%
  \par\vfill\nopagebreak
  \begin{center}\popsquiggle{poporange}\end{center}\vspace{4pt}
  \signatureonbuch
}
\AtBeginDocument{\def\llbracket{\mathopen{[\mkern-3.5mu[}}\def\rrbracket{\mathclose{]\mkern-3.5mu]}}} % intervalles d'entiers (glyphe ⟦ absent de la police)
% Glyphes absents de Fira Math : redéfinis à partir de glyphes disponibles
\AtBeginDocument{%
  \def\setminus{\mathbin{\backslash}}\let\smallsetminus\setminus
  \def\vdots{\mathord{\vbox{\baselineskip3.2pt\lineskiplimit0pt\kern4pt\hbox{.}\hbox{.}\hbox{.}}}}%
  \def\ddots{\mathinner{\mkern1mu\raise6.5pt\hbox{.}\mkern2mu\raise3.6pt\hbox{.}\mkern2mu\raise0.7pt\hbox{.}\mkern1mu}}%
  \def\triangle{\mathord{\scalebox{0.9}{$\symup{\Delta}$}}}%
  \def\nabla{\mathord{\raisebox{\depth}{\scalebox{1}[-1]{$\symup{\Delta}$}}}}%
  \def\checkmark{\popcheck{popdark}}%
  \def\star{\mathbin{\ast}}\def\bigstar{\mathord{\ast}}%
  \def\diamond{\mathbin{\scalebox{0.6}{$\Diamond$}}}%
  \def\bigcup{\mathop{\mathchoice{\raisebox{-0.3em}{\scalebox{1.7}{$\cup$}}}{\scalebox{1.25}{$\cup$}}{\cup}{\cup}}\displaylimits}%
  \def\bigcap{\mathop{\mathchoice{\raisebox{-0.3em}{\scalebox{1.7}{$\cap$}}}{\scalebox{1.25}{$\cap$}}{\cap}{\cap}}\displaylimits}%
  \def\lesseqqgtr{\mathrel{\lessgtr}}\def\bigoplus{\mathop{\oplus}\displaylimits}%
}
\providecommand{\ssi}{if and only if} % abréviation fréquente
\AtBeginDocument{\providecommand{\wideparen}{\overparen}} % arc de cercle

% Fonctions usuelles que les modèles utilisent sans que amsmath les fournisse
\providecommand{\arsinh}{\operatorname{arsinh}}\providecommand{\arcosh}{\operatorname{arcosh}}
\providecommand{\artanh}{\operatorname{artanh}}\providecommand{\arsech}{\operatorname{arsech}}
\providecommand{\arcsch}{\operatorname{arcsch}}\providecommand{\arcoth}{\operatorname{arcoth}}
\providecommand{\arcsinh}{\operatorname{arsinh}}\providecommand{\arccosh}{\operatorname{arcosh}}
\providecommand{\arctanh}{\operatorname{artanh}}\providecommand{\sech}{\operatorname{sech}}
\providecommand{\csch}{\operatorname{csch}}\providecommand{\arcsec}{\operatorname{arcsec}}
\providecommand{\arccsc}{\operatorname{arccsc}}\providecommand{\arccot}{\operatorname{arccot}}
\providecommand{\sgn}{\operatorname{sgn}}\providecommand{\lcm}{\operatorname{lcm}}
\providecommand{\Var}{\operatorname{Var}}\providecommand{\Cov}{\operatorname{Cov}}

\begin{document}

\pagegarde{THE GEOLOGY OF CAMEROON}{Geology}{Upper Sixth}{Upper Sixth — Geology}{9}{19 periods}{This chapter presents the complete geology of Cameroon: its igneous provinces and the Cameroon Volcanic Line, its metamorphic basement, its sedimentary basins, lakes, structures, mineral resources, stratigraphy and geologic hazards. It connects national geology to everyday life, economic development and GCE Advanced Level examination requirements.
\vspace{4pt}
\begin{multicols}{2}\small
\begin{itemize}
\item By the end of this chapter I can describe the main igneous rock provinces of Cameroon and their products.
\item By the end of this chapter I can explain the origin, structure and activity of the Cameroon Volcanic Line.
\item By the end of this chapter I can identify the metamorphic rocks and belts of the Cameroonian basement and relate them to past orogenic events.
\item By the end of this chapter I can describe the sedimentary basins of Cameroon, their fills and their economic importance.
\item By the end of this chapter I can explain the limnology of Cameroonian lakes, including crater lakes and the Lake Nyos gas phenomenon.
\item By the end of this chapter I can interpret the major geologic structures of Cameroon (faults, shear zones, lineaments).
\end{itemize}\end{multicols}}

\begin{sommaire}
\begin{multicols}{2}
\begin{enumerate}
\item Section 1: Igneous Geology of Cameroon and the Cameroon Volcanic Line (Lessons 129-131)
\item Section 2: Metamorphic Geology of Cameroon (Lessons 132-133)
\item Section 3: Sedimentary Geology and Basins of Cameroon (Lessons 134-136)
\item Section 4: Limnology and Lacustrine Processes in Cameroon (Lessons 137-139)
\item Section 5: Geologic Structures of Cameroon (Lessons 140-141)
\item Section 6: Economic Minerals and Natural Resources of Cameroon (Lessons 142-144)
\item Section 7: Stratigraphy of Cameroon (Lesson 145)
\item Section 8: Geologic Hazards in Cameroon and Integration (Lessons 146-147)
\item Integration activity
\item Methods and worked examples
\item Exercises
\item Solutions to the exercises
\item Summary sheet
\end{enumerate}
\end{multicols}
\end{sommaire}

\begin{objectifs}
\begin{itemize}
\item By the end of this chapter I can describe the main igneous rock provinces of Cameroon and their products.
\item By the end of this chapter I can explain the origin, structure and activity of the Cameroon Volcanic Line.
\item By the end of this chapter I can identify the metamorphic rocks and belts of the Cameroonian basement and relate them to past orogenic events.
\item By the end of this chapter I can describe the sedimentary basins of Cameroon, their fills and their economic importance.
\item By the end of this chapter I can explain the limnology of Cameroonian lakes, including crater lakes and the Lake Nyos gas phenomenon.
\item By the end of this chapter I can interpret the major geologic structures of Cameroon (faults, shear zones, lineaments).
\item By the end of this chapter I can locate and evaluate the economic minerals and natural resources of Cameroon.
\item By the end of this chapter I can construct the stratigraphic column of Cameroon from the Precambrian basement to recent deposits.
\item By the end of this chapter I can assess geologic hazards in Cameroon and propose mitigation measures.
\end{itemize}
\end{objectifs}

\begin{prerequis}
\begin{itemize}
\item Rock classification (igneous, sedimentary, metamorphic) and the rock cycle from Lower Sixth
\item Plate tectonics, volcanism and magma types (Lower Sixth)
\item Principles of stratigraphy: superposition, unconformities, relative and absolute dating
\item Metamorphic processes, grades and index minerals
\item Reading and interpreting geologic maps and cross-sections
\end{itemize}
\end{prerequis}

\newpage

\coursec{Section 1: Igneous Geology of Cameroon and the Cameroon Volcanic Line (Lessons 129-131)}

In August 1986, the communities around Lake Nyos experienced a silent catastrophe: a dense cloud of carbon dioxide released from the lake bottom killed over 1\,700 people and thousands of livestock overnight. Yet only a few hundred kilometres away, the same volcanic system gives Cameroon fertile soils, the highest peak in West Africa, and significant geothermal energy potential. Why does one geological feature bring both wealth and disaster? And why is bauxite mined in the Adamaoua while oil is pumped off the coast of Limbe? This chapter answers these questions by reading the geology of Cameroon itself, and we begin with the most spectacular part of that story: the \cle{igneous geology} of the country and its famous volcanic line.

\courssub{Overview of Igneous Activity in Cameroon}

Igneous rocks form when \cle{magma} (molten rock below the surface) or \cle{lava} (molten rock erupted at the surface) cools and solidifies. Cameroon displays both great families of igneous activity:

\begin{itemize}
\item \textbf{Plutonic (intrusive) activity}: magma solidified slowly deep underground, producing coarse-grained rocks such as \cle{granite}, \cle{syenite} and \cle{gabbro}. These rocks are exposed today only because erosion has stripped away the kilometres of rock that once covered them.
\item \textbf{Volcanic (extrusive) activity}: magma reached the surface and cooled rapidly, producing fine-grained rocks such as \cle{basalt}, \cle{trachyte} and \cle{rhyolite}, together with pyroclastic deposits (ash, lapilli, volcanic bombs).
\end{itemize}

In age, Cameroon's igneous rocks fall into two broad groups:

\begin{enumerate}
\item \textbf{Ancient intrusions} (Precambrian to Palaeozoic, roughly 3\,000 to 300 million years ago), which form part of the crystalline basement.
\item \textbf{Young volcanism} (Cenozoic, essentially the last 65 million years, and still active today), concentrated along the Cameroon Volcanic Line.
\end{enumerate}

\begin{methode}[Distinguishing plutonic from volcanic rocks by texture]
To classify a hand specimen as plutonic or volcanic, proceed as follows:
\begin{enumerate}
\item Look at the rock with the naked eye. Can you see individual mineral crystals clearly?
\item If \textbf{yes} -- crystals are large, interlocking and visible without a lens -- the rock is \textbf{phaneritic} and \textbf{plutonic} (slow cooling at depth, e.g. granite, gabbro).
\item If \textbf{no} -- the rock looks uniform, or crystals are visible only with a hand lens -- the rock is \textbf{aphanitic} and \textbf{volcanic} (fast cooling at the surface, e.g. basalt, rhyolite).
\item If large crystals (\emph{phenocrysts}) sit in a fine groundmass, the rock is \textbf{porphyritic}: it began cooling slowly at depth, then finished cooling quickly at the surface.
\end{enumerate}
\end{methode}

\begin{attention}[Texture tells you the cooling history, not the chemistry]
A common mistake is to think granite and rhyolite are different because of texture alone. In fact they have the \emph{same} chemical composition; only their cooling rate (and therefore grain size) differs. Likewise basalt and gabbro are chemical twins: basalt is volcanic, gabbro is plutonic. Never identify a rock by colour alone.
\end{attention}

\courssub{Ancient Intrusions of the Basement}

Beneath the younger cover of Cameroon lies a \cle{crystalline basement} belonging to the Congo craton in the south and the Pan-African mobile belt in the centre and north. Within this basement, several generations of ancient intrusions are recognised:

\begin{itemize}
\item \textbf{Granites and granitoids}: very widespread, especially in the southern and central parts of the country (the Ntem Complex around Ebolowa and Sangmélima, the Yaoundé Complex, and the Lom Complex further east). Some are Archaean to Palaeoproterozoic in age (older than about 2\,000 million years); others were emplaced during the \cle{Pan-African orogeny} (roughly 600 million years ago), when colliding continents melted parts of the crust.
\item \textbf{Syenites}: silica-undersaturated, feldspar-rich intrusions, often forming ring complexes (e.g. the Lomié and Poli complexes). They are important because they are commonly associated with concentrations of rare metals (Nb, Ta, REE).
\item \textbf{Gabbros and dolerites}: dark, mafic intrusions emplaced as dykes, sills and small massifs. Dolerite dyke swarms cutting the basement record ancient episodes of crustal extension, notably during the breakup of Gondwana.
\end{itemize}

These ancient intrusions matter economically: intense tropical weathering of granitic and syenitic rocks in the Adamaoua region produced the world-class \cle{bauxite} deposits of Minim-Martap and Ngaoundal, while some granitoids host tin (cassiterite) and rare-metal mineralisation (e.g. the Akonolinga area). They also matter scientifically, because they tell us that Cameroon has been a site of magma production for at least three billion years.

\begin{center}
\begin{tabularx}{\linewidth}{|l|X|X|}
\hline
\rowcolor{popblueL} \textbf{Complex / Region} & \textbf{Main Rock Types} & \textbf{Age / Tectonic Setting} \\
\hline
Ntem Complex (South) & Granites, gneisses, charnockites & Archaean--Palaeoproterozoic ($>$ 2\,000 Ma), Congo Craton \\
\hline
Yaoundé Complex (Centre) & Granitoids, migmatites & Pan-African ($\sim$ 600 Ma), collisional orogeny \\
\hline
Lom Complex (East) & Granites, syenites & Pan-African, post-collisional \\
\hline
Adamaoua Massifs (North) & Syenites, granites, gabbros & Pan-African, anorogenic ring complexes \\
\hline
\end{tabularx}
\end{center}

\courssub{The Cameroon Volcanic Line: Definition and General Setting}

\begin{definition}[The Cameroon Volcanic Line (CVL)]
The \cle{Cameroon Volcanic Line} is a linear chain of Cenozoic volcanoes, about $1\,600\ \mathrm{km}$ long, trending roughly SW--NE from the island of Pagalu (Annobón) in the Gulf of Guinea, through the islands of São Tomé, Príncipe and Bioko, across mainland Cameroon (Mount Cameroon, Manengouba, Bambouto, Oku, the Bamenda Highlands and the Adamaoua Plateau) towards the Lake Chad basin. Its defining peculiarity is that it crosses \textbf{both oceanic and continental crust} with the same style of alkaline volcanism and \textbf{without any systematic age progression} along the line.
\end{definition}

\begin{popfigure}
\begin{tikzpicture}[scale=1.0]
% Ocean
\fill[popblueL] (-0.5,-0.5) rectangle (10.5,7.5);
% Continent (Cameroon, simplified)
\fill[popgreenL] (3.2,0.2) -- (3.6,1.2) -- (4.0,2.2) -- (4.6,3.2) -- (5.2,4.4) -- (6.2,5.6) -- (7.4,6.6) -- (9.2,7.3) -- (10.5,7.3) -- (10.5,-0.5) -- (3.2,-0.5) -- cycle;
\draw[popdark, thick] (3.2,0.2) -- (3.6,1.2) -- (4.0,2.2) -- (4.6,3.2) -- (5.2,4.4) -- (6.2,5.6) -- (7.4,6.6) -- (9.2,7.3);
% CVL trend line
\draw[poporange, very thick, dashed] (0.4,0.4) -- (9.6,6.9);
\node[poporange, font=\bfseries, rotate=38] at (5.1,2.6) {CVL};
% Oceanic islands
\fill[poppurple] (0.6,0.5) circle (0.09); \node[font=\scriptsize, anchor=west] at (0.75,0.35) {Annobón};
\fill[poppurple] (1.7,1.3) circle (0.09); \node[font=\scriptsize, anchor=west] at (1.85,1.1) {São Tomé};
\fill[poppurple] (2.5,1.9) circle (0.09); \node[font=\scriptsize, anchor=west] at (2.65,1.75) {Príncipe};
\fill[poppurple] (3.35,2.6) circle (0.09); \node[font=\scriptsize, anchor=west] at (2.2,2.9) {Bioko};
% Continental volcanoes (triangles)
\foreach \x/\y/\n in {3.9/3.1/{Mt Cameroon}, 4.5/3.7/{Manengouba}, 5.0/4.3/{Bambouto}, 5.6/4.9/{Mt Oku}, 6.3/5.5/{Bamenda}, 7.3/6.2/{Adamaoua}}{
\fill[popink] (\x,\y) -- ++(-0.13,-0.2) -- ++(0.26,0) -- cycle;
\node[font=\scriptsize, anchor=west] at (\x+0.12,\y-0.18) {\n};}
% Lake Chad
\fill[popblue] (9.3,7.0) ellipse (0.5 and 0.28);
\node[font=\scriptsize, anchor=south] at (9.3,7.3) {Lake Chad};
% Labels
\node[font=\small\itshape, popblue] at (1.6,5.8) {Gulf of Guinea};
\node[font=\small\itshape, popdark] at (8.6,2.0) {Continental sector};
\node[font=\small\itshape, popblue] at (2.2,4.6) {Oceanic sector};
% North arrow
\draw[->, thick, popdark] (0.2,6.8) -- (0.2,7.4) node[above, font=\scriptsize] {N};
\end{tikzpicture}
\legende{Simplified map of the Cameroon Volcanic Line: a SW--NE alignment of volcanic centres crossing from oceanic crust (islands) onto continental crust (mainland massifs) up to the Lake Chad basin.}
\end{popfigure}

\courssub{The Oceanic Sector of the CVL}

The oceanic sector comprises four main islands rising from the floor of the Gulf of Guinea:

\begin{itemize}
\item \textbf{Annobón (Pagalu)}: the southwesternmost and smallest centre, a young shield volcano.
\item \textbf{São Tomé}: a large island with dramatic phonolite plugs (the famous ``Great Dog Peak'', Pico Cão Grande, is a volcanic neck of phonolite).
\item \textbf{Príncipe}: an older, deeply eroded island, dominated by phonolitic and trachytic rocks.
\item \textbf{Bioko} (formerly Fernando Pó): the island closest to the mainland, built of three large shield volcanoes (Santa Isabel, Pico Basile, and Gran Caldera).
\end{itemize}

The lavas of these islands are dominantly \cle{alkaline basalts} and their differentiates (trachytes, phonolites). Because the islands sit on oceanic lithosphere, they prove that the CVL is not controlled by the nature of the crust it crosses.

\begin{definition}[Alkaline basalt]
An \cle{alkaline basalt} is a silica-undersaturated lava rich in sodium (Na) and potassium (K), typically containing minerals such as nepheline or alkali feldspar in its normative composition. It is the characteristic product of \cle{intraplate volcanism}, generated by small degrees of partial melting of the mantle at relatively great depth, far from plate boundaries.
\end{definition}

\courssub{The Continental Sector of the CVL}

On the mainland, the CVL builds some of the most imposing relief in Africa:

\begin{itemize}
\item \textbf{Mount Cameroon} (about $4\,095\ \mathrm{m}$): the highest peak in West Africa and an \textbf{active volcano}, a great basaltic shield rising directly from the Atlantic coast near Buea and Limbe.
\item \textbf{Mount Manengouba}: a stratovolcano crowned by two nested calderas containing crater lakes (the ``Male'' and ``Female'' lakes).
\item \textbf{Mount Bambouto}: a large caldera volcano west of Mbouda, built of trachytes and rhyolites over a basaltic base.
\item \textbf{Mount Oku} (about $3\,011\ \mathrm{m}$): the second highest peak of the line, in the Bamenda Highlands, surrounded by crater lakes including Lake Oku.
\item \textbf{The Bamenda and Adamaoua massifs}: broad volcanic plateaux made of stacked basalt flows capped by trachytic and rhyolitic domes and plugs. The Adamaoua Plateau forms a horst-like highland extending the line northeastward towards Lake Chad.
\end{itemize}

Scattered across this sector are numerous \cle{maars} and crater lakes -- Lakes Nyos, Monoun, Oku, and the Manengouba lakes -- formed by explosive encounters between rising magma and groundwater (\emph{phreatomagmatic} explosions).

\courssub{Petrology of the CVL Lavas and Their Mantle Origin}

CVL magmas form a classic \cle{alkaline differentiation series}. A parent magma of basanite or alkali basalt composition evolves by fractional crystallisation (crystals settle out as the magma cools, changing the composition of the remaining liquid) towards trachyte, phonolite and rhyolite.

\begin{center}
\begin{tabularx}{\linewidth}{|l|X|X|}
\hline
\rowcolor{popblueL} \textbf{Rock type} & \textbf{Approx. SiO$_2$ content} & \textbf{Character} \\
\hline
Basanite & 42--45\,\% & Strongly undersaturated, nepheline-bearing \\
\hline
Alkali basalt & 45--50\,\% & Dominant lava of Mount Cameroon \\
\hline
Trachyte & 60--65\,\% & Alkali feldspar rich, forms domes and flows \\
\hline
Phonolite & 55--60\,\% & Nepheline + alkali feldspar, forms plugs \\
\hline
Rhyolite & $>$ 70\,\% & Silica-rich, viscous, forms domes and ignimbrites \\
\hline
\end{tabularx}
\end{center}

\begin{popfigure}
\begin{tikzpicture}
\begin{axis}[
ybar, bar width=0.55cm,
width=11cm, height=6.2cm,
ylabel={SiO$_2$ content (wt \%)},
symbolic x coords={Basanite, Alkali basalt, Trachyte, Phonolite, Rhyolite},
xtick=data, x tick label style={rotate=20, anchor=east, font=\scriptsize},
ymin=0, ymax=80,
nodes near coords, every node near coord/.append style={font=\scriptsize},
axis lines=left,
]
\addplot[fill=poporange, draw=popdark] coordinates {(Basanite,44) (Alkali basalt,48) (Trachyte,62) (Phonolite,57) (Rhyolite,72)};
\end{axis}
\end{tikzpicture}
\legende{Typical silica contents of CVL lavas, showing the alkaline differentiation trend from basanite (about 44\,\% SiO$_2$) to rhyolite (over 70\,\% SiO$_2$).}
\end{popfigure}

A decisive clue to the origin of these magmas is the presence of \cle{xenoliths}: fragments of \emph{peridotite} (olivine- and pyroxene-rich rock) torn from the \cle{mantle} and carried up by the rising basalt. Finding pieces of mantle inside a lava proves that the magma originated in the mantle itself, at depths of several tens of kilometres, and rose quickly enough not to be completely transformed on the way.

\begin{popfigure}
\begin{tikzpicture}[scale=0.95]
% Volcano edifice
\fill[popgreenL] (0,0) -- (2.2,4.2) -- (3.0,4.6) -- (3.8,4.2) -- (6,0) -- cycle;
\draw[popdark, thick] (0,0) -- (2.2,4.2) -- (3.0,4.6) -- (3.8,4.2) -- (6,0);
% Crater
\fill[white] (2.55,4.2) -- (3.0,4.75) -- (3.45,4.2) -- cycle;
% Lava flows
\draw[poporange, very thick] (2.6,4.25) .. controls (1.8,3.4) and (1.2,2.2) .. (0.4,0.6);
\draw[poporange, very thick] (3.4,4.25) .. controls (4.3,3.3) and (5.0,1.8) .. (5.7,0.5);
\node[font=\scriptsize, poporange, anchor=west] at (4.9,2.6) {lava flows};
% Pyroclastic layers
\draw[poppurple, dashed, thick] (0.5,0.9) -- (2.6,3.9);
\draw[poppurple, dashed, thick] (3.4,3.9) -- (5.5,0.9);
\node[font=\scriptsize, poppurple, anchor=east, align=center, text width=2.2cm] at (1.1,1.7) {pyroclastic layers};
% Conduit
\draw[popink, very thick] (3.0,4.3) -- (3.0,1.1);
\node[font=\scriptsize, popink, anchor=west] at (3.1,2.6) {conduit (dyke)};
% Magma chamber
\fill[poppink] (3.0,0.55) ellipse (1.5 and 0.55);
\draw[popdark, thick] (3.0,0.55) ellipse (1.5 and 0.55);
\node[font=\scriptsize, popdark] at (3.0,0.55) {magma chamber};
% Basement
\fill[gray!20] (-0.5,-0.9) rectangle (6.5,0);
\draw[popdark] (-0.5,0) -- (6.5,0);
\node[font=\scriptsize, popmuted] at (1.0,-0.45) {crystalline basement};
% Xenoliths
\fill[popgold] (2.6,0.45) circle (0.07);
\fill[popgold] (3.4,0.7) circle (0.07);
\node[font=\scriptsize, popgold, anchor=west] at (4.6,0.55) {mantle xenoliths};
% Ash cloud
\node[font=\scriptsize, popmuted, align=center] at (3.0,5.3) {ash and gas plume};
\draw[popmuted, ->] (3.0,4.8) -- (3.0,5.1);
\end{tikzpicture}
\legende{Cross-section of a CVL volcano: mantle-derived magma ponds in a crustal chamber, rises through a conduit, and builds an edifice of alternating lava flows and pyroclastic layers.}
\end{popfigure}

\courssub{Theories of the Origin of the CVL}

Why does this line of volcanoes exist at all, in the middle of the African plate? Two competing models dominate the debate:

\begin{enumerate}
\item \textbf{The hot line / mantle plume model}: a narrow, elongated upwelling of hot mantle (a ``hot line'') beneath the lithosphere supplies the magma. This explains the volume and the alkaline chemistry of the lavas.
\item \textbf{The lithospheric fracture model}: the CVL follows ancient zones of weakness -- the \cle{Pan-African shear zones}, such as the Central African Shear Zone (CASZ) and the Foumban--Mbere shear system -- which were reactivated in the Cenozoic. Magma from modest mantle melting simply used these fractures as pathways. This explains the perfect straightness of the line and its alignment with mapped faults.
\end{enumerate}

Most geologists today accept a \textbf{combined model}: reactivated fractures localise the volcanism, while anomalously hot mantle beneath provides the melt.

\begin{propriete}[No age progression along the CVL]
Unlike a classic hotspot track such as Hawaii, where volcano ages increase steadily with distance from the active centre (because the plate moves over a fixed plume), the CVL shows \textbf{no systematic age migration}: old and young volcanoes occur side by side along the line, and several centres have been active simultaneously. This is a key argument against a simple fixed-plume origin and in favour of structural (fracture) control. (The comparison itself is examinable; no formal proof is required.)
\end{propriete}

\begin{attention}[The CVL is NOT a subduction-zone arc]
A frequent exam error is to compare the CVL with the Andes or Japan. Cameroon lies \textbf{far from any active plate boundary}; there is no subducting slab beneath it. CVL volcanism is \emph{intraplate} and alkaline, whereas arc volcanism is typically calc-alkaline (andesites) and explosive because of water released from the slab. Confusing the two tectonic settings costs marks.
\end{attention}

\begin{savaistu}[Exam tip: compare the CVL with the East African Rift]
A favourite GCE question asks you to compare the CVL with East African Rift volcanism. Similarities: both are Cenozoic, intraplate, alkaline, mantle-derived. Differences: East African volcanism is tied to active continental \emph{rifting} (the plate is splitting apart, with large grabens), whereas the CVL is a linear alignment on a \emph{stable} plate, controlled by reactivated shear zones, and it uniquely crosses the ocean--continent boundary without changing character.
\end{savaistu}

\courssub{Mount Cameroon: an Active Volcano}

Mount Cameroon is one of Africa's most active volcanoes. Recorded historical eruptions include those of \textbf{1922, 1959, 1982, 1999 and 2000}, with earlier events reported in the nineteenth century (e.g. 1800, 1865). The typical eruption style is \emph{effusive to mildly explosive}: fissures open on the flanks, small \cle{spatter cones} and cinder cones form, and fluid basaltic lava flows travel many kilometres -- in 1999 and 2000, flows descended towards the coast and threatened the road and plantations near Limbe and Idenau. This Hawaiian-style behaviour reflects the low viscosity of the hot, silica-poor alkali basalt.

Monitoring is carried out by the \cle{Ekona Geophysical Observatory} (near Buea), which operates seismometers to detect the earthquakes that announce magma movement, together with ground-deformation (GPS, tiltmeters) and gas measurements. Because the volcano rises directly above densely populated towns (Buea, Limbe, Mutengene), eruption early warning is a national priority.

\begin{exoresolu}[Interpreting a CVL lava suite]
\textbf{Statement.} A student analyses three lavas from the continental CVL: sample A contains 47\,\% SiO$_2$ with olivine phenocrysts and mantle peridotite xenoliths; sample B contains 62\,\% SiO$_2$ and is rich in alkali feldspar; sample C contains 73\,\% SiO$_2$. (a) Name each rock. (b) Explain the relationship between the three samples. (c) What do the xenoliths in A tell us about the magma source?
\tcblower
\textbf{Solution.} (a) Sample A is an \textbf{alkali basalt} (silica-poor, olivine-bearing); sample B is a \textbf{trachyte}; sample C is a \textbf{rhyolite}. (b) The three rocks belong to the same \textbf{alkaline differentiation series}: the basaltic parent magma, rising from the mantle, evolved by \emph{fractional crystallisation} -- early-formed crystals (olivine, pyroxene, plagioclase) were removed, progressively enriching the residual liquid in silica and alkalis to produce trachyte and finally rhyolite. (c) The peridotite xenoliths are fragments of the \textbf{upper mantle} ripped up during ascent; they prove that the parent magma originated by \emph{partial melting of the mantle} and rose rapidly enough to carry these dense fragments to the surface.
\end{exoresolu}

\begin{exercice}[Check your understanding]
\begin{enumerate}
\item Define the Cameroon Volcanic Line and state two features that make it unique among the world's volcanic chains.
\item Using texture only, explain how you would distinguish a granite from a rhyolite in the field.
\item Give one argument in favour of the mantle plume model and one in favour of the fracture-control model for the CVL.
\item Why are Mount Cameroon eruptions mostly effusive rather than violently explosive?
\end{enumerate}
\end{exercice}

\begin{corrige}[Exercise 1]
\begin{enumerate}
\item The CVL is a $\sim 1\,600\ \mathrm{km}$ SW--NE chain of Cenozoic alkaline volcanoes from Annobón to Lake Chad. Unique features: it crosses both oceanic and continental crust with the same magma type, and it shows no age progression along its length.
\item Granite is phaneritic: large interlocking crystals visible to the naked eye, indicating slow cooling at depth (plutonic). Rhyolite is aphanitic: fine-grained or glassy, indicating rapid cooling at the surface (volcanic). Both have the same chemistry; only the cooling history differs.
\item Plume model: the large magma volume and alkaline chemistry suggest anomalously hot upwelling mantle. Fracture model: the line is perfectly straight and coincides with reactivated Pan-African shear zones, and the absence of age migration fits structural rather than plume control.
\item Its lavas are hot, silica-poor alkali basalts of low viscosity, which allow gases to escape gently; eruptions therefore produce fissure-fed lava flows and small spatter cones rather than explosive Plinian columns.
\end{enumerate}
\end{corrige}

\begin{aretenir}[Key points of Section 1]
\begin{itemize}
\item Cameroon has two igneous families: ancient Precambrian--Palaeozoic intrusions (granites, syenites, gabbros of the basement) and Cenozoic volcanism of the CVL.
\item The CVL is a $\sim 1\,600\ \mathrm{km}$ intraplate alkaline volcanic line crossing oceanic and continental crust with no age progression.
\item CVL lavas form an alkaline series: basanite $\rightarrow$ alkali basalt $\rightarrow$ trachyte/phonolite $\rightarrow$ rhyolite; mantle xenoliths prove a mantle source.
\item Its origin combines hot mantle upwelling with reactivated Pan-African shear zones; it is \textbf{not} a subduction arc.
\item Mount Cameroon ($4\,095\ \mathrm{m}$) is an active basaltic volcano (1922, 1959, 1982, 1999, 2000) monitored from the Ekona observatory.
\end{itemize}
\end{aretenir}

\coursec{Section 2: Metamorphic Geology of Cameroon (Lessons 132-133)}

In Section 1 we studied the igneous geology of Cameroon, from the ancient granites of the basement to the young volcanoes of the Cameroon Volcanic Line. We saw that beneath the volcanic cover and the sedimentary basins lies a much older foundation of crystalline rocks. It is time to look closely at that foundation. Much of it is not igneous at all: it is made of rocks that have been \emph{transformed} deep inside the Earth by heat, pressure and chemically active fluids. These are the \cle{metamorphic rocks}, and in Cameroon they occupy more than half of the national territory. Understanding them is essential, because they record more than 3 billion years of Earth history and they host some of the country's most important mineral resources, such as the iron of Mbalam.

\courssub{The Precambrian Basement of Cameroon: Extent and Subdivision}

\begin{definition}[The basement]
The \cle{basement} is the ancient crystalline foundation of a continent, made of igneous and metamorphic rocks, which underlies the younger sedimentary cover. In Cameroon the basement is of \cle{Precambrian} age (older than 541 Ma) and crops out over most of the country, from the southern border with Gabon and Congo to the Adamawa Plateau.
\end{definition}

Imagine a house: the sedimentary basins we shall study in Section 3 are like the furniture and the floor coverings, while the basement is the concrete slab and the foundations. In Cameroon, the ``slab'' is exposed almost everywhere except in the coastal basins (Douala, Kribi--Campo, Rio del Rey) and in the northern basins (Garoua, Logone--Birni).

The Cameroonian basement is classically subdivided into two great domains:

\begin{itemize}
\item An \cle{Archaean} domain (older than 2.5 Ga) in the south: the \textbf{Ntem Complex}, with the \textbf{Nyong Group} at its north-western margin. This is the northern edge of the great \cle{Congo Craton}.
\item A \cle{Proterozoic} domain, essentially \textbf{Pan-African} (Neoproterozoic, about 750--550 Ma), covering the centre and north of the country: the \textbf{Adamawa--Yade domain} and the \textbf{Lom domain}, which belong to the \cle{Central African Fold Belt}.
\end{itemize}

\begin{popfigure}
\begin{tikzpicture}[scale=1.0]
% Simplified outline of Cameroon
\draw[thick, popink] (0,0) -- (0.3,2.2) -- (0.8,4.2) -- (1.2,5.6) -- (2.2,6.6) -- (3.4,6.2) -- (4.2,5.0) -- (4.6,3.6) -- (4.4,2.0) -- (3.6,0.6) -- (2.2,-0.4) -- (1.0,-0.2) -- cycle;
% Ntem Complex (south)
\fill[popblueL, opacity=0.85] (0,0) -- (0.3,2.2) -- (4.4,2.0) -- (3.6,0.6) -- (2.2,-0.4) -- (1.0,-0.2) -- cycle;
% Nyong Group (NW margin band)
\fill[popgreenL] (0.3,2.2) -- (0.8,4.2) -- (1.6,4.0) -- (1.4,2.3) -- cycle;
% Pan-African domains (centre-north)
\fill[popgold, opacity=0.5] (0.8,4.2) -- (1.2,5.6) -- (2.2,6.6) -- (3.4,6.2) -- (4.2,5.0) -- (4.6,3.6) -- (4.4,2.0) -- (1.4,2.3) -- (1.6,4.0) -- cycle;
% Labels
\node[align=center, font=\small\bfseries, popdark] at (2.3,0.9) {NTEM COMPLEX\\ \normalfont\scriptsize Archaean ($>$2.5 Ga)};
\node[align=center, font=\scriptsize\bfseries, popdark, rotate=62] at (0.95,3.2) {NYONG GROUP};
\node[align=center, font=\small\bfseries, popdark] at (2.9,4.3) {PAN-AFRICAN DOMAINS\\ \normalfont\scriptsize Adamawa--Yade \& Lom ($\sim$600 Ma)};
% Cities
\fill[popink] (1.1,1.6) circle (0.06); \node[right, font=\scriptsize] at (1.15,1.6) {Yaound\'e};
\fill[popink] (0.5,1.3) circle (0.06); \node[left, font=\scriptsize] at (0.45,1.3) {Douala};
\fill[popink] (2.0,0.4) circle (0.06); \node[right, font=\scriptsize] at (2.05,0.4) {Ebolowa};
\fill[popink] (2.6,4.9) circle (0.06); \node[right, font=\scriptsize] at (2.65,4.9) {Ngaound\'er\'e};
\fill[popink] (3.3,1.5) circle (0.06); \node[right, font=\scriptsize] at (3.35,1.5) {Bertoua};
\fill[popink] (3.0,0.3) circle (0.06); \node[right, font=\scriptsize] at (3.05,0.3) {Mbalam};
% North arrow
\draw[->, thick, popink] (4.9,5.4) -- (4.9,6.2); \node[above, font=\scriptsize\bfseries] at (4.9,6.2) {N};
% Legend
\fill[popblueL] (5.1,1.9) rectangle (5.5,2.2); \node[right, font=\scriptsize] at (5.55,2.05) {Archaean craton};
\fill[popgreenL] (5.1,1.4) rectangle (5.5,1.7); \node[right, font=\scriptsize] at (5.55,1.55) {Nyong Group};
\fill[popgold, opacity=0.5] (5.1,0.9) rectangle (5.5,1.2); \node[right, font=\scriptsize] at (5.55,1.05) {Pan-African belt};
\end{tikzpicture}
\legende{Sketch map of Cameroon showing the three great metamorphic subdivisions of the Precambrian basement: the Archaean Ntem Complex, the Nyong Group at the Congo Craton margin, and the Pan-African (Neoproterozoic) domains of the centre and north.}
\end{popfigure}

\begin{aretenir}[Key point]
Southern Cameroon belongs to the \cle{Congo Craton} (stable Archaean nucleus), while central and northern Cameroon belong to the \cle{Pan-African mobile belt}, a zone of Neoproterozoic mountain building that reactivated and surrounded the old craton.
\end{aretenir}

\courssub{The Archaean Ntem Complex (South Cameroon)}

The \cle{Ntem Complex} occupies the whole of South Cameroon (South, East and parts of the Centre and Littoral regions) and continues into Gabon, Equatorial Guinea and Congo. It is the Cameroonian portion of the Congo Craton, one of the oldest pieces of continental crust on Earth.

Its main rock types are:

\begin{itemize}
\item \textbf{Charnockites}: dark, greenish, hypersthene-bearing granitic gneisses formed under very high-grade (granulite facies), dry conditions. They give the Ntem Complex its characteristic massive, grey-green outcrops, well exposed along the Yaound\'e--Ebolowa and Kribi roads.
\item \textbf{TTG gneisses}: banded grey gneisses of tonalite--trondhjemite--granodiorite composition. TTG suites are the typical ``grey gneisses'' of Archaean cratons worldwide and represent some of the earliest continental crust.
\item \textbf{Greenstone belts}: elongated belts of metamorphosed volcanic and sedimentary rocks (metabasalts, amphibolites, banded iron formations) squeezed between the granitic gneisses, for example around Ebolowa and Sangmelima.
\end{itemize}

Radiometric dating (U--Pb on zircons) gives ages \textbf{greater than 2.5 Ga}, with many ages around 2.8--2.9 Ga. These rocks therefore formed in the \cle{Archaean Eon}, when the Earth was less than half its present age.

\begin{savaistu}[Did you know?]
The name ``Ntem'' comes from the Ntem river of South Cameroon, along whose valley many of the characteristic charnockitic outcrops were first described by geologists of the colonial geological surveys. The complex is sometimes called the ``Ntem--Chaillu block'' when followed into Gabon and Congo.
\end{savaistu}

\courssub{The Congo Craton Margin and the Nyong Group}

At the north-western edge of the Ntem Complex, west and south-west of Yaound\'e, lies the \cle{Nyong Group} (also called Nyong Series). It represents the reworked margin of the Congo Craton and is of Palaeoproterozoic age (about 2.1--2.4 Ga), affected by the Eburnean orogeny and later by the Pan-African event.

The Nyong Group consists of:

\begin{itemize}
\item \textbf{High-grade gneisses} and migmatitic gneisses;
\item \textbf{Amphibolites} (metamorphosed basalts and dolerites);
\item \textbf{Itabirites}, that is \cle{banded iron formations} (BIF): rocks made of alternating millimetric bands of silica (quartz) and iron oxides (magnetite, haematite);
\item Quartzites, mica schists and meta-conglomerates.
\end{itemize}

The BIF of the Nyong Group is of great economic importance: it is the protore of the iron deposits of the South (Mbalam region and its extensions towards Nkout and Bikoula), which we shall meet again below.

\courssub{The Pan-African Mobile Belt in Cameroon}

\begin{definition}[The Pan-African orogeny]
The \cle{Pan-African orogeny} (also called the Pan-African thermotectonic event) is the great Neoproterozoic mountain-building event (roughly 750--550 Ma, culminating around 600 Ma) during which several ancient cratons collided and were welded together to form the supercontinent \cle{Gondwana}. In Cameroon it produced the \textbf{Central African Fold Belt}, which wraps around the northern margin of the Congo Craton.
\end{definition}

During this orogeny, the rocks of central and northern Cameroon were buried, folded, sheared and heated, then intruded by granites. Two main domains are distinguished:

\begin{itemize}
\item The \textbf{Adamawa--Yade domain} (Centre and Adamawa regions): a vast assemblage of \textbf{schists, gneisses and migmatites}, intruded by abundant \textbf{syn- to post-tectonic granites} (e.g. the granites of the Yaound\'e, Bafia and Ngaound\'er\'e regions). Syn-tectonic granites were emplaced \emph{during} deformation (they are foliated); post-tectonic granites came after (they are massive).
\item The \textbf{Lom domain} (East region, Lom river basin): a lower-grade domain of \textbf{mica schists}, meta-sandstones and meta-volcanics, interpreted as a Neoproterozoic sedimentary basin that was folded and metamorphosed in the greenschist facies.
\end{itemize}

\courssub{Metamorphic Facies and Grades in Cameroon}

\begin{propriete}[Metamorphic grade and index minerals]
Metamorphic grade increases with temperature and pressure. As grade rises, \cle{index minerals} appear in a fixed, predictable sequence in pelitic (clay-rich) rocks: \textbf{chlorite} $\rightarrow$ \textbf{biotite} $\rightarrow$ \textbf{garnet} $\rightarrow$ \textbf{staurolite} $\rightarrow$ \textbf{kyanite} $\rightarrow$ \textbf{sillimanite}. Each index mineral marks the crossing of a metamorphic reaction (an \cle{isograd}). The order of appearance is examinable knowledge; the thermodynamic derivation is not required.
\end{propriete}

Geologists group P--T conditions into \cle{metamorphic facies}. In Cameroon the following facies are represented:

\begin{center}
\begin{tabularx}{\linewidth}{|l|X|X|}
\hline
\rowcolor{popblueL} \textbf{Facies} & \textbf{Typical conditions / minerals} & \textbf{Cameroonian example} \\
\hline
Greenschist & Low grade; chlorite, muscovite, albite, epidote & Lom domain mica schists \\
\hline
Amphibolite & Medium grade; garnet, hornblende, staurolite, kyanite & Adamawa--Yade schists and gneisses; Nyong Group \\
\hline
Granulite & High grade, dry; hypersthene, sillimanite, orthopyroxene & Ntem Complex charnockites \\
\hline
\end{tabularx}
\end{center}

\begin{popfigure}
\begin{tikzpicture}[scale=1.05]
\draw[->, thick, popink] (0,0) -- (8.6,0) node[below left, font=\small] {Temperature (\textdegree C)};
\draw[->, thick, popink] (0,0) -- (0,5.6) node[rotate=90, above left=0.55cm and -1.3cm, font=\small] {Pressure (GPa)};
% Facies fields (schematic)
\fill[popgreenL] (0.2,0.2) -- (3.2,0.2) -- (4.0,2.2) -- (0.2,2.2) -- cycle;
\fill[popgold, opacity=0.45] (3.2,0.2) -- (6.4,0.2) -- (7.0,3.4) -- (4.0,2.2) -- cycle;
\fill[popblueL] (6.4,0.2) -- (8.3,0.2) -- (8.3,4.6) -- (7.0,3.4) -- cycle;
\fill[popgreenL, opacity=0.6] (0.2,2.2) -- (4.0,2.2) -- (4.6,4.8) -- (0.2,4.8) -- cycle;
% Labels
\node[align=center, font=\scriptsize\bfseries, popdark] at (1.9,1.1) {GREENSCHIST\\ \normalfont\scriptsize chlorite, biotite};
\node[align=center, font=\scriptsize\bfseries, popdark] at (5.3,1.5) {AMPHIBOLITE\\ \normalfont\scriptsize garnet, kyanite};
\node[align=center, font=\scriptsize\bfseries, popdark] at (7.5,2.2) {GRANULITE\\ \normalfont\scriptsize hypersthene};
\node[align=center, font=\scriptsize\bfseries, popdark] at (2.2,3.5) {BLUESCHIST /\\ HP fields\\ \normalfont\scriptsize (rare in Cameroon)};
% Cameroonian examples
\fill[popink] (2.0,0.8) circle (0.09); \node[below, font=\scriptsize, align=center] at (2.0,0.62) {Lom schists};
\fill[popink] (5.4,1.1) circle (0.09); \node[below, font=\scriptsize, align=center] at (5.4,0.85) {Adamawa--Yade\\ gneisses};
\fill[popink] (7.6,1.2) circle (0.09); \node[below, font=\scriptsize, align=center] at (7.6,0.95) {Ntem\\ charnockites};
\fill[popink] (4.8,2.6) circle (0.09); \node[right, font=\scriptsize, align=left] at (4.9,2.6) {Nyong Group\\ amphibolites};
\end{tikzpicture}
\legende{Schematic pressure--temperature (P--T) diagram of metamorphic facies, with typical Cameroonian terrains placed in their approximate fields.}
\end{popfigure}

\courssub{Progressive Metamorphism: From Schist to Gneiss to Migmatite}

One of the most instructive field observations in Cameroon is made along the roads of the Centre region, where one can follow, over a few kilometres, a progressive transformation of the same original sedimentary pile:

\begin{itemize}
\item \textbf{Schist}: fine-grained, strongly foliated (\cle{schistosity}), splits into sheets; chlorite and biotite visible.
\item \textbf{Gneiss}: coarser, with alternating light (quartz--feldspar) and dark (biotite--hornblende) bands: the \cle{gneissic banding}; garnet porphyroblasts common.
\item \textbf{Migmatite}: the rock begins to melt partially; light granitic veins (\cle{leucosome}) segregate from a darker refractory residue (\cle{melanosome}), with an unmelted relic (\cle{mesosome}).
\end{itemize}

\begin{popfigure}
\begin{tikzpicture}[scale=1.0]
% Schist zone
\fill[popgreenL] (0,0) -- (3.4,0) -- (3.4,3.4) -- (0,3.4) -- cycle;
\foreach \y in {0.3,0.7,1.1,1.5,1.9,2.3,2.7,3.1} \draw[popgreen, thin] (0.1,\y) .. controls (1.2,\y+0.15) and (2.2,\y-0.15) .. (3.3,\y);
\node[align=center, font=\small\bfseries, popdark] at (1.7,-0.4) {SCHIST};
\node[align=center, font=\scriptsize] at (1.7,3.7) {chlorite, biotite};
% Gneiss zone
\fill[popgold, opacity=0.4] (3.4,0) -- (6.8,0) -- (6.8,3.4) -- (3.4,3.4) -- cycle;
\foreach \y in {0.35,1.05,1.75,2.45,3.15} {\draw[popdark, thin] (3.5,\y) -- (6.7,\y); \draw[popmuted, thick] (3.5,\y+0.28) -- (6.7,\y+0.28);}
\node[align=center, font=\small\bfseries, popdark] at (5.1,-0.4) {GNEISS};
\node[align=center, font=\scriptsize] at (5.1,3.7) {garnet, banding};
% Migmatite zone
\fill[popblueL] (6.8,0) -- (10.2,0) -- (10.2,3.4) -- (6.8,3.4) -- cycle;
\foreach \y in {0.5,1.4,2.3,3.0} \draw[white, very thick] (6.9,\y) .. controls (7.8,\y+0.4) and (9.0,\y-0.4) .. (10.1,\y);
\foreach \y in {0.5,1.4,2.3,3.0} \draw[popblue, thin] (6.9,\y) .. controls (7.8,\y+0.4) and (9.0,\y-0.4) .. (10.1,\y);
\node[align=center, font=\small\bfseries, popdark] at (8.5,-0.4) {MIGMATITE};
\node[align=center, font=\scriptsize] at (8.5,3.7) {partial melting};
% Arrow of increasing grade
\draw[->, ultra thick, poporange] (0.3,-1.0) -- (9.9,-1.0);
\node[below, font=\small\bfseries, poporange] at (5.1,-1.05) {Increasing temperature and pressure (increasing metamorphic grade)};
\end{tikzpicture}
\legende{Progressive metamorphism of a clay-rich sediment: schist (foliated) $\rightarrow$ gneiss (banded) $\rightarrow$ migmatite (partially molten, with light leucosome veins).}
\end{popfigure}

\courssub{Polymetamorphism and Migmatisation}

Many Cameroonian basement rocks show evidence of \cle{polymetamorphism}: they have suffered \emph{several} tectono-thermal events. For example, the Nyong Group rocks carry an Eburnean ($\sim$2.1 Ga) metamorphic imprint overprinted by a Pan-African ($\sim$600 Ma) one; and the northern edge of the Ntem Complex was locally re-heated and re-foliated during the Pan-African orogeny. Evidence includes folded older foliations cut by younger ones, several generations of garnet, and zircons with old cores and young rims.

\cle{Migmatisation} --- the partial melting of high-grade gneisses --- is widespread in the Adamawa--Yade domain and marks the transition between metamorphism and magmatism: the leucosomes of migmatites are literally the first granitic melts, which, if they gather and migrate, form the Pan-African granitic plutons.

\begin{attention}[Common mistake]
Do not confuse \textbf{migmatites} with purely \textbf{igneous granites}. A migmatite is a \emph{mixed} rock: it still contains its unmelted metamorphic residue (mesosome/melanosome) with a preserved foliation, and the granitic part (leucosome) forms veins and patches \emph{within} it. A granite is a wholly magmatic rock that crystallised from a melt, is generally massive and homogeneous, and cuts across the surrounding structures. In the GCE practical, look for the inherited banding: that is the signature of the migmatite.
\end{attention}

\begin{methode}[Identifying a metamorphic rock in hand specimen]
\begin{enumerate}
\item \textbf{Look for foliation}: does the rock split along parallel planes (schistosity $\rightarrow$ schist) or show alternating light and dark bands (gneissic banding $\rightarrow$ gneiss)?
\item \textbf{Look for lineation}: stretched minerals, aligned needles or rods on the foliation planes (indicates the direction of tectonic transport).
\item \textbf{Estimate grain size}: fine (slate, phyllite), medium (schist), coarse (gneiss, migmatite).
\item \textbf{Identify the mineral assemblage}: chlorite/biotite (low grade), garnet/staurolite/kyanite (medium to high grade), sillimanite/hypersthene (very high grade).
\item \textbf{Check for melting textures}: light granitic veins in a darker foliated host $=$ migmatite.
\item \textbf{Conclude} with rock name, facies and probable protolith (e.g. garnet mica schist, amphibolite facies, from a shale).
\end{enumerate}
\end{methode}

\courssub{Economic Significance of Metamorphic Terrains}

The metamorphic basement of Cameroon is far from being of purely academic interest:

\begin{itemize}
\item \textbf{Iron}: the itabirites (BIF) of the Nyong Group and its southern extensions host the large iron deposits of \textbf{Mbalam} and neighbouring prospects (Nkout), among the most important mining projects in the country.
\item \textbf{Kyanite}: an aluminium silicate used in refractory ceramics, known in Pan-African schists of the Centre region.
\item \textbf{Marble}: metamorphosed limestones occur in the northern domains and are used for cement and ornamental stone.
\item \textbf{Construction materials}: gneisses, charnockites and granites are quarried everywhere for aggregate, blocks and lateritic gravels; the garnet-bearing gneisses of the Yaound\'e area supply much of the capital's crushed stone.
\item \textbf{Gold}: artisanal gold washing in the East region (Batouri, B\'etar\'e Oya) exploits weathered Pan-African schists and associated quartz veins.
\end{itemize}

\begin{exoresolu}[Worked exercise: recognising metamorphic terrains]
\textbf{Statement.} A student on a field trip between Yaound\'e and Ebolowa collects two samples. Sample A is a dark, coarse, grey-green rock with visible hypersthene and no foliation; radiometric data give an age of 2.8 Ga. Sample B is a banded rock with alternating red-brown (haematite) and white (quartz) millimetric bands. (1) Identify each rock and its metamorphic facies or type. (2) Assign each sample to its geological domain in Cameroon. (3) State the economic importance of sample B.
\tcblower
\textbf{Solution.}
\begin{enumerate}
\item Sample A is a \textbf{charnockite}: the presence of hypersthene (an orthopyroxene) and the absence of hydrous minerals indicate \textbf{granulite facies} (high temperature, dry conditions). Sample B is an \textbf{itabirite} (banded iron formation, BIF), a chemically banded meta-sedimentary rock.
\item Sample A, with its Archaean age of 2.8 Ga, belongs to the \textbf{Ntem Complex} (Congo Craton, South Cameroon). Sample B is characteristic of the \textbf{Nyong Group} and its extensions at the craton margin.
\item Sample B (BIF) is the protore of the \textbf{iron deposits of the Mbalam region}: enrichment of the itabirite (removal of silica, concentration of iron oxides) produces high-grade iron ore of major economic value.
\end{enumerate}
\end{exoresolu}

\begin{savaistu}[Exam tip: link to Gondwana]
A frequent GCE question asks you to place the Pan-African belt of Cameroon in its global context. Remember: the Pan-African orogeny corresponds to the \textbf{collision and assembly of the cratons that formed the supercontinent Gondwana} around 600--550 Ma. In Cameroon, the Central African Fold Belt is the suture zone between the Congo Craton (to the south) and the Saharan/Nigerian shields (to the north). Drawing a simple sketch with the Congo Craton, the collision zone and the Pan-African belt earns full marks.
\end{savaistu}

\begin{exercice}[Consolidation exercise]
\begin{enumerate}
\item Define the terms \emph{basement}, \emph{craton} and \emph{mobile belt}, and give one Cameroonian example of each.
\item Draw a labelled table comparing the Ntem Complex and the Adamawa--Yade domain (age, main rock types, metamorphic facies, tectonic event).
\item Explain, with the index-mineral sequence, how a geologist recognises increasing metamorphic grade when travelling from the Lom basin towards the Adamawa Plateau.
\item Why is a migmatite considered as evidence of the transition between metamorphism and magmatism?
\item Sketch a cross-section showing the relationship between the Congo Craton, the Nyong Group and the Pan-African belt.
\end{enumerate}
\end{exercice}

\begin{corrige}[Exercise n\textdegree{}1]
\begin{enumerate}
\item \textbf{Basement}: ancient crystalline (igneous + metamorphic) foundation beneath the sedimentary cover --- e.g. the Precambrian basement of South Cameroon. \textbf{Craton}: an old, stable, undeformed part of continental crust --- e.g. the Congo Craton (Ntem Complex). \textbf{Mobile belt}: a long zone of deformation and metamorphism between cratons --- e.g. the Pan-African Central African Fold Belt.
\item Comparison table:
\begin{center}
\begin{tabularx}{\linewidth}{|l|X|X|}
\hline
\rowcolor{popblueL} \textbf{Criterion} & \textbf{Ntem Complex} & \textbf{Adamawa--Yade domain} \\
\hline
Age & Archaean, $>$2.5 Ga & Neoproterozoic, $\sim$600 Ma \\
\hline
Main rocks & Charnockites, TTG gneisses, greenstones & Schists, gneisses, migmatites, granites \\
\hline
Facies & Granulite & Greenschist to amphibolite (migmatisation) \\
\hline
Event & Archaean crust formation & Pan-African orogeny \\
\hline
\end{tabularx}
\end{center}
\item From Lom to Adamawa the grade rises: chlorite-zone schists (greenschist facies) pass to biotite then garnet schists, then to garnet--sillimanite gneisses and migmatites (amphibolite facies). Each new index mineral (chlorite $\rightarrow$ biotite $\rightarrow$ garnet $\rightarrow$ sillimanite) marks the crossing of an isograd, i.e. higher temperature and pressure.
\item A migmatite contains both an unmelted metamorphic residue (foliated mesosome/melanosome) and a newly formed granitic melt (leucosome). It therefore records the very moment when metamorphic temperatures became high enough for partial melting: the leucosomes are the first magmas, which can later collect into granitic plutons. It is the natural bridge between the metamorphic and igneous realms.
\item The cross-section should show the Congo Craton (Ntem Complex) to the south, the Nyong Group as a narrow zone at its margin, and the Pan-African belt (Adamawa--Yade and Lom domains) to the north, with appropriate labels for rock types and metamorphic grade increasing northwards.
\end{enumerate}
\end{corrige}

\begin{aretenir}[To conclude this section]
The metamorphic geology of Cameroon rests on three pillars: the \cle{Archaean Ntem Complex} (charnockites, TTG, greenstones, $>$2.5 Ga) forming the Congo Craton; the \cle{Nyong Group} at its margin (gneisses, amphibolites, iron-bearing itabirites); and the \cle{Pan-African belt} (Adamawa--Yade and Lom domains, $\sim$600 Ma) with its schists, gneisses, migmatites and granites. Metamorphic grade increases with temperature and pressure, revealed by the index-mineral sequence, and many rocks are polymetamorphic. These terrains host iron, kyanite, marble, gold and construction materials. In the next section we shall see what lies \emph{on top} of this basement: the sedimentary basins of Cameroon.
\end{aretenir}

\coursec{Section 3: Sedimentary Geology and Basins of Cameroon (Lessons 134--136)}

In Section 2 we saw that the metamorphic basement of Cameroon --- the gneisses, schists and migmatites of the Congo craton and the Pan-African mobile belts --- forms the ancient, crystalline ``floor'' of the country. But this basement is not everywhere exposed. Over large areas it is buried beneath a cover of younger, unmetamorphosed \emph{sedimentary rocks}: sandstones, shales, limestones and evaporites, deposited layer upon layer in subsiding depressions. These depressions, and the rocks that fill them, are the subject of this section. They are not merely a geological curiosity: they host Cameroon's oil and gas (Rio del Rey), its salt springs (Mamfe), its artesian water (Chad Basin) and much of its agricultural land. Understanding them is understanding the economic geology of the country.

\courssub{The Concept of a Sedimentary Basin}

Imagine a region of the Earth's crust that slowly sinks over millions of years, like a broad, shallow bowl being gently pressed down. Rivers, wind and the sea continuously deliver sediment --- sand, mud, shells, salts --- into this depression. Because the floor keeps subsiding, space is always available for new sediment, and an enormous thickness of strata, sometimes several kilometres, can accumulate in one place while neighbouring regions receive almost nothing. Such a region is called a \cle{sedimentary basin}.

\begin{definition}[Sedimentary basin]
A \cle{sedimentary basin} is a region of the Earth's crust that has undergone \textbf{prolonged subsidence} relative to its surroundings, allowing the accumulation and preservation of a thick sequence of sedimentary rocks. The sediment fill records the history of the basin: its climate, its water bodies, its tectonic events and its life.
\end{definition}

Two ingredients are therefore essential: \textbf{subsidence} (to create accommodation space) and \textbf{sediment supply} (to fill it). Subsidence may be caused by crustal stretching (rifting), by cooling of the lithosphere after a thermal event, or by the weight of the sediment itself.

Cameroonian basins are conveniently classified into three families according to their tectonic setting:

\begin{center}
\begin{tabularx}{\linewidth}{|l|X|X|}
\hline
\rowcolor{popblueL} \textbf{Type of basin} & \textbf{Origin of subsidence} & \textbf{Cameroonian examples} \\
\hline
Coastal / passive margin basin & Cooling and sediment loading after Atlantic opening & Douala, Kribi--Campo, Rio del Rey \\
\hline
Intracratonic basin & Slow sagging of stable continental crust & Chad Basin (Logone--Birni), Bertoua, Batouri \\
\hline
Rift basin & Crustal stretching along faults (grabens) & Mamfe Basin, Benue Trough (Garoua) \\
\hline
\end{tabularx}
\end{center}

\begin{popfigure}
\begin{tikzpicture}[scale=0.9, font=\small]
% Outline of Cameroon (stylised)
\draw[thick, popink] (0,0) -- (0.4,1.6) -- (0.2,3.2) -- (1.2,4.6) -- (1.6,6.2)
 -- (2.4,7.4) -- (3.8,8.2) -- (5.6,8.0) -- (6.4,6.8) -- (6.2,5.2)
 -- (7.0,4.2) -- (6.6,2.6) -- (5.4,1.4) -- (4.2,0.4) -- (2.6,-0.4) -- (1.2,-0.3) -- cycle;
% Ocean
\fill[popblueL] (-1.6,-0.8) rectangle (0.1,3.4);
\node[popblue, rotate=75] at (-0.7,1.4) {Atlantic Ocean};
% Basins
\fill[popblue, opacity=0.45] (0.5,0.6) ellipse (0.9 and 0.7);
\node[align=center] at (0.5,0.6) {\textbf{Douala}\\ \textbf{Rio del Rey}};
\fill[popblue, opacity=0.35] (1.6,-0.1) ellipse (0.9 and 0.45);
\node at (1.6,-0.35) {\textbf{Kribi--Campo}};
\fill[poporange, opacity=0.5] (1.9,2.6) ellipse (0.85 and 0.5);
\node at (1.9,2.6) {\textbf{Mamfe}};
\fill[popgreen, opacity=0.45] (3.6,5.6) ellipse (1.3 and 0.6);
\node[align=center] at (3.6,5.6) {\textbf{Garoua / Benue}};
\fill[poppurple, opacity=0.35] (4.6,7.6) ellipse (1.5 and 0.7);
\node[align=center] at (4.6,7.6) {\textbf{Chad Basin}\\ \textbf{(Logone--Birni)}};
\fill[popgold, opacity=0.5] (5.6,2.2) ellipse (0.8 and 0.5);
\node[align=center] at (5.6,2.2) {\textbf{Bertoua}\\ \textbf{Batouri}};
% Basement label
\node[align=center, popmuted] at (3.4,3.6) {Precambrian\\ basement};
% North arrow
\draw[->, thick] (7.4,7.6) -- (7.4,8.4);
\node at (7.4,8.6) {N};
\end{tikzpicture}
\legende{Schematic map of Cameroon locating the principal sedimentary basins: Douala--Rio del Rey and Kribi--Campo (coastal), Mamfe (rift), Garoua--Benue (rift), Chad Basin and the small eastern basins (intracratonic).}
\end{popfigure}

\courssub{The Douala / Kribi--Campo Coastal Basin}

When the South Atlantic Ocean began to open in the Early Cretaceous, the future coast of Cameroon subsided as the stretched crust cooled. The sea invaded this margin, and from the Cretaceous to the present day a thick wedge of \textbf{marine and deltaic sediments} accumulated: this is the \cle{Douala Basin} and its southern continuation, the \cle{Kribi--Campo Basin}.

The stratigraphic succession records a general story of alternating marine transgressions and deltaic progradations:

\begin{itemize}
\item \textbf{Cretaceous}: basal sandstones and conglomerates resting on the basement, overlain by marine shales and limestones deposited as the sea deepened;
\item \textbf{Palaeogene}: shales with interbedded sandstones, deposited in shallow marine to deltaic settings;
\item \textbf{Neogene to Present}: thick deltaic sands and muds of the Niger--Rio del Rey delta system, prograding seawards.
\end{itemize}

\begin{popfigure}
\begin{tikzpicture}[scale=1.0, font=\small]
% Stratigraphic log frame
\draw[thick] (0,0) rectangle (2.2,8);
% Basement
\fill[popmuted!60] (0,0) rectangle (2.2,0.7);
\foreach \x in {0.2,0.7,1.2,1.7} {\draw (\x,0.1) -- (\x+0.4,0.6);}
\node[align=center] at (3.9,0.35) {Precambrian basement (gneiss)};
% Basal sandstone
\fill[popgold!70] (0,0.7) rectangle (2.2,1.5);
\foreach \x in {0.15,0.55,0.95,1.35,1.75} {\fill (\x,1.0) circle (0.03); \fill (\x+0.2,1.25) circle (0.03);}
\node[align=center] at (3.9,1.1) {Basal sandstone (Cretaceous)};
% Marine shale
\fill[popblueL] (0,1.5) rectangle (2.2,2.9);
\foreach \y in {1.7,2.0,2.3,2.6} {\draw[popblue] (0.1,\y) -- (2.1,\y);}
\node[align=center] at (3.9,2.2) {Marine shale (source rock)};
% Limestone
\fill[popgreenL] (0,2.9) rectangle (2.2,3.7);
\draw (0.3,3.1) rectangle (0.7,3.4); \draw (1.2,3.2) rectangle (1.6,3.5);
\node[align=center] at (3.9,3.3) {Limestone (marine transgression)};
% Shale
\fill[popblueL] (0,3.7) rectangle (2.2,4.8);
\foreach \y in {3.95,4.25,4.55} {\draw[popblue] (0.1,\y) -- (2.1,\y);}
\node[align=center] at (3.9,4.25) {Shale (seal rock)};
% Deltaic sandstone
\fill[popgold!70] (0,4.8) rectangle (2.2,6.0);
\foreach \x in {0.15,0.55,0.95,1.35,1.75} {\fill (\x,5.1) circle (0.03); \fill (\x+0.2,5.5) circle (0.03);}
\node[align=center] at (3.9,5.4) {Deltaic sandstone (reservoir)};
% Top shale/sand alternation
\fill[popblueL] (0,6.0) rectangle (2.2,6.6);
\fill[popgold!70] (0,6.6) rectangle (2.2,7.3);
\fill[popblueL] (0,7.3) rectangle (2.2,8.0);
\node[align=center] at (3.9,7.0) {Alternating deltaic\\ sands and muds (Neogene)};
% Age axis
\node[rotate=90] at (-0.5,4) {younger $\longrightarrow$};
\end{tikzpicture}
\legende{Schematic stratigraphic log of the Douala Basin, showing the sandstone--shale--limestone alternations above the Precambrian basement and the petroleum roles of each unit.}
\end{popfigure}

\courssub{The Rio del Rey and the Niger Delta Connection: Hydrocarbons}

The \cle{Rio del Rey} estuary, at the head of the Bight of Biafra near the Nigerian border, is the heart of Cameroon's petroleum industry. Geologically it is the western lobe of the great \textbf{Niger Delta} system: an enormous pile of deltaic sandstones and shales, several kilometres thick, built out into the ocean since the Palaeogene.

\begin{propriete}[Conditions for a hydrocarbon accumulation]
A commercial accumulation of oil or gas requires four elements to coexist:
\begin{enumerate}
\item a \cle{source rock} --- organic-rich shale in which kerogen is converted to oil and gas by burial and heating;
\item a \cle{reservoir rock} --- a porous and permeable rock (typically deltaic sandstone) that stores the hydrocarbons;
\item a \cle{seal} (cap rock) --- an impermeable shale that prevents upward escape;
\item a \cle{trap} --- a geometric configuration (anticline, fault closure, pinch-out) that concentrates the accumulation.
\end{enumerate}
All four are present in the Rio del Rey: mature marine shales (source), deltaic sandstones (reservoir), thick mudstones (seal), and rollover anticlines associated with growth faults (traps). (The proof of these roles is descriptive, not examinable as a formal demonstration, but you must be able to \emph{apply} them to a given section.)
\end{propriete}

A characteristic feature of deltaic margins is the \textbf{growth fault}: a listric (curved) normal fault that moves while sediment is being deposited, so that the beds thicken towards the fault. As the hanging-wall block slides down, the beds roll over into an anticline --- a classic oil trap.

\courssub{The Mamfe Basin: a Cretaceous Intracontinental Rift}

Crossing the Manyu Division of the South-West Region, the \cle{Mamfe Basin} is a small but instructive basin. It is an \textbf{intracontinental rift basin}: a fault-bounded trough (a graben or half-graben) formed when the crust stretched during the Early Cretaceous, at the time the South Atlantic was beginning to open. It is in fact a side-branch of the Benue Trough system.

Its fill is entirely \textbf{continental}: fluvial sandstones and conglomerates, floodplain shales, and --- crucially --- \textbf{evaporites}. In a closed basin under an arid climate, lake waters evaporated and precipitated salts. Today these evaporites feed the famous \textbf{salt springs} of the Mamfe area, exploited traditionally for centuries. Traces of \textbf{lead--zinc} mineralisation (galena, sphalerite) occur in the basin, related to brines circulating through the rift sediments.

\begin{popfigure}
\begin{tikzpicture}[scale=1.05, font=\small]
% Basement blocks
\fill[popmuted!50] (-0.5,-0.5) rectangle (2.0,2.2);
\fill[popmuted!50] (5.6,-0.5) rectangle (8.0,2.2);
\fill[popmuted!50] (2.0,-0.5) rectangle (5.6,0.4);
% Faults
\draw[very thick, popdark] (2.0,2.2) -- (2.6,-0.5);
\draw[very thick, popdark] (5.6,2.2) -- (5.0,-0.5);
% Basin fill layers (half-graben, thickening to the left fault)
\fill[popgold!60] (2.0,2.2) -- (2.6,-0.5) -- (5.0,-0.5) -- (5.6,2.2) -- cycle;
\fill[popblueL] (2.02,2.2) -- (2.35,0.6) -- (5.25,0.6) -- (5.58,2.2) -- cycle;
\fill[popgreenL] (2.05,2.2) -- (2.2,1.3) -- (5.4,1.3) -- (5.55,2.2) -- cycle;
% Labels
\node[align=center] at (0.7,2.6) {Basement\\ horst};
\node[align=center] at (6.9,2.6) {Basement\\ horst};
\node[align=center] at (3.8,1.85) {Fluvial sandstone};
\node[align=center] at (3.8,0.95) {Shale + evaporites (salt)};
\node[align=center] at (3.8,0.05) {Basal conglomerate};
\node[popdark] at (2.15,2.75) {border fault};
\node[popdark] at (5.75,2.75) {border fault};
\draw[->, thick, popdark] (2.3,2.35) -- (2.6,1.9);
\draw[->, thick, popdark] (5.4,2.35) -- (5.1,1.9);
% Arrows showing extension
\draw[->, very thick, poporange] (-0.3,3.1) -- (0.8,3.1);
\draw[->, very thick, poporange] (7.9,3.1) -- (6.8,3.1);
\node[poporange] at (3.8,3.2) {crustal extension};
\end{tikzpicture}
\legende{Block diagram (cross-section) of a half-graben rift basin applied to the Mamfe Basin: crustal extension drops a fault-bounded block, which fills with conglomerates, sandstones, shales and evaporites.}
\end{popfigure}

\begin{attention}[Common mistake: Mamfe is NOT marine]
Candidates frequently describe the Mamfe Basin as a marine basin ``because it contains salt''. This is wrong. The Mamfe Basin is a \textbf{continental rift basin}; its evaporites formed by evaporation of \emph{lake} water in a closed depression, not from the sea. The Douala Basin, by contrast, is genuinely marine to deltaic. Always ask: \emph{what do the fossils and structures say?} Marine fossils (ammonites, foraminifera) indicate the sea; evaporites with fluvial sandstones indicate a continental sabkha or salt lake.
\end{attention}

\courssub{The Garoua and Benue Trough Basins}

In North Cameroon, around Garoua, Cretaceous sediments outcrop along the upper Benue valley. They belong to the \cle{Benue Trough}, a major NE--SW rift system extending from the Niger Delta deep into Nigeria and Cameroon. The Garoua Basin contains fluvial and marine sandstones, shales and limestones of Cretaceous age, folded and faulted during a later compressional event (Santonian deformation).

\begin{exemplebox}[The Benue Trough and the South Atlantic triple junction]
When the South Atlantic opened, the rift system had the geometry of a \textbf{triple junction}: three rift arms meeting at a point. Two arms --- the future South Atlantic margins --- continued to open and became an ocean. The third arm, the Benue Trough, \emph{failed}: opening stopped, and the arm was left as a sediment-filled rift within the continent. Such an aborted arm is called a \cle{failed rift} or \emph{aulacogen}. This is why marine Cretaceous sediments are found today far inland at Garoua: the sea once penetrated up the trough.
\end{exemplebox}

\begin{attention}[Exam tip]
If a question asks you to ``account for the presence of Cretaceous marine sediments in the interior of West Africa'', the expected answer links the Benue Trough to the \textbf{failed arm of the South Atlantic triple junction}: marine waters invaded the rift during the Cretaceous, and the sediments were preserved when the arm stopped opening.
\end{attention}

\courssub{The Chad Basin (Logone--Birni)}

The extreme north of Cameroon, around Logone--Birni, lies on the southern edge of the vast \cle{Chad Basin}, an intracratonic sag basin. Its Cenozoic fill is \textbf{fluvio-lacustrine}: sands and clays deposited by rivers (the Logone, Chari) and in a huge lake. During humid climatic phases of the Quaternary, this lake expanded into \cle{Mega-Lake Chad}, far larger than today's shrunken remnant. The ancient shorelines, diatomites and lacustrine clays of the basin record these climatic oscillations. The basin's sands form important \textbf{aquifers}, a vital water resource for the arid Far North.

\courssub{Continental Interiors and Lateritic Covers}

In the east and centre of the country, small basins such as \textbf{Bertoua} and \textbf{Batouri} preserve continental sediments resting on the basement. More widespread, however, are the \cle{lateritic covers} and \cle{duricrusts}: under the hot, humid climate, intense chemical weathering leaches silica and concentrates iron and aluminium oxides, producing red laterite that can harden into an iron-rich crust (\emph{ferricrete}). These duricrusts protect plateaus, armour the landscape, and are locally enriched in \textbf{bauxite} (e.g.\ the Adamaoua and western plateaux) --- a resource we shall meet again in Section 6.

\courssub{Reconstructing Depositional Environments}

How do we know that the Douala Basin was marine while Mamfe was a continental rift? The evidence lies in the rocks themselves.

\begin{methode}[Reading a sedimentary rock to reconstruct its environment]
\begin{enumerate}
\item \textbf{Grain size and sorting}: conglomerates and coarse, poorly sorted sands indicate high-energy rivers or alluvial fans near the basin margin; well-sorted sands suggest beaches or delta channels; fine muds indicate quiet, deep or distal water.
\item \textbf{Sedimentary structures}: cross-bedding records migrating dunes or ripples in a current; graded bedding records turbidity currents on deep slopes; mud cracks and salt casts prove exposure and evaporation in a continental setting; ripple marks indicate shallow agitated water.
\item \textbf{Fossils}: marine fossils (foraminifera, ammonites, marine bivalves) prove a sea; plant remains, freshwater ostracods and fish indicate continental or lacustrine conditions.
\item \textbf{Synthesis}: combine the three lines of evidence to name the environment (fluvial, deltaic, shallow marine, deep marine, lacustrine, evaporitic).
\end{enumerate}
\end{methode}

\begin{exoresolu}[Identifying a depositional environment]
\textbf{Statement.} A borehole in the Douala Basin cuts the following succession, from base to top: (a) dark shale rich in planktonic foraminifera; (b) fine sandstone with graded bedding; (c) cross-bedded sandstone with fragments of driftwood; (d) mudstone with root traces. Interpret each bed and describe the evolution of the environment.
\tcblower
\textbf{Solution.}
\begin{itemize}
\item (a) Planktonic foraminifera are marine microfossils living in open water: the shale was deposited in an \textbf{open marine}, relatively deep environment.
\item (b) Graded bedding is the signature of \textbf{turbidity currents}: submarine flows depositing sand on the deeper slope or basin floor.
\item (c) Cross-bedding indicates strong currents; driftwood proves continental input: a \textbf{deltaic channel or mouth bar} environment.
\item (d) Root traces show a soil developed on emergent land: a \textbf{delta plain / coastal swamp}.
\end{itemize}
The succession passes from deep marine (a) up to emerged land (d): the sea became progressively shallower as the delta \textbf{prograded} (advanced seawards). This is a typical regressive, delta-building sequence of the Rio del Rey--Niger Delta system.
\end{exoresolu}

\begin{exercice}[Comparing two Cameroonian basins]
Copy and complete a table comparing the \textbf{Mamfe Basin} and the \textbf{Douala Basin} under the headings: tectonic type; age of main fill; dominant sediments; fossils expected; economic resources. Then justify, in five lines, why hydrocarbons are exploited in the Douala Basin but not in the Mamfe Basin.
\end{exercice}

\begin{corrige}[Exercise 1]
\begin{itemize}
\item \textbf{Mamfe Basin}: intracontinental rift (graben); Cretaceous; fluvial sandstones, conglomerates, shales, evaporites; continental fossils (plants, freshwater forms); salt springs, traces of Pb--Zn.
\item \textbf{Douala Basin}: coastal passive-margin basin; Cretaceous to Cenozoic; marine shales, limestones, deltaic sandstones; marine fossils (foraminifera, ammonites); oil and gas.
\item \textbf{Justification}: the Douala Basin possesses the complete petroleum system --- organic-rich marine source shales, porous deltaic sandstone reservoirs, impermeable shale seals and growth-fault/anticline traps, all buried deeply enough for oil generation. The Mamfe Basin is a small, shallow continental rift lacking mature marine source rocks and adequate burial, so no commercial hydrocarbons have formed.
\end{itemize}
\end{corrige}

\begin{aretenir}[Key points of Section 3]
\begin{itemize}
\item A sedimentary basin is a region of prolonged subsidence filled with thick sediments; Cameroon has coastal (Douala, Kribi--Campo), rift (Mamfe, Benue/Garoua) and intracratonic (Chad, Bertoua--Batouri) basins.
\item The Rio del Rey, linked to the Niger Delta, holds Cameroon's oil: source shale, reservoir sandstone, seal and growth-fault traps.
\item The Mamfe Basin is a \emph{continental} Cretaceous rift with evaporites (salt springs) and Pb--Zn traces --- never describe it as marine.
\item The Benue Trough is the failed arm of the South Atlantic triple junction, which explains marine Cretaceous sediments at Garoua.
\item The Chad Basin preserves fluvio-lacustrine deposits of Mega-Lake Chad and important aquifers.
\item Depositional environments are reconstructed from grain size, sedimentary structures and fossils.
\end{itemize}
\end{aretenir}

\coursec{Section 4: Limnology and Lacustrine Processes in Cameroon (Lessons 137-139)}

In Section 3 we studied the sedimentary basins of Cameroon --- the Douala and Rio del Rey basins in the south, the Benue trough and the Logone--Birni basin in the north. Many of those basins once held, or still hold, large bodies of standing water whose sediments record the history of the basin. We now turn from the \emph{sedimentary record} of lakes to the \emph{living lakes} themselves: how they formed, how their waters behave, and why some of them --- like Lake Nyos --- can kill. This is the science of \cle{limnology}, and Cameroon, with its crater lakes, its share of Lake Chad and its artificial reservoirs, is one of the finest natural laboratories in Africa for studying it.

\courssub{What is Limnology? Scope and Classification of Cameroonian Lakes}

\begin{definition}[Limnology]
\cle{Limnology} is the scientific study of inland waters --- lakes, ponds, reservoirs and wetlands --- including their physical properties (temperature, light, stratification), chemical properties (dissolved gases, nutrients, pH), biological communities, and the sediments deposited within them. \cle{Lacustrine processes} are the physical, chemical and biological processes that operate specifically within lakes.
\end{definition}

Why should a geologist care about lakes? Because a lake is a \emph{sediment trap}: everything eroded from its catchment, everything that grows and dies in its waters, and every chemical change in the water column is recorded layer by layer in the bottom mud. Reading lake sediments is one of the best ways to reconstruct past climates --- as we shall see with Mega-Chad.

Cameroon possesses a remarkable diversity of lakes, which we classify by their \textbf{mode of origin}:

\begin{center}
\begin{tabularx}{\linewidth}{|l|X|X|}
\hline
\rowcolor{popblueL} \textbf{Origin} & \textbf{Process} & \textbf{Cameroonian examples} \\
\hline
Volcanic (crater/maar lakes) & Water fills explosion craters (maars) or calderas of the Cameroon Volcanic Line & Nyos, Monoun, Barombi (Kumba), Bambili, Oku, Manengouba twin lakes \\
\hline
Tectonic / endorheic & Subsiding continental basin with internal drainage & Lake Chad \\
\hline
River-dammed (natural) & Lava flow or landslide dams a valley & Lake Oku (lava-dammed crater); small CVL valley lakes \\
\hline
Artificial (reservoirs) & Dam built across a river for hydropower/irrigation & Lagdo, Bamendjing, Mape, Mbakaou \\
\hline
\end{tabularx}
\end{center}

\begin{popfigure}
\begin{center}
\begin{tikzpicture}[scale=1.0]
% Cameroon outline (stylised)
\draw[thick, popink, fill=popblueL!40] (2.0,0) -- (3.6,0) -- (4.4,0.9) -- (4.6,2.2) -- (4.2,3.4) -- (4.6,4.6) -- (3.4,5.6) -- (3.5,6.6) -- (2.6,7.4) -- (1.7,6.8) -- (1.4,5.6) -- (0.6,4.8) -- (0.2,3.4) -- (0.9,2.2) -- (0.7,1.0) -- cycle;
% CVL axis
\draw[very thick, poporange, dashed] (0.5,0.3) -- (1.6,2.4) -- (2.4,4.4);
\node[poporange, font=\small, align=center] at (0.35,2.6) {CVL};
% Lakes
\fill[popblue] (2.5,5.6) circle (0.09); \node[right, font=\small] at (2.65,5.6) {Nyos};
\fill[popblue] (2.35,5.2) circle (0.09); \node[left, font=\small] at (2.2,5.2) {Monoun};
\fill[popblue] (1.5,2.3) circle (0.09); \node[left, font=\small] at (1.35,2.3) {Barombi};
\fill[popblue] (2.5,4.5) circle (0.09); \node[right, font=\small] at (2.65,4.5) {Oku};
\fill[popblue] (2.7,3.9) circle (0.09); \node[right, font=\small] at (2.85,3.9) {Bambili};
\fill[popblue] (3.0,6.9) circle (0.11); \node[right, font=\small] at (3.15,6.9) {Lagdo};
\fill[popblue] (2.7,7.5) circle (0.14); \node[right, font=\small] at (2.85,7.55) {L. Chad};
\fill[popblue] (2.2,4.9) circle (0.08); \node[left, font=\small] at (2.05,4.85) {Bamendjing};
% Compass
\draw[->, thick] (5.3,6.6) -- (5.3,7.3); \node[font=\small] at (5.3,7.5) {N};
\end{tikzpicture}
\end{center}
\legende{Sketch map of Cameroon locating the principal lakes. The crater lakes (Nyos, Monoun, Barombi, Oku, Bambili) lie along or near the Cameroon Volcanic Line (dashed axis); Lake Chad and the Lagdo reservoir occupy the northern lowlands.}
\end{popfigure}

\courssub{The Volcanic Crater Lakes of the Cameroon Volcanic Line}

Along the Cameroon Volcanic Line (Section 1), explosive encounters between rising magma and groundwater produced wide, low-rimmed explosion craters called \cle{maars}. Once the eruption ended, rain and groundwater filled these craters, giving deep, steep-sided, almost circular lakes. The best known are:

\begin{itemize}
\item \textbf{Lake Nyos} (North West Region): a maar lake about \SI{208}{m} deep, infamous for the 1986 disaster.
\item \textbf{Lake Monoun} (North West Region): a smaller maar lake (about \SI{96}{m} deep), site of the 1984 gas disaster.
\item \textbf{Lake Barombi} near Kumba (South West Region): a large crater lake (about \SI{110}{m} deep) renowned for its endemic fishes.
\item \textbf{Lake Bambili} and the \textbf{twin lakes of Mount Manengouba} (Muanenguba lakes), and \textbf{Lake Oku} on the Oku massif --- a beautiful crater lake surrounded by montane forest.
\end{itemize}

These lakes share a dangerous secret: the magma that dug their craters still leaks gas upwards. Carbon dioxide of \cle{magmatic origin} seeps through the crater floor, dissolves in the cold, deep, high-pressure bottom water, and accumulates there for decades --- provided the lake does not mix.

\courssub{Stratification, Meromixis and Limnic Eruptions}

\begin{definition}[Thermal stratification]
In deep tropical lakes, sunlight warms the surface layer while the depths remain cold. Because warm water is less dense than cold water, the lake divides into a warm, wind-mixed surface layer (the \cle{epilimnion}) and a cold, dense bottom layer (the \cle{hypolimnion}), separated by a zone of rapid temperature change (the \cle{thermocline}). This density layering resists vertical mixing.
\end{definition}

\begin{definition}[Meromictic lake]
A \cle{meromictic lake} is a lake whose deep waters \emph{never} mix with the surface waters, even seasonally. The permanent stratification allows dissolved gases and salts to build up in the bottom layer (the \cle{monimolimnion}) year after year. Nyos and Monoun are meromictic.
\end{definition}

In an ordinary (\emph{holomictic}) lake, seasonal cooling of the surface makes the surface water denser, so it sinks and the whole lake \emph{overturns}, carrying dissolved gases harmlessly to the surface in small amounts. In a meromictic crater lake, this safety valve is missing. Magmatic \ce{CO2} keeps dissolving at the bottom until the water approaches saturation. Then any trigger --- a landslide, an earthquake, a violent storm, or simply the water becoming saturated --- can set off disaster.

\begin{definition}[Limnic eruption (lake overturn)]
A \cle{limnic eruption} is the sudden, massive release of dissolved gas (chiefly \ce{CO2}) from the deep water of a stratified lake when the stratification breaks down. As deep water is displaced upwards, the pressure on it falls, the dissolved gas comes out of solution (like opening a bottle of fizzy drink), the bubbles drive more water upwards, and the process runs away explosively, erupting a fountain of water and a dense cloud of gas.
\end{definition}

\begin{propriete}[Behaviour of the \ce{CO2} cloud]
Carbon dioxide is about 1.5 times denser than air (molar mass \SI{44}{g.mol^{-1}} against about \SI{29}{g.mol^{-1}} for air). The gas cloud released in a limnic eruption therefore does not rise and disperse: it \emph{hugs the ground} and flows downhill along valleys like an invisible river. \ce{CO2} is colourless, odourless and displaces oxygen, so people and animals in its path are asphyxiated without warning. (The mechanism is examinable as an explanation, not a formal proof.)
\end{propriete}

\begin{experience}[Case study: the disasters of Monoun (1984) and Nyos (1986)]
On 15 August 1984, Lake Monoun released a \ce{CO2} cloud that killed 37 people on a nearby road. On 21 August 1986, Lake Nyos underwent a far larger limnic eruption: a fountain of water and gas tens of metres high burst from the lake, and a \ce{CO2} cloud flowed down the surrounding valleys, killing about 1\,700 people and thousands of livestock, with victims recorded more than \SI{20}{km} from the lake. The lake water turned reddish-brown as iron-rich deep water mixed to the surface and oxidised. These remain the only two recorded limnic eruptions in human history --- both in Cameroon.
\end{experience}

\begin{attention}[Common mistake]
Lake Nyos did \emph{not} erupt lava or ash in 1986, and no volcanic eruption occurred. The disaster was a \emph{limnic} eruption --- a burst of dissolved gas from the lake water itself. Writing that ``the volcano Nyos erupted'' in the examination is a serious error.
\end{attention}

\begin{popfigure}
\begin{center}
\begin{tikzpicture}[scale=1.05]
% Crater walls
\draw[thick, popink, fill=popgold!25] (-4.6,0) -- (-2.6,0) .. controls (-2.3,-0.6) and (-2.1,-1.2) .. (-2.0,-2.6) -- (2.0,-2.6) .. controls (2.1,-1.2) and (2.3,-0.6) .. (2.6,0) -- (4.6,0);
% Water body
\fill[popblue!30] (-2.0,-0.35) rectangle (2.0,-2.6);
\fill[popblue!55] (-2.0,-1.7) rectangle (2.0,-2.6);
\draw[popblue, thick] (-2.0,-0.35) -- (2.0,-0.35);
\draw[popblue, dashed] (-2.0,-1.7) -- (2.0,-1.7);
\node[font=\small, popdark] at (0,-0.95) {Epilimnion (warm, low \ce{CO2})};
\node[font=\small, popdark] at (-0.7,-1.45) {thermocline};
\node[font=\small, popdark, align=center] at (0,-2.2) {Hypolimnion: cold, dense,\\\ce{CO2}-rich (monimolimnion)};
% Gas seep
\foreach \x in {-0.4,-0.15,0.1,0.35}{\fill[poporange] (\x,-2.55) circle (0.045);}
\node[font=\small, poporange, align=center] at (0,-2.95) {magmatic \ce{CO2} seep};
% Degassing pipe
\draw[very thick, popgreen] (1.2,-2.55) -- (1.2,0.55);
\draw[very thick, popgreen, ->] (1.2,0.55) -- (1.2,1.05);
\node[font=\small, popgreen, align=center] at (2.9,0.8) {degassing pipe:\\fountain of\\bottom water};
% raft
\draw[thick, popink, fill=popgreenL] (0.85,-0.35) rectangle (1.55,-0.2);
% labels
\node[font=\small, popink] at (-3.4,0.25) {crater rim};
\draw[->, popink] (-3.4,0.1) -- (-2.7,-0.05);
\end{tikzpicture}
\end{center}
\legende{Cross-section of Lake Nyos: a permanently stratified (meromictic) water column with a \ce{CO2}-charged bottom layer fed by magmatic gas, and the degassing pipe that lifts bottom water so the gas escapes safely at the surface.}
\end{popfigure}

\begin{popfigure}
\begin{center}
\begin{tikzpicture}
\begin{axis}[
  width=0.72\linewidth, height=7cm,
  xlabel={Dissolved \ce{CO2} concentration (arbitrary units)},
  ylabel={Depth (m)},
  y dir=reverse,
  xmin=0, xmax=100, ymin=0, ymax=210,
  xtick={0,20,40,60,80,100},
  ytick={0,50,100,150,200},
  grid=both, grid style={popline!50},
  axis line style={popink},
  legend style={at={(0.97,0.35)}, anchor=east, font=\small}
]
\addplot[very thick, popblue, smooth] coordinates {(2,0) (3,30) (4,60) (6,90) (10,110) (25,130) (55,155) (85,180) (95,200)};
\addlegendentry{\ce{CO2} vs depth}
\addplot[dashed, thick, poporange] coordinates {(0,110) (100,110)};
\node[font=\small, poporange] at (axis cs:70,103) {chemocline};
\end{axis}
\end{tikzpicture}
\end{center}
\legende{Dissolved \ce{CO2} concentration versus depth in a meromictic crater lake (schematic, based on Lake Nyos profiles): concentration is low in the mixed surface layer and rises steeply below the chemocline towards near-saturation at the bottom.}
\end{popfigure}

\courssub{Degassing Engineering at Lake Nyos}

\begin{methode}[Degassing a gas-charged crater lake]
\begin{enumerate}
\item \textbf{Survey:} measure temperature and \ce{CO2} concentration profiles to locate the gas-rich layer and estimate the total gas inventory.
\item \textbf{Install a pipe:} anchor a vertical pipe from a floating raft so its intake sits in the deep, gas-rich layer.
\item \textbf{Prime the flow:} pump bottom water up the pipe briefly. Once rising water decompresses and \ce{CO2} bubbles form, the bubbly water inside the pipe becomes \emph{lighter} than the surrounding lake water, so the flow becomes \emph{self-sustaining} --- a natural fountain (the same principle as a gas-lift pump).
\item \textbf{Vent safely:} the gas escapes to the atmosphere at the surface in small, harmless amounts; the degassed water falls back and, being oxygenated and less gas-charged, gradually reduces the deep inventory.
\item \textbf{Monitor:} repeat profile measurements to confirm that deep \ce{CO2} is falling and that the lake is moving towards stability.
\end{enumerate}
\end{methode}

Degassing pipes installed at Lake Nyos (the first in 2001, with additional pipes later) and at Lake Monoun have substantially lowered the deep-water \ce{CO2} content, reducing --- though not yet eliminating --- the risk of another limnic eruption. A reinforced dam also protects the natural weak outlet wall of Lake Nyos from failure, which would cause both a flood and a sudden depressurisation of the gas-rich layer.

\begin{methode}[Routine monitoring of crater lakes]
\begin{enumerate}
\item Measure the \textbf{gas content} of water samples collected at different depths.
\item Record \textbf{temperature and conductivity profiles} to check the strength of stratification.
\item Track \textbf{stability}: estimate the energy needed to mix the water column; a falling stability margin signals rising danger.
\item Inspect the \textbf{crater walls and outlet} for landslide or erosion hazards that could trigger an overturn.
\end{enumerate}
\end{methode}

\begin{savaistu}[Lake Kivu: the giant cousin of Nyos]
Lake Kivu, between Rwanda and the Democratic Republic of Congo, is a vastly larger meromictic lake containing enormous quantities of dissolved \ce{CO2} \emph{and} methane (\ce{CH4}). A limnic eruption there would threaten millions of people. Engineers now pump the deep water to extract methane for electricity generation --- turning a hazard into a resource. In the examination, comparing Nyos (small, \ce{CO2} only, degassing pipes) with Kivu (large, \ce{CO2} + \ce{CH4}, industrial gas extraction) earns high marks.
\end{savaistu}

\courssub{Lake Chad: a Shallow Endorheic Lake and the Mega-Chad Palaeolake}

Lake Chad, at the northern tip of Cameroon, is the complete opposite of a crater lake: extremely \textbf{shallow} (only a few metres deep on average), very \textbf{broad}, and \cle{endorheic} --- it has no outlet to the sea, so water leaves only by evaporation and seepage.

Its \textbf{hydrological balance} can be written simply:
\[
\Delta S = (P + R) - (E + I)
\]
where $\Delta S$ is the change in stored water, $P$ is precipitation on the lake, $R$ is river inflow (dominated by the Chari--Logone system, with a small contribution from the Komadougou Yobe), $E$ is evaporation (very large under the Sahelian sun) and $I$ is infiltration.

Because $R$ depends on rainfall over the Chari--Logone catchment, any drought shrinks the lake dramatically. Lake Chad has contracted from a much larger area in the 1960s to a small fraction of that in recent decades, through a combination of reduced rainfall, droughts and water abstraction for irrigation. Shorelines have retreated by many kilometres, affecting fishing and farming communities of Cameroon's Far North, Nigeria, Niger and Chad.

Lake sediments and ancient shorelines (beach ridges, wave-cut features) prove that during wetter climatic phases of the Quaternary, the lake swelled into \cle{Mega-Chad}, a palaeolake comparable in area to a small sea. Studying its sediments allows geologists to reconstruct past climate swings of the Sahara--Sahel belt.

\courssub{River-Dammed and Reservoir Lakes: Lagdo, Bamendjing, Mape}

Cameroon has dammed several rivers for hydropower, irrigation and flood control:

\begin{itemize}
\item \textbf{Lagdo Dam} (Benue River, North Region): a large reservoir supplying irrigation water, a fishery and hydroelectric power, and moderating Benue floods.
\item \textbf{Bamendjing and Mape reservoirs} (upper Noun/Mifi system, West Region): regulate flow towards the Mbakaou reservoir and the Songloulou and Edea hydropower schemes on the Sanaga River.
\end{itemize}

Reservoirs, however, age. Two lacustrine processes dominate their evolution:

\begin{itemize}
\item \textbf{Sedimentation:} rivers carry suspended sand, silt and clay into the still water of the reservoir, where the sediment settles and gradually fills the basin, reducing storage capacity year by year.
\item \textbf{Eutrophication:} nutrients (nitrates, phosphates) washed in from farms and villages fertilise the water, causing algal blooms and invasive water weeds; decaying algae consume dissolved oxygen, stressing fish and degrading water quality.
\end{itemize}

\courssub{Lacustrine Processes: Stratification, Turnover, Sedimentation, Productivity}

Let us gather the key processes that operate in \emph{any} lake, using Cameroonian examples:

\begin{enumerate}
\item \textbf{Thermal stratification and turnover.} Deep lakes stratify (epilimnion / thermocline / hypolimnion). Where surface cooling or strong winds eventually mix the whole column, the lake \emph{overturns}, renewing oxygen at depth and releasing accumulated gases. Shallow Lake Chad mixes easily and repeatedly (it is \emph{polymictic}); deep crater lakes like Nyos resist mixing (meromictic) --- with the dangerous consequences seen above.
\item \textbf{Sedimentation.} Lakes trap clastic sediment from their catchments (deltas build into Lake Chad from the Chari), chemical precipitates (carbonates, evaporite salts in endorheic settings), and organic remains (diatom shells, plant debris). The result is a layered, datable archive of environmental change.
\item \textbf{Biological productivity.} Sunlit surface waters support plankton, the base of the food chain. Productivity is controlled by nutrient supply: too little gives clear, poor water; too much (eutrophication) gives blooms and oxygen depletion.
\end{enumerate}

\begin{exoresolu}[Comparing two Cameroonian lakes]
\textbf{Statement.} A candidate states: ``Lake Nyos and Lake Chad are both dangerous because both can overturn.'' Using your knowledge of limnology, (a) explain why this statement is misleading, and (b) identify the real hazard associated with each lake.
\tcblower
\textbf{Solution.}
\begin{enumerate}
\item[(a)] Lake Nyos is deep ($\approx\SI{208}{m}$), meromictic and charged with magmatic \ce{CO2}; its stratification normally \emph{prevents} overturn, and it is precisely the rare, catastrophic breakdown of that stratification (a limnic eruption) that is dangerous. Lake Chad is shallow and polymictic: it mixes frequently and contains no significant deep gas reservoir, so an overturn there is a routine, harmless event. The two lakes therefore cannot be classed together.
\item[(b)] The real hazard at Nyos is a limnic eruption releasing a dense, asphyxiating \ce{CO2} cloud. The real hazards at Lake Chad are hydrological and environmental: shrinkage through drought and abstraction, loss of fisheries and farmland, and the socio-economic stress this causes in the Far North.
\end{enumerate}
\end{exoresolu}

\courssub{Ecological and Economic Roles of Cameroonian Lakes}

\begin{itemize}
\item \textbf{Fisheries:} Lake Chad, Lagdo, Barombi and the crater lakes support important artisanal fisheries that feed local populations.
\item \textbf{Water supply and irrigation:} reservoirs such as Lagdo store water for dry-season irrigation and domestic use; crater lakes supply some highland communities.
\item \textbf{Hydropower regulation:} Bamendjing, Mape and Mbakaou regulate flows feeding the Sanaga power stations, the backbone of Cameroon's electricity supply.
\item \textbf{Tourism:} the twin lakes of Manengouba, Lake Oku and the ring-road crater lakes attract growing eco-tourism.
\item \textbf{Science and conservation:} Lake Barombi is world-famous for its \cle{endemic species} --- cichlid fishes and other organisms found nowhere else on Earth, evolved in isolation within the crater. It is a natural laboratory of evolution, like a miniature ``Galapagos''.
\end{itemize}

\begin{aretenir}[Key points of Section 4]
\begin{itemize}
\item Limnology studies inland waters; Cameroonian lakes are classified by origin: volcanic (crater/maar), tectonic-endorheic, river-dammed and artificial.
\item Nyos and Monoun are meromictic crater lakes charged with magmatic \ce{CO2}; breakdown of stratification caused the limnic eruptions of 1984 and 1986.
\item \ce{CO2} is denser than air: the gas cloud flows down valleys and suffocates silently.
\item Degassing pipes exploit the gas-lift principle to vent deep \ce{CO2} safely; monitoring tracks gas content, temperature profiles and stability.
\item Lake Chad is a shallow endorheic lake whose balance $\Delta S = (P+R)-(E+I)$ explains its shrinkage; Mega-Chad records wetter Quaternary climates.
\item Reservoirs (Lagdo, Bamendjing, Mape) provide power, irrigation and fisheries but suffer sedimentation and eutrophication.
\item Lakes provide fisheries, water, tourism and unique endemic biodiversity (Barombi).
\end{itemize}
\end{aretenir}

\begin{exercice}[Section 4 practice]
\begin{enumerate}
\item Define a meromictic lake and explain why meromixis is a necessary condition for a limnic eruption.
\item With the aid of a labelled diagram, explain how a degassing pipe reduces the hazard at Lake Nyos.
\item Using the water balance equation, explain two reasons for the shrinkage of Lake Chad and one piece of sedimentary evidence for the former Mega-Chad.
\item State one ecological and one economic importance of Lake Barombi.
\end{enumerate}
\end{exercice}

\begin{corrige}[Exercise 1]
A meromictic lake is one whose deep waters never mix with the surface waters. Meromixis is necessary because only permanent stratification allows magmatic \ce{CO2} to accumulate at depth to near-saturation over many years; a lake that overturns regularly releases gas continuously in small, harmless quantities and can never store enough for an eruption.
\end{corrige}

\begin{corrige}[Exercise 2]
The pipe's intake sits in the deep, gas-rich layer. After priming, rising bottom water decompresses, \ce{CO2} bubbles form, and the bubbly column inside the pipe becomes lighter than the surrounding water, so the flow is self-sustaining (gas-lift principle). Gas vents at the surface in small safe amounts while degassed water returns; repeated operation lowers the deep \ce{CO2} inventory and hence the risk of a limnic eruption. (Diagram: stratified lake, pipe from hypolimnion to surface fountain on a raft.)
\end{corrige}

\begin{corrige}[Exercise 3]
Reasons for shrinkage: (i) reduced rainfall and drought over the Chari--Logone catchment, lowering inflow $R$; (ii) high evaporation $E$ plus increasing abstraction of water for irrigation. Evidence for Mega-Chad: ancient beach ridges and shoreline terraces far above the present lake, and extensive lacustrine sediments (diatomites, lake muds) across the basin.
\end{corrige}

\begin{corrige}[Exercise 4]
Ecological: Lake Barombi hosts endemic cichlid fishes found nowhere else, making it a unique site of evolution and conservation. Economic: it supports an artisanal fishery and local tourism around Kumba.
\end{corrige}

\coursec{Section 5: Geologic Structures of Cameroon (Lessons 140-141)}

In the previous section we studied Cameroon's lakes and lacustrine processes. We saw that the crater lakes of the Western Highlands (Nyos, Monoun, Oku) sit on volcanic structures, while Lakes Barombi Mbo and Barombi Koto occupy volcanic maars. But why are these lakes, volcanoes and basins located exactly where they are? The answer lies in the \cle{structural framework} of the country: the faults, shear zones, folds and joints that dissect the Cameroonian crust. In this section we move from surface processes to the deep architecture of Cameroon. We will see that nearly everything we have studied so far -- the alignment of the Cameroon Volcanic Line, the shape of the sedimentary basins, the seismicity of the Mount Cameroon region -- is controlled by \cle{geologic structures}.

\courssub{The Structural Framework of Cameroon: Cratonic South versus Mobile Belt North}

A useful first observation comes from a simple relief map of Cameroon. The southern half of the country is a vast, low, gently undulating plateau covered by dense equatorial forest and deeply weathered rocks. The northern half, by contrast, is cut by mountain chains (the Adamawa Plateau margins, the Mandara Mountains) and by elongate basins. This contrast in the landscape reflects a fundamental contrast in the \emph{structural behaviour} of the crust.

\begin{definition}[Craton and mobile belt]
A \cle{craton} is an old, cold, thick and rigid part of the continental crust that has remained tectonically stable (undeformed) for a very long time -- typically more than a billion years. A \cle{mobile belt} (or orogenic belt) is a zone of crust that has been weakened, heated, deformed and metamorphosed during a mountain-building episode (orogeny); it remains mechanically weaker than a craton and can be reactivated later.
\end{definition}

Cameroon is divided into two great structural domains:

\begin{enumerate}
\item \textbf{The Congo Craton (southern Cameroon).} South of roughly the latitude of Yaound\'e--Bertoua (approximately along the Sanaga River corridor), the basement belongs to the northern edge of the \cle{Congo Craton}. It is made of Archaean to Palaeoproterozoic rocks (older than about 2.0 Ga): greenstone belts, granites, gneisses and charnockites of the Ntem Complex (also called the Ntem Series). This crust is rigid, poorly foliated, and deformed only by ancient structures. It forms the stable ``foundation'' on which the rest of the country was built.
\item \textbf{The Pan-African mobile belt (central and northern Cameroon).} North of the craton boundary, the basement consists of Neoproterozoic rocks (roughly 800--550 Ma) that were deformed, metamorphosed and intruded by granites during the \cle{Pan-African orogeny}, the great collision event that assembled the supercontinent Gondwana. This belt is intensely foliated, folded and cut by major shear zones. Key lithological units include the Poli, Lom and Yaound\'e series.
\end{enumerate}

\begin{aretenir}[The two structural domains]
Southern Cameroon = rigid \cle{Congo Craton} (Archaean--Palaeoproterozoic, stable). Central and northern Cameroon = \cle{Pan-African mobile belt} (Neoproterozoic, deformed, foliated, cut by shear zones). The boundary between them is itself a major structural feature, close to the Sanaga Fault zone.
\end{aretenir}

\begin{savaistu}[Why this division matters for everyday life]
The rigid craton of the south is a poor seismic zone but hosts the iron deposits of the south-east (Mbalam region) in its ancient greenstone belts. The weaker mobile belt of the north is where faults were later reactivated, guiding both the volcanism of the Cameroon Volcanic Line and the opening of sedimentary basins. Structure controls resources, hazards and even the road network.
\end{savaistu}

\courssub{The Great Shear Zones: CCSZ, Foumban and Tchollire-Banyo}

If you look at a geological map of central Cameroon, you notice that the foliation of the gneisses and the outlines of the granite plutons are drawn out into long, parallel, NE--SW trending bands. These are not ordinary faults but \cle{shear zones}.

\begin{definition}[Shear zone]
A \cle{shear zone} is a broad zone of intense ductile deformation, from a few hundred metres to several kilometres wide, along which crustal blocks slide past each other. Unlike a narrow brittle fault, a shear zone deforms rocks plastically at depth: minerals are stretched and recrystallised, producing strongly foliated rocks called \cle{mylonites}. When the movement is horizontal (blocks sliding laterally), we speak of a \cle{strike-slip} shear zone.
\end{definition}

The Pan-African belt of Cameroon is cut by a family of great NE--SW strike-slip shear zones, formed about 600--550 Ma during the final stages of the Pan-African collision:

\begin{itemize}
\item The \cle{Central Cameroon Shear Zone (CCSZ)}, the most important, running NE--SW across the Adamawa region. It is a dextral (right-lateral) strike-slip zone: a block on the north-west side has moved towards the north-east relative to the south-east side. Width up to 10 km in places.
\item The \cle{Foumban shear zone}, a related branch passing near Foumban in the West Region, often considered a splay of the CCSZ.
\item The \cle{Tchollire-Banyo shear zone}, further north, linking Cameroon to the shear systems of Nigeria (e.g. the Kaltungo shear zone) and Chad.
\end{itemize}

These zones belong to a continental-scale network of Pan-African strike-slip faults that continues into north-eastern Brazil (the Pernambuco and Patos shear zones) -- direct evidence that Cameroon and Brazil were joined before the opening of the South Atlantic.

\begin{propriete}[Structural control of magmatism]
The volcanoes of the \cle{Cameroon Volcanic Line (CVL)} are aligned along, or at the intersections of, \emph{reactivated Pan-African fractures}, particularly the CCSZ and the Sanaga Fault. This demonstrates the principle of \cle{structural control of magmatism}: magma rises preferentially along pre-existing zones of crustal weakness rather than piercing intact crust. (Proof not examinable; the map evidence is examinable.)
\end{propriete}

\courssub{The Sanaga Fault: An Ancient Lineament Reactivated}

The \cle{Sanaga Fault} is a major ENE--WSW to E--W lineament crossing central Cameroon, roughly along the course of the Sanaga River between Nanga-Eboko and Ed\'ea. It originated as a Pan-African structure near the boundary between the Congo Craton and the mobile belt, and it was \cle{reactivated} during the Cretaceous and Cenozoic.

Why is reactivation so important? Old faults are permanent scars in the crust. When a new stress field is applied -- for example the stresses related to the opening of the South Atlantic (Cretaceous) and the uplift of the Adamawa Plateau (Cenozoic) -- the crust breaks again along the old scars rather than creating new ones. The Sanaga Fault thus became a pathway for Cenozoic magmas: several volcanic centres of the continental CVL (including the Manengouba--Bambouto sector and the Mount Cameroon area) lie close to the intersection of the Sanaga trend with the NE--SW shear zones. The straight lower course of the Sanaga River itself is guided by the fault -- an example of \cle{structural control of drainage}.

\begin{attention}[A lineament is not automatically a fault]
A common exam mistake is to identify every straight line on a satellite image or map as a fault. A \cle{lineament} is simply a linear feature (a straight valley, an aligned series of hills, a vegetation boundary). It \emph{may} mark a fault, but it can also be a dyke, a joint set, a lithological contact or even a human feature. To call it a fault you need \textbf{field evidence}: offset of rock units, fault breccia, slickensides, mylonites, or aligned springs and seismicity. Always write: ``the lineament \emph{suggests} a fault, to be confirmed in the field''.
\end{attention}

\courssub{Faults Controlling the Sedimentary Basins}

The sedimentary basins studied in Section 3 did not open randomly: they are bounded and shaped by faults.

\begin{itemize}
\item \textbf{The Mamfe Basin} (South West Region) is a small Cretaceous rift basin, an offshoot of the Benue Trough. It is bounded by \cle{normal faults} along which the basin floor subsided while sediments (sandstones, shales, evaporites) accumulated. Its rhombic shape is typical of a \emph{pull-apart} geometry along strike-slip faults.
\item \textbf{The Benue Trough} (extending into northern Cameroon as the Garoua and Yola arms) is a major Cretaceous rift system controlled by large boundary normal faults. Subsidence along these faults created the space for thousands of metres of marine and continental Cretaceous sediments.
\item \textbf{The Douala and Rio del Rey basins} (Atlantic coast) are passive-margin and deltaic basins whose structure is dominated by listric (curved) normal faults linked to sediment loading and gravity sliding toward the ocean.
\end{itemize}

\begin{definition}[Rift and boundary fault]
A \cle{rift} is an elongate depression formed where the crust is stretched and thinned. Extension is taken up by \cle{normal faults}: faults along which the hanging wall moves down relative to the footwall. The largest normal faults at the edge of a rift are called \cle{boundary faults}; they control the depth of the basin and the thickness of its sedimentary fill.
\end{definition}

\begin{popfigure}
\begin{center}
\begin{tikzpicture}[scale=1.0]
% Cameroon outline (stylised)
\draw[thick, popink] (0,0) -- (0.6,2.2) -- (1.2,3.4) -- (1.0,4.6) -- (1.8,5.6) -- (2.6,6.4) -- (3.4,6.9) -- (4.4,6.6) -- (5.6,6.2) -- (6.4,5.2) -- (6.2,4.0) -- (5.6,3.0) -- (5.2,1.8) -- (4.0,0.8) -- (2.2,0.2) -- cycle;
% Congo craton
\fill[popgreenL, opacity=0.5] (0.4,0.3) -- (0.6,2.0) -- (2.6,2.6) -- (4.2,1.6) -- (4.0,0.8) -- (2.2,0.2) -- cycle;
\node[popdark, align=center, text width=2.2cm] at (2.0,1.2) {\footnotesize Congo Craton (S)};
% Mobile belt label
\node[popdark, align=center, text width=2.6cm] at (4.6,5.4) {\footnotesize Pan-African mobile belt (N)};
% CCSZ
\draw[very thick, poporange] (1.6,2.4) -- (5.9,5.0);
\node[poporange, rotate=24, align=center] at (4.0,4.15) {\footnotesize CCSZ};
% Foumban shear
\draw[thick, poporange, dashed] (1.4,3.0) -- (4.6,4.6);
\node[poporange, rotate=24] at (2.6,4.15) {\footnotesize Foumban SZ};
% Tchollire-Banyo
\draw[thick, poporange, dashed] (3.2,5.6) -- (6.2,6.3);
\node[poporange, rotate=8] at (4.9,6.35) {\footnotesize Tchollire-Banyo SZ};
% Sanaga fault
\draw[very thick, poppurple] (0.7,2.9) -- (4.6,2.4);
\node[poppurple] at (2.6,2.85) {\footnotesize Sanaga Fault};
% CVL
\draw[very thick, popblue] (0.15,1.9) -- (2.2,3.5);
\foreach \p in {(0.3,2.05),(0.8,2.4),(1.3,2.75),(1.8,3.1),(2.1,3.4)} \node[popblue] at \p {\scriptsize $\blacktriangle$};
\node[popblue, rotate=35, align=center] at (0.85,3.15) {\footnotesize CVL};
% Mamfe basin
\draw[thick, popgreen] (0.9,2.9) ellipse (0.5 and 0.3);
\node[popgreen!60!black] at (0.95,3.35) {\footnotesize Mamfe};
% Benue / Garoua basin
\draw[thick, popgreen] (4.6,5.9) ellipse (0.7 and 0.35);
\node[popgreen!60!black] at (4.7,5.55) {\footnotesize Garoua (Benue)};
% Douala basin
\draw[thick, popgreen] (0.35,1.55) ellipse (0.42 and 0.28);
\node[popgreen!60!black, align=center, text width=1.3cm] at (-0.35,1.15) {\footnotesize Douala};
% Mount Cameroon
\node[popblue] at (0.35,2.1) {\scriptsize $\blacktriangle$};
\node[popink, align=center, text width=1.6cm] at (-0.45,2.35) {\footnotesize Mt Cameroon};
% Rio del Rey
\draw[thick, popgreen] (0.1,1.0) ellipse (0.35 and 0.2);
\node[popgreen!60!black] at (-0.5,0.7) {\footnotesize Rio del Rey};
% legend
\node[align=left, text width=4.2cm, anchor=west] at (6.7,4.2) {\footnotesize \textcolor{poporange}{\rule{0.5cm}{2pt}} shear zones\\[2pt] \textcolor{poppurple}{\rule{0.5cm}{2pt}} Sanaga Fault\\[2pt] \textcolor{popblue}{$\blacktriangle$} CVL volcanoes\\[2pt] \textcolor{popgreen!60!black}{\rule{0.4cm}{1.5pt}} sedimentary basins};
\end{tikzpicture}
\end{center}
\legende{Simplified structural map of Cameroon: the Central Cameroon Shear Zone (CCSZ), Foumban and Tchollire-Banyo shear zones, the Sanaga Fault, the Cameroon Volcanic Line and the fault-bounded sedimentary basins.}
\end{popfigure}

\courssub{Fault Types and Their Cameroonian Examples}

Three fundamental fault geometries occur in Cameroon. Recognising them on maps and in the field is a core GCE skill.

\begin{popfigure}
\begin{center}
\begin{tikzpicture}[scale=0.85]
% Normal fault
\begin{scope}
\fill[popblueL] (0,0) rectangle (3.4,1.4);
\fill[popgreenL] (0,1.4) rectangle (3.4,2.2);
\draw[thick, popink] (1.7,2.6) -- (2.1,0);
\draw[thick] (0,1.4) -- (1.78,1.4);
\draw[thick] (2.02,1.4) -- (3.4,1.4);
\draw[thick] (0,2.2) -- (1.72,2.2);
\draw[thick] (2.06,2.2) -- (3.4,2.2);
\draw[->, thick, popdark] (1.35,2.0) -- (1.35,1.2);
\draw[->, thick, popdark] (2.45,0.9) -- (2.45,1.7);
\node[align=center] at (1.7,-0.5) {\footnotesize \textbf{Normal fault}\\ \footnotesize Mamfe \& Benue rifts};
\end{scope}
% Reverse fault
\begin{scope}[xshift=4.6cm]
\fill[popblueL] (0,0) rectangle (3.4,1.4);
\fill[popgreenL] (0,1.4) rectangle (3.4,2.2);
\draw[thick, popink] (2.1,2.6) -- (1.7,0);
\draw[thick] (0,1.4) -- (1.62,1.4);
\draw[thick] (1.86,1.4) -- (3.4,1.4);
\draw[thick] (0,2.2) -- (1.78,2.2);
\draw[thick] (1.98,2.2) -- (3.4,2.2);
\draw[->, thick, popdark] (1.3,1.1) -- (1.3,1.9);
\draw[->, thick, popdark] (2.5,1.9) -- (2.5,1.1);
\node[align=center] at (1.7,-0.5) {\footnotesize \textbf{Reverse fault}\\ \footnotesize Pan-African thrusts, N Cameroon};
\end{scope}
% Strike-slip
\begin{scope}[xshift=9.2cm]
\fill[popblueL] (0,0) rectangle (3.4,1.1);
\fill[popgreenL] (0,1.1) rectangle (3.4,2.2);
\draw[thick, popink] (1.7,0) -- (1.7,2.2);
\draw[->, thick, popdark] (0.4,1.75) -- (1.4,1.75);
\draw[->, thick, popdark] (3.0,0.55) -- (2.0,0.55);
\draw[thick] (0,1.1) -- (1.7,1.1);
\draw[thick] (1.7,1.1) -- (3.4,1.1);
\node[align=center] at (1.7,-0.5) {\footnotesize \textbf{Strike-slip fault}\\ \footnotesize CCSZ, Sanaga (map view)};
\end{scope}
\end{tikzpicture}
\end{center}
\legende{Block diagrams of the three fault types found in Cameroon. Normal fault: hanging wall moves down (extension, rift basins). Reverse fault: hanging wall moves up (compression, Pan-African belt). Strike-slip fault: horizontal sliding (shear zones).}
\end{popfigure}

\begin{exemplebox}[Identifying a fault type in the field]
Near the margin of the Mamfe Basin, a geologist finds a road cut where red Cretaceous sandstones lie against grey Pan-African gneisses, with a zone of crushed rock (fault breccia) between them, and the sandstone beds dip steeply toward the basin centre. The young sediments are \emph{downthrown} against the old basement: this is a \cle{normal boundary fault}, consistent with the extensional origin of the basin.
\end{exemplebox}

\courssub{Folds in the Pan-African Belt and in the Cretaceous Basins}

\cle{Folds} are bends in layered rocks produced by ductile deformation under compression. Two families of folds are distinguished in Cameroon:

\begin{enumerate}
\item \textbf{Pan-African folds} affect the metamorphic basement of central and northern Cameroon. They are typically \emph{tight} to \emph{isoclinal} (limbs nearly parallel), with axes trending N--S to NNE--SSW, often refolded by later deformation. On the map they appear as repeated, elongated bands of gneiss, schist and quartzite -- for example in the Poli and Lom series of the North Region.
\item \textbf{Cretaceous basin folds} affect the sedimentary fill of the Benue Trough. During the Santonian compressional event (about 85 Ma), the trough sediments were folded into broad, open \emph{anticlines and synclines} with axes parallel to the trough (roughly NE--SW). These folds are important because they can trap hydrocarbons, as in the neighbouring Benue Trough of Nigeria.
\end{enumerate}

On a geological map, an anticline shows the \emph{oldest} beds in the centre with dips directed \emph{away} from the axis; a syncline shows the \emph{youngest} beds in the centre with dips directed \emph{toward} the axis.

\courssub{Joints, Lineaments, Drainage and Mineralisation}

\cle{Joints} are fractures along which there has been no visible displacement. They form in sets (parallel families) and are ubiquitous in both the basement and the sedimentary cover. Their importance in Cameroon is threefold:

\begin{itemize}
\item \textbf{Drainage control.} Rivers preferentially follow joints and faults because fractured rock erodes faster. The rectangular (right-angled) drainage pattern of many streams on the Adamawa Plateau, and the straight course of the lower Sanaga, reflect joint and fault control.
\item \textbf{Weathering and landforms.} Jointing guides the spheroidal weathering of granites, producing the inselbergs and castle koppies typical of the Centre and North Regions.
\item \textbf{Mineralisation.} Joints and shear zones channel hydrothermal fluids. Gold mineralisation in the East Region (Batouri, B\'etar\'e Oya) is associated with quartz veins emplaced along fractures in and near Pan-African shear zones -- a direct link between structure and economic geology that we will exploit in Section 6.
\end{itemize}

\courssub{Neotectonics: The Structure of Cameroon Is Still Active}

\cle{Neotectonics} is the study of recent and ongoing crustal movements. Cameroon is not a dead piece of crust:

\begin{itemize}
\item The Mount Cameroon region is one of the most seismically active zones in West Africa. Earthquakes there are linked to \emph{reactivation of faults} of the CVL and to magmatic pressure beneath the volcano. The eruptions of Mount Cameroon (e.g. 1999, 2000) were accompanied by earthquake swarms and ground fissuring along the volcano's fracture zones.
\item Historical seismicity is also recorded along the Sanaga Fault zone and around the Adamawa shear zones, showing that the Pan-African scars remain mechanically active.
\item The 1986 Lake Nyos disaster (Section 4) reminds us that the deep CO$_2$ feeding the crater lakes rises along faults and fractures -- another expression of active structure.
\end{itemize}

\begin{attention}[Exam tip: one concept, two hazards]
Examiners love the link: \textbf{fault reactivation} $\rightarrow$ (1) \emph{volcanism}, because fractures give magma a pathway (CVL alignment), and (2) \emph{seismicity}, because reactivated faults slip and generate earthquakes (Mount Cameroon region). Whenever a question mentions CVL structure, mention both consequences -- it shows integrated understanding.
\end{attention}

\courssub{Reading Structural Data: Strike, Dip and Map Interpretation}

All the structures above are described quantitatively using two measurements taken on any planar feature (bedding, foliation, fault plane, joint).

\begin{definition}[Strike and dip]
The \cle{strike} of a plane is the direction (bearing) of the horizontal line contained in that plane -- the line along which the plane intersects a horizontal surface. The \cle{dip} is the angle between the plane and the horizontal, measured \emph{perpendicular to the strike}, in the direction of maximum slope. A measurement is written, for example, as ``strike N40$^\circ$E, dip 30$^\circ$ SE'' or in quadrant notation ``040/30 SE''.
\end{definition}

\begin{methode}[Measuring strike and dip with a compass-clinometer]
\begin{enumerate}
\item \textbf{Expose the plane.} Find a clean, representative surface of bedding or foliation on the outcrop; avoid loose or rotated blocks.
\item \textbf{Measure the strike.} Place the long edge of the compass flat on the plane and level it (bubble centred). Read the bearing of the horizontal edge: this is the strike.
\item \textbf{Measure the dip.} Turn the clinometer 90$^\circ$ from the strike direction, place it on the plane in the direction of steepest slope, and read the angle: this is the dip. Note the dip \emph{direction} (the quadrant toward which the plane slopes).
\item \textbf{Record and plot.} Write the data in your notebook (e.g. 040/30 SE). On the map, draw the standard symbol: a long line in the strike direction with a short tick on the dip side, and write the dip value beside the tick.
\item \textbf{Check consistency.} Repeat at several stations; consistent strikes with changing dips suggest a fold, while abrupt changes in rock type across a line suggest a fault.
\end{enumerate}
\end{methode}

\begin{popfigure}
\begin{center}
\begin{tikzpicture}[scale=1.0]
% outcrop plane (parallelogram)
\fill[popblueL] (0,0) -- (4.2,0) -- (5.0,1.6) -- (0.8,1.6) -- cycle;
\draw[thick, popink] (0,0) -- (4.2,0) -- (5.0,1.6) -- (0.8,1.6) -- cycle;
% horizontal reference plane (dashed)
\draw[dashed, popmuted] (0,0) -- (4.2,0);
\draw[dashed, popmuted] (4.2,0) -- (5.0,0.9);
% strike line
\draw[very thick, poporange] (0.6,0.55) -- (4.6,0.55);
\node[poporange] at (2.6,0.85) {\footnotesize strike line (horizontal)};
% dip arrow
\draw[->, very thick, poppurple] (2.6,0.55) -- (2.95,1.35);
\node[poppurple, align=center] at (3.6,1.25) {\footnotesize dip direction};
% dip angle arc
\draw[popdark] (3.4,0.55) arc (0:33:0.8);
\node[popdark] at (3.75,0.72) {\footnotesize $\delta$ = dip angle};
% compass
\draw[fill=white] (1.3,0.55) circle (0.28);
\node at (1.3,0.55) {\scriptsize N};
% map symbol inset
\begin{scope}[xshift=7.2cm, yshift=0.4cm]
\draw[thick, popink] (0,0.6) -- (2.2,0.6);
\draw[thick, popink] (1.1,0.6) -- (1.1,0.1);
\node at (1.35,0.15) {\footnotesize 30};
\node[align=center, text width=3.2cm] at (1.1,-0.55) {\footnotesize map symbol: strike N40$^\circ$E, dip 30$^\circ$ SE};
\end{scope}
\end{tikzpicture}
\end{center}
\legende{Strike and dip on an outcrop: the strike is the horizontal line on the plane; the dip is the maximum slope angle measured perpendicular to strike. Right: the standard map symbol (long line = strike, tick = dip direction, number = dip value).}
\end{popfigure}

\begin{exoresolu}[Interpreting structural data from a traverse]
\textbf{Statement.} During a field traverse across a Cretaceous outcrop near Garoua, a student records the following bedding attitudes along an E--W line: at station A, 090/25 N; at station B, 090/25 N; at station C, 090/25 S; at station D, 090/25 S. Fossils show that the beds at A are older than those at B, and those at D are older than those at C. (a) What structure is present? (b) Where is its axis? (c) What is the regional significance of such folds?
\tcblower
\textbf{Solution.}
\begin{enumerate}
\item All strikes are identical (E--W, 090), so the beds form a single folded series. Between B and C the dip direction reverses from north to south: the beds dip \emph{away} from a central zone located between B and C.
\item Beds dipping away from a central axis define an \cle{anticline}. The fossil evidence confirms it: the oldest beds (A, D) are on the flanks and the structure closes outward -- wait: oldest beds at A (west) and D (east) with younger beds at B and C near the axis would indicate a \emph{syncline}. Re-reading: A older than B on the western limb, D older than C on the eastern limb, so the youngest beds are at B and C, nearest the axis. Youngest beds in the centre $+$ dips toward the axis (B dips N, C dips S, i.e. toward each other) $=$ \cle{syncline}.
\item The axis lies between stations B and C, trending E--W (parallel to the strike, 090).
\item Such open folds in the Benue Trough fill formed during the Santonian compressional event; synclines preserve the youngest Cretaceous sediments, while the associated anticlines are potential structural traps for hydrocarbons.
\end{enumerate}
\textbf{Answer:} (a) an E--W trending syncline; (b) axis between B and C; (c) Santonian folding of the Benue Trough, important for hydrocarbon trapping.
\end{exoresolu}

\begin{exercice}[Fault recognition on a sketch map]
On a map of the Mamfe region, granite gneiss of the basement is in contact with Cretaceous sandstone along a straight, 15 km long boundary. Along this line you find crushed rock, aligned springs, and the sandstone is repeated nowhere else. (a) What type of structure is the boundary? Justify with two pieces of evidence. (b) Is it more likely normal or reverse, given that the Mamfe Basin is a rift? Explain. (c) Why would a straight river segment along this boundary support your interpretation, and why is it not sufficient on its own?
\end{exercice}

\begin{corrige}[Exercise: Fault recognition on a sketch map]
\begin{enumerate}
\item The boundary is a \cle{fault}. Evidence: (i) crushed rock (fault breccia) along the contact; (ii) aligned springs, which rise along impermeable fault zones; the straightness of the contact over 15 km and the abrupt juxtaposition of unrelated rock units (basement against young sediment) are additional arguments.
\item It is most likely a \cle{normal fault}: the Mamfe Basin is an extensional rift, and rift boundary faults are normal faults along which the sedimentary hanging wall subsided against the basement footwall.
\item A straight river segment indicates \emph{structural control of drainage}: running water exploits the weak, fractured fault zone. It supports the fault interpretation but is not sufficient alone, because a lineament can also be a dyke, a joint set or a lithological contact -- only field evidence (breccia, offset, springs) proves faulting.
\end{enumerate}
\end{corrige}

\begin{aretenir}[Key points of Section 5]
\begin{itemize}
\item Cameroon = rigid \cle{Congo Craton} in the south $+$ deformed \cle{Pan-African mobile belt} in the centre and north.
\item The \cle{CCSZ}, Foumban and Tchollire-Banyo shear zones are Pan-African strike-slip zones; the \cle{Sanaga Fault} is a reactivated lineament near the craton boundary.
\item \cle{Structural control}: CVL volcanoes, sedimentary basins (Mamfe, Benue), drainage and gold mineralisation all follow ancient fractures.
\item Fault types: \cle{normal} (rifts), \cle{reverse} (Pan-African compression), \cle{strike-slip} (shear zones). Neotectonic reactivation explains both CVL volcanism and Mount Cameroon seismicity.
\item Field skills: measure \cle{strike and dip} with a compass-clinometer; a lineament is only a \emph{suspected} fault until field evidence confirms it.
\end{itemize}
\end{aretenir}

With the structural architecture of Cameroon now in place, we are ready to see how these structures have localised the country's mineral wealth -- the subject of Section 6, Economic Minerals and Natural Resources of Cameroon.

\coursec{Section 6: Economic Minerals and Natural Resources of Cameroon (Lessons 142-144)}

\courssub{Introduction: From Structure to Resources}

\begin{definition}[Economic Mineral]
A mineral or rock that can be extracted, processed, and marketed at a profit. The classification depends on geological occurrence, concentration (grade), tonnage, accessibility, and market conditions.
\end{definition}

\begin{aretenir}[Key Controls on Mineralisation in Cameroon]
\begin{itemize}
\item \textbf{Geological Setting:} The Cameroon Volcanic Line (CVL) for magmatic/hydrothermal deposits; the Congo Craton basement for gold, iron, diamonds; the Benue Trough / Coastal Basin for sedimentary resources (limestone, salt, hydrocarbons).
\item \textbf{Tectonics:} Pan-African shear zones (e.g., Central Cameroon Shear Zone) act as fluid pathways for gold and base metals. Rift structures control hydrocarbon traps and volcanic vents.
\item \textbf{Climate/Weathering:} Tropical lateritisation forms bauxite, nickel laterite, and supergene gold enrichment.
\end{itemize}
\end{aretenir}

\coursec{Section 7: Stratigraphy of Cameroon (Lesson 145)}

In Section 6 we catalogued the mineral wealth of Cameroon --- bauxite at Minim-Martap, iron at Mbalam, gold in the eastern greenstone belts, petroleum in the Douala and Rio del Rey basins. You may have noticed a pattern: each resource belongs to a particular family of rocks formed at a particular time. Gold sits in Archaean greenstones, petroleum in Cretaceous--Cenozoic basin fill, bauxite on Cenozoic volcanic plateaux. To organise all of this knowledge into a single coherent picture, geologists use \cle{stratigraphy}: the science that arranges rock units in their order of formation. This section builds the synthetic stratigraphic column of Cameroon, from the ancient Ntem Complex to the youngest ashes of Mount Cameroon, and shows how the gaps in that column are just as informative as the rocks themselves.

\courssub{Principles: superposition, unconformities, correlation and the geologic time scale}

Imagine the pages of a book in which the history of Cameroon is written. Each page is a rock layer; the oldest pages lie at the bottom, the youngest at the top. But this book has been torn, burnt and partly erased: whole chapters (the Palaeozoic, for instance) are missing. Stratigraphy is the art of reading what remains and honestly acknowledging what is lost.

\begin{definition}[Stratigraphic column]
A \cle{stratigraphic column} is a vertical diagram showing the succession of rock units in a region, from the oldest at the bottom to the youngest at the top, with their lithologies, ages, thicknesses and the events that affected them.
\end{definition}

Four principles allow us to construct such a column.

\begin{enumerate}
\item \textbf{Superposition.} In an undisturbed sedimentary succession, each bed is younger than the one beneath it. This is the oldest and most fundamental rule of stratigraphy.
\item \textbf{Unconformities.} Where deposition stopped and erosion planed the land before sedimentation resumed, the contact between the old and new rocks is an \cle{unconformity}. It marks a gap in time.
\item \textbf{Correlation.} Rocks of the same age in different places are matched using fossils, lithological similarity and radiometric dates.
\item \textbf{The geologic time scale.} All dated units are placed within the standard framework of eons, eras and periods (Archaean, Proterozoic, Palaeozoic, Mesozoic, Cenozoic).
\end{enumerate}

\begin{propriete}[Meaning of an unconformity]
An \cle{unconformity} represents a gap in the rock record, corresponding to a period of erosion or of non-deposition. The time interval missing at an unconformity is called a \cle{hiatus}. In Cameroon, the greatest hiatus of all separates the Precambrian basement from the Cretaceous cover: nearly 500 million years with almost no preserved rock.
\end{propriete}

\begin{methode}[Correlating units between basins]
To show that, for example, the Cretaceous sandstones of the Mamfe Basin are the same age as those of the Benue Trough, proceed as follows.
\begin{enumerate}
\item Compare \textbf{lithology}: similar rock types in similar successions suggest the same depositional episode.
\item Look for \textbf{fossils}: identical fossil assemblages (for instance Cretaceous ammonites, bivalves or pollen) indicate the same age.
\item Use \textbf{radiometric ages} where available: dated volcanic or intrusive rocks bracket the age of the sediments they cut or underlie.
\item Check the \textbf{structural setting}: units deposited in the same rift event share the same relation to faults and unconformities.
\item Conclude with a correlation chart, drawing horizontal time-lines between the columns of each basin.
\end{enumerate}
\end{methode}

\courssub{The Archaean foundation: the Ntem Complex ($>2.5$ Ga)}

The base of the Cameroonian column is the \cle{Ntem Complex} of the South Region, the northwestern edge of the great Congo craton. It consists of grey gneisses (tonalitic--trondhjemitic gneisses), charnockites and granulites, all older than about $2.5$ billion years. These rocks formed deep in the roots of an ancient mountain chain and were later uplifted and exposed.

Within and around the Ntem Complex occur \cle{greenstone belts} --- elongated strips of metamorphosed volcanic and sedimentary rocks, typically rich in iron formations and gold-bearing quartz veins. The Batouri and Bétaré Oya regions of the East Region host such belts; they explain the artisanal gold workings we met in Section 6.

\begin{aretenir}[The Archaean record]
The Ntem Complex ($>2.5$ Ga) is the oldest known rock unit in Cameroon. Together with its greenstone belts it forms the Archaean basement onto which all younger units were deposited or accreted.
\end{aretenir}

\courssub{The Palaeoproterozoic: the Nyong Group ($\sim 2.1$ Ga)}

Resting on and around the Archaean core, the \cle{Nyong Group} of the Centre--South regions records a Palaeoproterozoic cycle of sedimentation and mountain building at about $2.1$ Ga. It comprises metasediments --- micaschists, quartzites, itabirites (banded iron formations) and marbles --- deposited in a basin at the margin of the Congo craton and then metamorphosed during the Eburnean orogeny. The iron formations of the Nyong Group are the deep roots of the great iron districts (Mbalam, Nkout) discussed in Section 6.

\courssub{The Neoproterozoic: the Pan-African belt ($\sim 800$--$550$ Ma)}

The largest part of Cameroon --- the Centre, Adamawa, North and Far North regions --- is built of \cle{Pan-African} rocks. Between roughly $800$ and $550$ Ma, a major orogeny deformed and metamorphosed sedimentary and volcanic sequences into schists, gneisses and migmatites, while large volumes of \cle{syn-tectonic granite} were intruded into the rising mountain belt. Units such as the Lom, Mbalmayo--Bengbis and Yaoundé groups (metasediments, including the famous Yaoundé garnet gneisses) belong to this belt. The Pan-African belt is a \emph{mobile belt} caught between the Congo craton to the south and the West African craton to the north.

\begin{savaistu}[Three basements in one country]
Cameroon is one of the few African countries where you can travel in a single day from an Archaean craton (Ntem), across a Palaeoproterozoic belt (Nyong), into a Neoproterozoic mobile belt (Pan-African). This triple basement controls everything else in the country's geology: where basins later opened, where faults cut, and where minerals are found.
\end{savaistu}

\courssub{The Palaeozoic gap: a missing chapter}

Between the end of the Pan-African orogeny (about $550$ Ma) and the Cretaceous (about $145$ Ma), Cameroon has almost \emph{no preserved rock record}. No Cambrian limestones, no Devonian shales, no Carboniferous coals are known in the country.

\begin{attention}[The Palaeozoic gap does not mean ``nothing happened'']
A very common mistake is to write that ``during the Palaeozoic, nothing happened in Cameroon''. This is wrong. The absence of Palaeozoic rocks means there is \textbf{no preserved record}: either the region stood high as an eroding landmass (non-deposition), or Palaeozoic sediments were deposited and later completely eroded away. Tectonic events certainly occurred --- but the rocks that would testify to them have been destroyed or were never laid down. In stratigraphy, absence of evidence is not evidence of absence of events.
\end{attention}

\courssub{The Mesozoic record: Cretaceous rift and marine sediments}

The record resumes in the \cle{Cretaceous}, when the break-up of Gondwana opened rift basins across Cameroon:

\begin{itemize}
\item The \textbf{Benue Trough} (North Region, around Garoua): continental sandstones and conglomerates passing up into marine shales and limestones deposited during a Cretaceous marine transgression from the south Atlantic.
\item The \textbf{Mamfe Basin} (South West Region): a smaller rift basin with continental clastic sediments, famous for dinosaur and other vertebrate remains and for its evaporites (salt springs).
\item The \textbf{Douala Basin} (Littoral): marine and deltaic Cretaceous sediments that thicken seaward and form the deeper part of the coastal petroleum basin.
\end{itemize}

These basins record the birth of the South Atlantic Ocean: first continental rifting, then shallow seas as the ocean widened.

\courssub{The Cenozoic record: coastal basins, the Chad Basin and the CVL}

The Cenozoic column of Cameroon has three distinct expressions.

\begin{enumerate}
\item \textbf{Coastal marine and deltaic sequences.} In the Douala and Rio del Rey basins, Cenozoic shales, sandstones and deltaic deposits accumulated as the Niger-derived and Sanaga-derived deltas prograded into the Gulf of Guinea. These beds host the country's petroleum.
\item \textbf{Chad Basin fluvio-lacustrine beds.} In the Far North, around Kousséri and the Logone floodplain, Neogene to Quaternary river and lake sediments (sands, clays, diatomites) record the fluctuations of ancient Lake Chad.
\item \textbf{CVL volcanic rocks.} From about $65$ Ma to the present, the Cameroon Volcanic Line poured out basalts, trachytes and phonolites, building Mount Cameroon, the Manengouba and Bambouto calderas, and the Adamawa Plateau. These are the youngest rocks in the column --- Mount Cameroon last erupted in the years 1999--2000.
\end{enumerate}

\courssub{The synthetic stratigraphic column of Cameroon}

We can now assemble the full column. Read it from the bottom up, as a geologist reads time.

\begin{popfigure}
\begin{center}
\begin{tikzpicture}[scale=0.92]
% column background
\fill[popgreenL] (0,0) rectangle (4,1.3);
\fill[popblueL] (0,1.3) rectangle (4,2.3);
\fill[poporange!35] (0,2.3) rectangle (4,3.7);
\fill[white] (0,3.7) rectangle (4,4.5);
\fill[popgold!40] (0,4.5) rectangle (4,5.9);
\fill[poppink!35] (0,5.9) rectangle (4,7.3);
% outer frame
\draw[thick] (0,0) rectangle (4,7.3);
% internal boundaries
\draw (0,1.3) -- (4,1.3);
\draw (0,2.3) -- (4,2.3);
\draw (0,3.7) -- (4,3.7);
\draw[dashed] (0,4.5) -- (4,4.5);
\draw (0,5.9) -- (4,5.9);
% wavy unconformity lines at the great hiatus
\draw[thick,decorate,decoration={snake,amplitude=1.2pt}] (0,3.7) -- (4,3.7);
\draw[thick,decorate,decoration={snake,amplitude=1.2pt}] (0,4.5) -- (4,4.5);
% unit labels
\node[align=center,font=\small] at (2,0.65) {Ntem Complex \& greenstones\\ gneiss, granulite, charnockite};
\node[align=center,font=\small] at (2,1.8) {Nyong Group\\ quartzite, itabirite, marble};
\node[align=center,font=\small] at (2,3.0) {Pan-African belt\\ schists, gneisses, syn-tectonic granites};
\node[align=center,font=\small\itshape] at (2,4.1) {Palaeozoic hiatus\\ (550--145 Ma)};
\node[align=center,font=\small] at (2,5.2) {Cretaceous basins\\ Benue, Mamfe, Douala};
\node[align=center,font=\small] at (2,6.6) {Cenozoic cover \& CVL\\ coastal, Chad Basin, volcanics};
% age labels on right
\node[font=\footnotesize,anchor=west] at (4.15,0.65) {$>2.5$ Ga Archaean};
\node[font=\footnotesize,anchor=west] at (4.15,1.8) {$\sim 2.1$ Ga Palaeoproterozoic};
\node[font=\footnotesize,anchor=west] at (4.15,3.0) {$800$--$550$ Ma Neoproterozoic};
\node[font=\footnotesize,anchor=west] at (4.15,4.1) {Palaeozoic gap};
\node[font=\footnotesize,anchor=west] at (4.15,5.2) {$\sim 145$--$66$ Ma Mesozoic};
\node[font=\footnotesize,anchor=west] at (4.15,6.6) {$66$ Ma--today Cenozoic};
% arrow of time
\draw[->,thick,popink] (-0.5,0) -- (-0.5,7.3) node[midway,rotate=90,above,font=\small] {younger};
\end{tikzpicture}
\end{center}
\legende{Synthetic stratigraphic column of Cameroon, from the Archaean Ntem Complex at the base to the Cenozoic cover and Cameroon Volcanic Line rocks at the top. Wavy lines mark the great unconformities; the Palaeozoic interval is a hiatus with no preserved rocks.}
\end{popfigure}

\begin{popfigure}
\begin{center}
\begin{tikzpicture}[scale=1.0]
% basement
\fill[poporange!35] (0,0) -- (0,2.2) -- (1.2,2.6) -- (2.4,2.0) -- (3.6,2.5) -- (4.8,2.1) -- (6,2.4) -- (6,0) -- cycle;
\draw[thick] (0,0) -- (0,2.2) -- (1.2,2.6) -- (2.4,2.0) -- (3.6,2.5) -- (4.8,2.1) -- (6,2.4) -- (6,0);
% cover
\fill[popblueL] (0,2.2) -- (1.2,2.6) -- (2.4,2.0) -- (3.6,2.5) -- (4.8,2.1) -- (6,2.4) -- (6,3.6) -- (0,3.6) -- cycle;
\draw[thick] (0,3.6) -- (6,3.6);
% bedding in cover
\draw[popmuted] (0,2.9) -- (6,2.9);
\draw[popmuted] (0,3.25) -- (6,3.25);
% granite veins in basement
\draw[poppurple,thick] (1.5,0) -- (1.9,2.35);
\draw[poppurple,thick] (4.2,0) -- (4.0,2.25);
% labels
\node[align=center,font=\small] at (3,1.1) {Pan-African basement\\ schists, gneisses, granites ($\sim 600$ Ma)};
\node[align=center,font=\small] at (3,4.15) {Cretaceous cover: sandstones and shales ($\sim 120$ Ma)};
\draw[<-,thick,popdark] (4.9,2.18) -- (5.6,1.5) node[below,align=center,font=\small] {unconformity\\ hiatus $\sim 480$ Ma};
\end{tikzpicture}
\end{center}
\legende{Cross-section showing the angular/erosional unconformity between the folded Pan-African basement and the flat-lying Cretaceous sedimentary cover. The wavy erosion surface represents a hiatus of nearly half a billion years.}
\end{popfigure}

\courssub{Correlation with neighbouring countries}

The Cameroonian column does not stop at the borders; stratigraphy is international by nature.

\begin{itemize}
\item The \textbf{Ntem Complex} continues southward into Gabon, Equatorial Guinea and the Republic of Congo as part of the Congo craton; the same Archaean gneisses and the same Palaeoproterozoic (Francevillian-age) basins are found there.
\item The \textbf{Benue Trough} extends northeastward into Nigeria, where its Cretaceous marine limestones and shales are well dated by ammonites --- the same transgression that flooded the Garoua area.
\item The \textbf{Mamfe Basin} is the direct westward continuation of the Nigerian Benue rift system; its Cretaceous beds correlate with those of the lower Benue in Nigeria.
\item The \textbf{Chad Basin} sediments of the Far North are continuous with those of Chad and Nigeria, recording the same history of Lake Chad's expansion and shrinkage.
\item The \textbf{Cameroon Volcanic Line} continues offshore into the Gulf of Guinea islands (Bioko, Príncipe, São Tomé), showing that the same magmatic lineament crosses national boundaries.
\end{itemize}

\begin{exoresolu}[Placing a rock unit in the column]
\textbf{Statement.} During a field trip near Garoua (North Region), a student collects a fossiliferous marine limestone containing ammonites. (a) To which era does this rock most likely belong? (b) On what does it rest, and what is the nature of that contact? (c) Name one basin in a neighbouring country where the same unit can be correlated.

\tcblower
\textbf{Solution.}
\begin{enumerate}
\item (a) Ammonites are characteristic Mesozoic fossils, and marine limestones with ammonites in the Garoua region belong to the \textbf{Cretaceous} deposits of the Benue Trough. Era: \textbf{Mesozoic}.
\item (b) The Cretaceous limestone rests on the \textbf{Precambrian basement} (Pan-African gneisses and granites, or locally older units). The contact is an \textbf{unconformity}: an erosion surface representing a hiatus of several hundred million years (the Palaeozoic gap).
\item (c) The same Cretaceous marine beds correlate with the \textbf{Benue Trough of Nigeria}, where identical ammonite-bearing limestones and shales are found.
\end{enumerate}
\end{exoresolu}

\begin{exercice}[Building and reading the column]
\begin{enumerate}
\item List the five main Precambrian-to-Recent divisions of the stratigraphic column of Cameroon, giving one rock unit and one approximate age for each.
\item Explain why the contact between the Pan-African basement and the Cretaceous sediments of the Benue Trough is described as an unconformity.
\item A classmate claims: ``Cameroon has no Palaeozoic rocks, so Cameroon did not exist during the Palaeozoic.'' Correct this statement in two or three sentences.
\item Give two criteria you would use to correlate the Cretaceous sediments of the Mamfe Basin with those of the Nigerian Benue Trough.
\end{enumerate}
\end{exercice}

\begin{corrige}[Exercise: Building and reading the column]
\begin{enumerate}
\item Archaean: Ntem Complex gneisses, $>2.5$ Ga. Palaeoproterozoic: Nyong Group quartzites and itabirites, $\sim 2.1$ Ga. Neoproterozoic: Pan-African schists and syn-tectonic granites, $\sim 800$--$550$ Ma. Mesozoic: Cretaceous sediments of the Benue, Mamfe and Douala basins, $\sim 145$--$66$ Ma. Cenozoic: coastal deltaic beds, Chad Basin fluvio-lacustrine beds and CVL volcanics, $66$ Ma to present.
\item The basement rocks were deformed, metamorphosed and eroded long before the Cretaceous seas arrived. The Cretaceous beds lie horizontally on an erosion surface cut across the old basement, so the contact marks a long hiatus of erosion and non-deposition --- an unconformity.
\item Cameroon existed during the Palaeozoic, but no rocks of that age are preserved there. Either the region was an eroding landmass where no sediments accumulated, or Palaeozoic deposits were later eroded away. Absence of rocks means absence of record, not absence of events.
\item Any two of: identical or comparable fossil assemblages (Cretaceous ammonites, bivalves, pollen); similar lithological successions (continental sandstones passing to marine beds); radiometric ages from associated volcanic or intrusive rocks; the same structural relation to the rift faults and the basal unconformity.
\end{enumerate}
\end{corrige}

\begin{attention}[Exam tip: placing units in the correct era]
In map and data questions, examiners frequently ask you to assign a Cameroonian rock unit to its era. Memorise these anchors: \textbf{Archaean} = Ntem Complex and greenstone belts (gold, iron); \textbf{Palaeoproterozoic} = Nyong Group; \textbf{Neoproterozoic} = Pan-African schists, gneisses and granites (most of the country); \textbf{Mesozoic} = Cretaceous of Benue, Mamfe, Douala; \textbf{Cenozoic} = coastal basins, Chad Basin, all CVL volcanics. And remember: there is \emph{no} Palaeozoic column in Cameroon --- if a question seems to demand one, the correct answer is the hiatus.
\end{attention}

\begin{aretenir}[Key points of Section 7]
\begin{itemize}
\item A stratigraphic column arranges rock units from oldest (bottom) to youngest (top); unconformities mark gaps caused by erosion or non-deposition.
\item Cameroon's column: Archaean Ntem Complex and greenstones ($>2.5$ Ga) $\rightarrow$ Palaeoproterozoic Nyong Group ($\sim 2.1$ Ga) $\rightarrow$ Neoproterozoic Pan-African belt ($800$--$550$ Ma) $\rightarrow$ Palaeozoic hiatus $\rightarrow$ Cretaceous rift and marine basins $\rightarrow$ Cenozoic coastal, Chad Basin and CVL rocks.
\item Correlation with Gabon/Congo (craton), Nigeria (Benue Trough, Mamfe) and Chad (Chad Basin) shows that geological units cross political borders.
\item Missing rocks mean a missing record, not a missing history.
\end{itemize}
\end{aretenir}

With the stratigraphic framework now assembled, we hold the complete timetable of Cameroon's geological past. The final section of this chapter turns to what happens when those same geological forces --- volcanoes, faults, lakes and unstable slopes --- act on human timescales: the geologic hazards of Cameroon.

\coursec{Section 8: Geologic Hazards in Cameroon and Integration (Lessons 146-147)}

In the previous section we established the stratigraphic framework that records Cameroon's geological history from the Archaean Congo Craton through the Cretaceous basins to the recent Cameroon Volcanic Line (CVL). That history is not merely academic: the same tectono-magmatic processes that built our relief and concentrated our mineral wealth also generate the natural hazards that threaten communities today. This final section synthesises the entire course by examining how geological processes translate into hazards, how those hazards are distributed across the national territory, and how an integrated geological understanding supports sustainable development and disaster risk reduction.

\courssub{Fundamental Concepts: Hazard, Vulnerability, Risk and Exposure}

Before analysing specific Cameroonian hazards we must distinguish four terms that are often confused in everyday language but have precise meanings in disaster risk reduction.

\begin{definition}[Hazard, Vulnerability, Exposure and Risk]
A \cle{geological hazard} is a potentially damaging natural process or phenomenon of geological origin (e.g. volcanic eruption, earthquake, landslide, limnic gas burst). \cle{Vulnerability} is the susceptibility of people, infrastructure, economic activities or ecosystems to suffer harm when exposed to a hazard of a given intensity. \cle{Exposure} is the inventory of elements (population, buildings, crops, lifelines) located in the hazard zone. \cle{Risk} is the expected loss (fatalities, injuries, economic damage, disruption) resulting from the interaction of hazard, vulnerability and exposure over a specified time period. In its simplest quantitative form:
\[
\text{Risk} = \text{Hazard} \times \text{Vulnerability} \times \text{Exposure}
\]
where each factor is expressed in compatible units (e.g. annual probability of exceedance for hazard, fragility curves for vulnerability, replacement value for exposure).
\end{definition}

\begin{popfigure}
\begin{tikzpicture}[scale=0.9, every node/.style={font=\small}]
% Risk equation flowchart
\node[draw, rounded corners, fill=popblueL, thick, align=center] (hazard) at (0,3){Hazard\\(Probability \& Intensity)};
\node[draw, rounded corners, fill=popgreenL, thick, align=center] (vuln) at (0,0){Vulnerability\\(Fragility of elements)};
\node[draw, rounded corners, fill=poporange, thick, align=center] (expo) at (0,-3){Exposure\\(People, assets, GDP)};
\node[draw, rounded corners, fill=popdark, thick, text=white, align=center] (risk) at (5,0){RISK\\(Expected annual loss)};
\draw[->, thick, popblue] (hazard.east) -- node[above, pos=0.5] {$\times$} (risk.west);
\draw[->, thick, popgreen] (vuln.east) -- node[left, pos=0.5] {$\times$} (risk.west);
\draw[->, thick, poporange] (expo.east) -- node[below, pos=0.5] {$\times$} (risk.west);
% Example annotations
\node[align=center, text width=3.5cm, popmuted] at (-4,3) {Mt Cameroon\\eruption probability\\per year};
\node[align=center, text width=3.5cm, popmuted] at (-4,0) {Houses on\\steep lateritic\\slopes};
\node[align=center, text width=3.5cm, popmuted] at (-4,-3) {Population of\\Buea $\approx$ 150\,000};
\node[align=center, text width=4cm, popdark] at (9,0) {Annualised\\fatalities,\\economic loss\\in CFA francs};
\end{tikzpicture}
\legende{The risk equation: Risk = Hazard $\times$ Vulnerability $\times$ Exposure. Reducing any factor lowers risk.}
\end{popfigure}

\courssub{Spatial Distribution of Hazards in Cameroon}

Cameroon's hazards are not randomly scattered; they cluster along well-defined geological structures. The Cameroon Volcanic Line (CVL) is the dominant hazard corridor, concentrating volcanic, seismic, limnic and landslide risks. The coastal sedimentary basins add flooding and coastal erosion, while the Sanaga Fault introduces a distinct seismic source.

\begin{propriete}[Hazard Clustering along the CVL]
The majority of Cameroon's high-impact geological hazards occur within a $\sim$50 km wide corridor centred on the Cameroon Volcanic Line (Mount Cameroon $\to$ Manengouba $\to$ Bamboutos $\to$ Oku $\to$ Ngaoundéré). This clustering reflects a common deep-seated cause: the CVL is a lithospheric-scale weakness zone where mantle upwelling, crustal fracturing and magmatic intrusion simultaneously drive volcanism, seismicity, CO$_2$ accumulation in crater lakes, and steep, unstable topography prone to landslides.
\end{propriete}

\begin{popfigure}
\begin{tikzpicture}[scale=0.75]
% Simplified hazard map of Cameroon
% Coastline
\draw[thick, popblue] (0,0) -- (0,10) -- (1,11) -- (2,10) -- (3,10.5) -- (4,10) -- (5,9.5) -- (6,9) -- (7,8.5) -- (8,8) -- (9,7.5) -- (10,7) -- (11,6.5) -- (12,6) -- (13,5.5) -- (14,5) -- (15,4.5) -- (16,4) -- (17,3.5) -- (18,3) -- (19,2.5) -- (20,2) -- (21,1.5) -- (22,1) -- (23,0.5) -- (24,0);
% CVL axis
\draw[dashed, thick, poporange] (2,1) -- (4,3) -- (6,5) -- (8,7) -- (10,8.5) -- (12,9) -- (14,9.5) -- (16,9) -- (18,8) -- (20,6.5) -- (22,4.5);
% Volcanic hazard zones (orange circles)
\foreach \x/\y/\r/\name in {3/2/0.8/Mt Cameroon, 7/6/0.6/Manengouba, 10/8.5/0.6/Bamboutos, 13/9.5/0.6/Oku, 17/8/0.5/Ngaoundere} {
    \fill[poporange, opacity=0.3] (\x,\y) circle (\r);
    \draw[poporange, thick] (\x,\y) circle (\r);
    \node[poporange, font=\tiny] at (\x,\y+1.1) {\name};
}
% Limnic hazards (blue triangles)
\foreach \x/\y/\name in {10/8.2/Lake Nyos, 9/7.8/Lake Monoun} {
    \fill[popblue, opacity=0.3] (\x,\y) -- ++(-0.3,-0.5) -- ++(0.6,0) -- cycle;
    \draw[popblue, thick] (\x,\y) -- ++(-0.3,-0.5) -- ++(0.6,0) -- cycle;
    \node[popblue, font=\tiny] at (\x,\y-0.8) {\name};
}
% Seismic zones (orange squares)
\foreach \x/\y in {5/4, 12/8, 18/5} {
    \fill[poporange, opacity=0.3] (\x-0.4,\y-0.4) rectangle (\x+0.4,\y+0.4);
    \draw[poporange, thick] (\x-0.4,\y-0.4) rectangle (\x+0.4,\y+0.4);
}
% Landslide prone areas (green hatching)
\foreach \x/\y in {8/7, 11/8, 15/7.5} {
    \fill[popgreen, opacity=0.2, pattern=north east lines] (\x-0.6,\y-0.4) rectangle (\x+0.6,\y+0.4);
    \draw[popgreen, thick] (\x-0.6,\y-0.4) rectangle (\x+0.6,\y+0.4);
}
% Coastal flooding/erosion (purple band)
\fill[poppurple, opacity=0.2] (0,0) rectangle (24,1.2);
\draw[poppurple, thick] (0,1.2) -- (24,1.2);
\node[poppurple, font=\tiny, rotate=90] at (-0.5,0.6) {Coastal flooding \& erosion};
% Cities
\foreach \x/\y/\name in {3/1.5/Buea, 5/2.5/Douala, 10/7.5/Fundong, 14/9.5/Nkambe, 18/5.5/Meiganga} {
    \fill[black] (\x,\y) circle (0.1);
    \node[font=\tiny, anchor=west] at (\x+0.2,\y) {\name};
}
% Legend (manual)
\begin{scope}[xshift=25cm, yshift=2cm]
\fill[poporange, opacity=0.3] (0,0) circle (0.2); \draw[poporange, thick] (0,0) circle (0.2);
\node[anchor=west, font=\tiny] at (0.4,0) {Volcanic};
\pgftransformyshift{-0.6cm};
\fill[popblue, opacity=0.3] (0,0) -- ++(-0.2,-0.3) -- ++(0.4,0) -- cycle; \draw[popblue, thick] (0,0) -- ++(-0.2,-0.3) -- ++(0.4,0) -- cycle;
\node[anchor=west, font=\tiny] at (0.4,0) {Limnic};
\pgftransformyshift{-0.6cm};
\fill[poporange, opacity=0.3] (-0.2,-0.2) rectangle (0.2,0.2); \draw[poporange, thick] (-0.2,-0.2) rectangle (0.2,0.2);
\node[anchor=west, font=\tiny] at (0.4,0) {Seismic};
\pgftransformyshift{-0.6cm};
\fill[popgreen, opacity=0.2, pattern=north east lines] (-0.2,-0.2) rectangle (0.2,0.2); \draw[popgreen, thick] (-0.2,-0.2) rectangle (0.2,0.2);
\node[anchor=west, font=\tiny] at (0.4,0) {Landslide};
\pgftransformyshift{-0.6cm};
\fill[poppurple, opacity=0.2] (-0.2,-0.2) rectangle (0.2,0.2); \draw[poppurple, thick] (-0.2,-0.2) rectangle (0.2,0.2);
\node[anchor=west, font=\tiny] at (0.4,0) {Coastal};
\end{scope}
\end{tikzpicture}
\legende{Simplified multi-hazard zonation map of Cameroon. Orange circles: volcanic centres; blue triangles: limnic hazard lakes; orange squares: seismic zones; green hatched rectangles: landslide-prone lateritic highlands; purple band: coastal flooding and erosion. Dashed orange line: CVL axis.}
\end{popfigure}

\courssub{Volcanic Hazards: Mount Cameroon and the CVL}

Mount Cameroon (4\,095 m) is West Africa's most active volcano, with $\sim$20 documented eruptions since 1787. Its activity is dominantly effusive (basaltic lava flows) but explosive phreatomagmatic and Strombolian episodes occur, producing ash falls and ballistic projectiles.

\begin{exemplebox}[1999 and 2000 Eruptions of Mount Cameroon]
The 1999 eruption opened a 2 km fissure on the south flank at $\sim$2\,600 m a.s.l., emitting $\sim$25 $\times$ 10$^6$ m$^3$ of lava that reached the Atlantic Ocean, cutting the Limbe--Idenau road. The 2000 eruption occurred on the north-east flank near 1\,800 m a.s.l., threatening Buea's water catchment and the CDC tea plantations. Both eruptions were preceded by detectable seismic swarms (ARGV network) allowing partial evacuation.
\end{exemplebox}

\begin{attention}[Common Mistake: Confusing Lava Flow Speed with Pyroclastic Flow Speed]
Basaltic lava flows on Mount Cameroon typically advance at 0.1--10 m/h on gentle slopes, allowing evacuation. Pyroclastic density currents (rare but possible during phreatomagmatic explosions) travel at $>$100 km/h and are unsurvivable. In examinations, always specify which hazard you are discussing and the corresponding mitigation (evacuation vs. exclusion zone).
\end{attention}

\courssub{Limnic Hazards: Lake Nyos and Lake Monoun}

The 1986 Lake Nyos disaster (1\,746 fatalities) and the 1984 Lake Monoun event (37 fatalities) are the only known lethal limnic eruptions worldwide. Both lakes occupy maar craters on the CVL (Oku Volcanic Group). Deep magmatic CO$_2$ dissolves in the bottom waters under hydrostatic pressure; stratification prevents mixing. A trigger (landslide, heavy rainfall, minor earthquake) can initiate overturn, releasing a dense CO$_2$ cloud that hugs the ground and asphyxiates life down-valley.

\begin{methode}[Degassing Pipe Principle (Kling et al., 1990s -- operational since 2001)]
\begin{enumerate}
\item Install a vertical HDPE pipe from a raft to the lake bottom (Nyos: 203 m; Monoun: 93 m).
\item Prime the pipe by pumping water upward; dissolved CO$_2$ exsolves, creating buoyancy that sustains a self-driven fountain.
\item The fountain degasses deep water at the surface, progressively reducing total dissolved gas content.
\item Monitor dissolved gas pressure (total gas pressure / hydrostatic pressure ratio); maintain ratio $<$ 0.6 for safety.
\end{enumerate}
\end{methode}

\begin{savaistu}[Cameroon's Unique Contribution to Volcanology]
The Nyos/Monoun degassing programme is the world's first successful mitigation of a limnic eruption hazard. It demonstrates how fundamental geochemistry (Henry's law, gas solubility vs. pressure) translates directly into life-saving engineering. The ARGV (Argentinean--Cameroonian--French--Japanese collaboration) continues to monitor Lake Kivu (DRC/Rwanda) using Cameroonian expertise.
\end{savaistu}

\courssub{Seismic Hazards: CVL and Sanaga Fault}

Cameroon experiences moderate seismicity (M$_w$ typically $<$ 5.5). Two source zones dominate:
\begin{itemize}
\item \textbf{CVL-related seismicity:} Shallow ($<$ 15 km), clustered beneath volcanic centres (Mt Cameroon, Manengouba, Oku). Often swarms preceding eruptions.
\item \textbf{Sanaga Fault zone:} NE--SW trending structure crossing the Adamawa Plateau. The 2005 Mbam earthquake (M$_w$ 5.1, depth $\sim$10 km) caused damage in Bafia and Yaoundé (intensity VI--VII EMS-98). Historical events (1909, 1960s) suggest recurrence intervals of $\sim$50--100 years for M $>$ 5.
\end{itemize}

\begin{exoresolu}[Estimating Peak Ground Acceleration (PGA) for Yaoundé from the 2005 Mbam Earthquake]
\textbf{Statement:} The 2005 Mbam earthquake (M$_w$ 5.1, hypocentral depth 10 km, epicentre $\sim$60 km NE of Yaoundé) was felt at intensity VI in Yaoundé. Using the empirical attenuation relation for stable continental crust $\log_{10}(\text{PGA}) = -1.0 + 0.3M_w - \log_{10}(R) - 0.002R$ (PGA in cm/s$^2$, $R$ in km), estimate the PGA in Yaoundé and comment on the intensity observation.
\tcblower
\textbf{Solution:}
$M_w = 5.1$, hypocentral distance $R = \sqrt{60^2 + 10^2} \approx 60.8$ km.
\[
\log_{10}(\text{PGA}) = -1.0 + 0.3(5.1) - \log_{10}(60.8) - 0.002(60.8)
\]
\[
\log_{10}(\text{PGA}) = -1.0 + 1.53 - 1.78 - 0.12 = -1.37
\]
\[
\text{PGA} = 10^{-1.37} \approx 0.043\ \text{cm/s}^2 = 0.43\ \text{mm/s}^2 \approx 0.00004\ g
\]
This very low PGA is consistent with intensity VI (felt by all, slight damage) only if local site amplification (thick lateritic cover, sedimentary basin edge effects) increased ground motion by a factor of 10--100. \emph{Examiner's note: Always mention site effects when macroseismic intensity exceeds attenuation predictions.}
\end{exoresolu}

\courssub{Landslides and Mass Movements: The Lateritic Slope Problem}

Cameroon's humid tropical climate produces thick lateritic regolith (10--30 m) on steep volcanic and basement slopes. High-intensity rainfall ($>$200 mm/24 h) saturates the regolith, reducing suction and shear strength, triggering shallow translational slides and debris flows. The 2019 Bafoussam landslide (42 fatalities) and recurrent Bamboutos caldera rim failures exemplify this.

\begin{popfigure}
\begin{tikzpicture}[scale=0.9]
% Landslide cross-section on lateritic slope
% Topography
\draw[thick, popdark] (0,0) -- (1,0.5) -- (3,2.5) -- (5,3.5) -- (7,4) -- (9,4.2) -- (11,4) -- (13,3.5) -- (15,2.5) -- (17,1.5) -- (18,1) -- (20,0.5) -- (21,0);
% Laterite layer
\draw[fill=poporange!30, pattern=north east lines, pattern color=poporange] (0,0) -- (1,0.5) -- (3,2.5) -- (5,3.5) -- (7,4) -- (9,4.2) -- (11,4) -- (13,3.5) -- (15,2.5) -- (17,1.5) -- (18,1) -- (20,0.5) -- (21,0) -- (21,-1.5) -- (0,-1.5) -- cycle;
% Bedrock
\draw[fill=popdark!30] (0,-1.5) -- (21,-1.5) -- (21,-3) -- (0,-3) -- cycle;
% Slip surface
\draw[thick, dashed, poporange] (4,3.2) .. controls (7,2.5) and (10,1.8) .. (14,2.2);
% Arrows for movement
\draw[->, thick, poporange] (8,3.5) -- (9,3.2);
\draw[->, thick, poporange] (10,3) -- (11,2.7);
% Rainfall
\foreach \x in {2,5,8,11,14,17} {
    \draw[popblue, thick] (\x,5) -- (\x-0.2,4.5);
    \draw[popblue, thick] (\x+0.3,5) -- (\x+0.1,4.5);
}
% Water table
\draw[popblue, thick, dashed] (3,2.8) .. controls (7,2.2) and (11,1.5) .. (15,2.0);
% Labels
\node[poporange, font=\small] at (10,2.5) {Lateritic regolith (10--30 m)};
\node[popdark, font=\small] at (10,-2.2) {Fractured bedrock (basalt/granite)};
\node[poporange, font=\small] at (9,1.5) {Slip surface};
\node[popblue, font=\small] at (10,4.8) {Intense rainfall};
\node[popblue, font=\small] at (18,2.5) {Perched water table};
% House
\draw[fill=popdark] (8.5,4.2) rectangle (9.5,4.8);
\draw[fill=poporange] (8.7,4.2) -- (9,4.5) -- (9.3,4.2) -- cycle;
\node[font=\tiny] at (9,4.0) {House};
\end{tikzpicture}
\legende{Translational landslide in lateritic regolith on a steep volcanic slope. Intense rainfall raises the perched water table, reducing effective stress along a planar slip surface at the laterite/bedrock interface. The dashed orange line marks the failure plane.}
\end{popfigure}

\begin{attention}[Examination Trap: Landslide Trigger in Cameroon]
In GCE questions, \textbf{never} state that Cameroonian landslides are primarily triggered by earthquakes. The vast majority are rainfall-induced on weathered lateritic slopes. Earthquakes (M $<$ 5.5) are too weak and infrequent to be the main trigger. Cite the 2019 Bafoussam event (trigger: 3 days of $>$300 mm rain) and the Bamboutos recurrent slides (trigger: seasonal saturation).
\end{attention}

\courssub{Flooding and Coastal Erosion}

Douala (Wouri estuary), Limbe (Atlantic coast) and Kribi (Nyong estuary) face compound flooding: fluvial (heavy rainfall, poor drainage), pluvial (urbanisation, impermeable surfaces) and coastal (storm surges, spring tides, sea-level rise $\sim$3 mm/yr). Mangrove loss and sand mining exacerbate coastal erosion (retreat rates 1--3 m/yr locally). Delta subsidence (Holocene sediment compaction) adds $\sim$1--2 mm/yr relative sea-level rise.

\courssub{Hazard Mitigation and Disaster Risk Reduction in Cameroon}

Cameroon's mitigation framework operates at three levels:

\begin{center}
\begin{tabularx}{\linewidth}{|l|X|X|}\hline
\rowcolor{popblueL}
\textbf{Level} & \textbf{Institutional Actor} & \textbf{Key Actions} \\ \hline
Monitoring & ARGV (Mt Cameroon), IRGM (seismic), MINRESI/ONUD (Nyos/Monoun) & Real-time seismic, geodetic, gas, lake-level networks; bulletins. \\ \hline
Land-use Planning & MINDUH, Local Councils & Hazard zonation maps (PNAT), building codes (Eurocode 8 adaptation), relocation of high-risk settlements (e.g. Nyos downstream villages). \\ \hline
Preparedness \& Response & DPC (Civil Protection), Red Cross, NGOs & Early warning sirens (Nyos), simulation exercises, community risk committees, contingency stocks. \\ \hline
\end{tabularx}
\end{center}

\begin{methode}[Drawing a Simple Hazard Zonation Map (Examinable Skill)]
\begin{enumerate}
\item \textbf{Collect data:} Geological map (lithology, structures), DEM (slope, aspect), historical event inventory (catalogues, newspapers, oral history), land-use map.
\item \textbf{Define hazard sources:} Map volcanic vents, fault traces, lake shores, steep slopes ($>$25°), floodplains.
\item \textbf{Assign intensity zones:} For each hazard, delineate high/medium/low zones based on distance from source, slope, lithology, recurrence interval (e.g. lava flow probability contours from vent).
\item \textbf{Combine into multi-hazard map:} Overlay single-hazard maps; classify each pixel by the highest hazard level present (conservative approach) or by a weighted sum.
\item \textbf{Validate:} Cross-check with known past events; adjust boundaries.
\item \textbf{Produce output:} Legend, scale, north arrow, coordinate grid, metadata (date, author, data sources).
\end{enumerate}
\end{methode}

\courssub{Integration 6: Geology, Resources and Hazards for Sustainable Development}

This course has traversed Cameroon's geology from the Archaean basement to the active CVL. The integration challenge is to manage the \emph{same} geological endowment that provides resources (oil, gas, bauxite, cobalt, nickel, geothermal, groundwater, fertile volcanic soils, tourism landscapes) while mitigating the hazards it generates.

\begin{aretenir}[Synthesis: The Geological Duality of Cameroon]
\begin{itemize}
\item \textbf{Resources:} CVL $\to$ fertile soils (coffee, tea, oil palm), geothermal potential, mineral deposits (Ni-Co laterites, bauxite), tourism (crater lakes, volcanoes). Sedimentary basins $\to$ hydrocarbons, groundwater, construction aggregates, coastal fisheries.
\item \textbf{Hazards:} CVL $\to$ eruptions, earthquakes, CO$_2$ bursts, landslides. Basins $\to$ flooding, subsidence, coastal erosion.
\item \textbf{Sustainable Management:} Requires \textbf{integrated land-use planning} informed by multi-hazard maps, \textbf{enforced building codes} (seismic, volcanic), \textbf{community-based monitoring} (rainfall thresholds for landslides, lake gas sensors), and \textbf{economic diversification} to reduce vulnerability of exposed populations.
\end{itemize}
\end{aretenir}

\begin{exemplebox}[Essay Question Model: "Discuss how geological knowledge contributes to sustainable development in Cameroon."]
\textbf{Structure your answer around three pillars:}
\begin{enumerate}
\item \textbf{Resource identification \& management:} Stratigraphy $\to$ basin analysis $\to$ hydrocarbon exploration; laterite profiles $\to$ Ni-Co/bauxite evaluation; CVL structure $\to$ geothermal targeting.
\item \textbf{Hazard mitigation:} Hazard zonation (this section) $\to$ land-use zoning, early warning, engineering design (seismic codes, lava diversion channels).
\item \textbf{Environmental stewardship:} Hydrogeology $\to$ aquifer protection; limnology $\to$ lake ecosystem management; coastal geology $\to$ mangrove restoration, setback lines.
\end{enumerate}
\textbf{Conclude:} Geology is not a luxury; it is the evidence base for the National Development Strategy (SND30). Every CFA franc invested in geological mapping returns multiples in avoided losses and optimised resource revenue.
\end{exemplebox}

\begin{exercice}[GCE-Style Practice Questions]
\begin{itemize}
\item Define hazard, vulnerability, exposure and risk. Using the 1986 Lake Nyos event, illustrate each term with a specific example from that disaster.
\item Draw a labelled cross-section showing the mechanism of a rainfall-triggered translational landslide in lateritic regolith on a volcanic slope. Indicate the slip surface, water table, and direction of movement.
\item The 2005 Mbam earthquake (M$_w$ 5.1) caused intensity VI damage in Yaoundé, 60 km away. Explain why the observed intensity exceeded the prediction from a standard attenuation law. Propose two mitigation measures for Yaoundé's building stock.
\item Sketch a multi-hazard zonation map of the Mount Cameroon region (radius 30 km) showing volcanic, seismic, landslide and coastal hazards. Justify the spatial extent of each zone.
\item "Cameroon's geological hazards are concentrated along the Cameroon Volcanic Line." Discuss this statement with reference to volcanic, limnic, seismic and landslide hazards.
\end{itemize}
\end{exercice}

\begin{corrige}[Exercise 1]
\textbf{Hazard:} Sudden overturn of stratified Lake Nyos releasing $\sim$1.6 $\times$ 10$^9$ m$^3$ of CO$_2$. \textbf{Vulnerability:} Villagers sleeping in low-lying valleys, no warning system, housing not gas-tight. \textbf{Exposure:} $\sim$2\,000 people in the gas path (Nyos, Cha, Subum, Fang). \textbf{Risk realised:} 1\,746 fatalities, 3\,000+ livestock lost, psychological trauma, economic disruption. Post-1986 risk reduction: degassing pipes (hazard reduction), relocation uphill (exposure reduction), sirens and education (vulnerability reduction).
\end{corrige}

\begin{corrige}[Exercise 3]
\textbf{Reasons for higher intensity:} (1) Site amplification: Yaoundé sits on thick lateritic cover and weathered granite overlying a sedimentary basin edge; soft soils amplify seismic waves (resonance). (2) Basin-edge effect: 2D/3D wave focusing at the Congo Craton / sedimentary basin boundary. (3) Building vulnerability: Pre-1990s non-engineered masonry, poor reinforcement. \textbf{Mitigation measures:} (1) Enforce Eurocode 8 / NBC 2019 seismic design for new construction and retrofitting of critical facilities (hospitals, schools). (2) Microzonation study for Yaoundé to identify high-amplification zones and restrict high-occupancy buildings there.
\end{corrige}

\begin{attention}[Final Examination Advice for This Chapter]
\begin{itemize}
\item Always locate hazards on a sketch map of Cameroon (CVL, basins, faults, lakes, cities).
\item Quantify where possible: recurrence intervals, volumes, fatalities, PGA values, rainfall thresholds.
\item Link mitigation to the \textbf{risk equation}: state explicitly whether a measure reduces hazard, vulnerability or exposure.
\item In essay questions, integrate across the syllabus: e.g. "Lateritic weathering (Section 2) creates both Ni-Co resources and landslide hazards (Section 8)."
\end{itemize}
\end{attention}

This concludes the Upper Sixth Geology course. You now possess a coherent framework — from mineral identification to national hazard integration — that meets the GCE Advanced Level syllabus and equips you for university geoscience or professional practice in Cameroon's resource and risk sectors. Revise actively: draw maps, sketch sections, calculate numbers, and teach the concepts to a peer. Success in the examination and beyond depends on \emph{understanding}, not memorisation. Bonne chance!

\newpage

\coursec{Integration activity}

\courssub{Aims of the activity}

This integration activity closes the chapter \emph{The Geology of Cameroon} by asking you to act as a professional geologist. By the end of the activity, you should be able to:

\begin{itemize}
\item \textbf{Mobilise} your knowledge of igneous geology (the Cameroon Volcanic Line, basaltic volcanism), geological structures (faults, fractures), stratigraphy, natural resources and geological hazards, and apply them \emph{simultaneously} to one real Cameroonian region;
\item \textbf{Read and interpret} a simplified geological map: identify rock types, structures, and the spatial relationship between geology, relief and human settlement;
\item \textbf{Zone} a territory according to volcanic, seismic and landslide risk levels, justifying each zone with geological evidence;
\item \textbf{Propose} a reasoned land-use plan (agriculture, quarrying, geothermal prospecting, settlement) that balances economic development against risk;
\item \textbf{Formulate} concrete mitigation and monitoring measures, and \textbf{defend} your plan orally, as geologists do before administrative authorities;
\item \textbf{Compare} your proposals with the real monitoring work carried out around Mount Cameroon by the \cle{ARGV} research programme based at the Ekona research centre.
\end{itemize}

\begin{attention}[Why this matters for the GCE]
Integration questions at the Advanced Level never test one isolated lesson. They present a situation (a map, a disaster, a project) and expect you to \emph{combine} several parts of the syllabus. This activity trains exactly that skill: moving from \emph{knowing} to \emph{using}. In the examination, marks are awarded for the \emph{link} between observation and decision, not for either one alone.
\end{attention}

\courssub{Situation}

\textbf{The problem.}\par
The South West Region of Cameroon is experiencing rapid population growth. Buea, on the south-eastern flank of Mount Cameroon (about \SI{4095}{m}, the highest peak in West and Central Africa), and Limbe, on the coast, are expanding fast. Land is becoming scarce, and pressure is mounting to open new farmland, new quarters and new quarries on the volcano's slopes. At the same time, Mount Cameroon is the most active volcano in West and Central Africa: it has erupted repeatedly in historical time (notably in 1922, 1959, 1982, 1999 and 2000), producing basaltic lava flows that descended the flanks — the 1999 and 2000 flows affected the upper slopes above Buea and towards the Limbe road. The region also records frequent felt earthquakes, and the steep, deeply weathered slopes around Buea suffered destructive landslides in 2001 and 2003, which killed people and destroyed homes.

The Governor of the South West Region has therefore commissioned your team of geologists to produce a \textbf{land-use and risk-management plan} for the Mount Cameroon--Buea--Limbe area. You are given the simplified geological map below and the data sheet that follows.

\begin{popfigure}
\begin{tikzpicture}[scale=1.0, font=\small]
% Sea
\fill[popblueL] (-0.5,-3.6) rectangle (10.5,-1.6);
\node[popblue] at (5.0,-2.6) {\textbf{Atlantic Ocean (Bight of Biafra)}};
% Land zones
\fill[popgreenL] (-0.5,3.4) rectangle (10.5,-1.6);
% Lava flow zones (young flows)
\fill[poporange!55] (4.2,3.4) -- (5.6,3.4) -- (6.4,-1.6) -- (3.4,-1.6) -- cycle;
% Older volcanic apron
\fill[popgold!40] (2.0,3.4) -- (4.2,3.4) -- (3.4,-1.6) -- (1.2,-1.6) -- cycle;
\fill[popgold!40] (5.6,3.4) -- (8.0,3.4) -- (8.8,-1.6) -- (6.4,-1.6) -- cycle;
% Summit
\fill[poppink!70] (4.55,3.4) circle (0.55);
\node[align=center] at (4.55,2.95) {\textbf{Summit}\\ about \SI{4095}{m}};
% Faults
\draw[popdark, very thick, dashed] (1.6,3.4) -- (2.6,-1.6);
\draw[popdark, very thick, dashed] (7.4,3.4) -- (8.4,-1.6);
\node[popdark, align=center] at (1.05,1.6) {Fault F1};
\node[popdark, align=center] at (8.95,1.6) {Fault F2};
% Towns
\fill[black] (3.0,0.4) circle (0.09); \node[anchor=south east] at (3.0,0.4) {\textbf{Buea}};
\fill[black] (7.6,-1.0) circle (0.09); \node[anchor=south west] at (7.6,-1.0) {\textbf{Limbe}};
\fill[black] (1.6,-0.6) circle (0.09); \node[anchor=south] at (1.6,-0.6) {Idenau};
% Road
\draw[thick] (3.0,0.4) -- (4.8,-0.2) -- (7.6,-1.0);
\node[rotate=-13] at (5.4,-0.45) {Buea--Limbe road};
% 1999/2000 flow path
\draw[poppurple, very thick, ->] (4.7,3.0) -- (4.4,1.6) -- (4.6,0.2) -- (4.3,-1.2);
\node[poppurple, align=center] at (6.0,1.9) {1999/2000\\ lava path};
% North arrow
\draw[->, thick] (10.0,2.4) -- (10.0,3.2); \node at (10.0,3.35) {N};
% Legend
\node[anchor=west, align=left] at (8.6,3.3) {\textcolor{poppink!70}{$\blacksquare$} Summit cone};
\node[anchor=west, align=left] at (8.6,2.85) {\textcolor{poporange!55}{$\blacksquare$} Young flows ($<$100 yr)};
\node[anchor=west, align=left] at (8.6,2.4) {\textcolor{popgold!40}{$\blacksquare$} Older volcanic apron};
\node[anchor=west, align=left] at (8.6,1.95) {\textcolor{popgreenL}{$\blacksquare$} Sedimentary plain};
\end{tikzpicture}
\legende{Simplified geological sketch map of the Mount Cameroon--Buea--Limbe area (schematic, not to scale).}
\end{popfigure}

\textbf{Data sheet provided to your team.}
\begin{enumerate}
\item \textbf{Rocks:} the volcano is built almost entirely of \cle{basaltic} lava flows, scoria cones and pyroclastic fall deposits (Strombolian to Hawaiian activity). The coastal plain around Limbe and Idenau consists of Cretaceous--Cenozoic sedimentary rocks (sandstones, shales, limestones) of the Douala basin, overlain by volcanic ash and young soils.
\item \textbf{Soils:} weathered basalt on the mid-slopes gives deep, fertile \cle{andosols} (rich volcanic soils) that support banana, oil-palm and market gardening; on steep upper slopes the same soils are thin and unstable.
\item \textbf{Structures:} Mount Cameroon lies on the ocean-continent segment of the \cle{Cameroon Volcanic Line}; two major fault families (F1, F2, trending roughly NE--SW) cut the flanks and channel both magma ascent and groundwater. Earthquake epicentres cluster along these faults.
\item \textbf{Hazard record:} historical eruptions about once every 15--20 years; lava flows typically follow the summit-to-coast corridors; felt earthquakes several times per year; destructive landslides on the steep, water-saturated slopes above Buea after heavy rains (2001, 2003); volcanic gas and ash fall affect villages downwind.
\item \textbf{Resources:} abundant basalt and scoria (aggregates, blocks, pozzolana for cement); fertile soils; hot springs and high heat flow along faults (geothermal potential); scenery and the summit (tourism); offshore and onshore sedimentary basin (petroleum in the Rio del Rey / Douala basin further south-east).
\end{enumerate}

\begin{methode}[How to attack an integration problem]
\begin{enumerate}
\item \textbf{Observe} the map and data sheet: list every rock unit, structure, town and hazard indicator without interpreting yet.
\item \textbf{Interpret} each observation geologically: what process produced it, and what does it imply (fertility, instability, permeability, heat)?
\item \textbf{Cross} the hazards with the human stakes: where do people, roads and farmland overlap with the dangerous zones?
\item \textbf{Decide}: propose land uses and mitigation measures, each justified by at least one geological argument from steps 1--2.
\item \textbf{Check} coherence: no settlement in a high-risk zone, no quarry on a fault, no contradiction between your zoning and your land-use plan.
\end{enumerate}
\end{methode}

\courssub{Tasks}

Work in groups of four to six. Answer the following tasks in order, writing a short justified report (2--3 pages) and preparing a 5-minute presentation.

\begin{enumerate}
\item \textbf{Reading the map.} Identify and list the rock types and geological structures shown on the map. For each, state the evidence on the map and the geological process that produced it. (Skills: igneous geology, structures.)
\item \textbf{Risk zoning.} Divide the mapped area into three zones (high, moderate, low risk) for each of the three hazards: lava flow, earthquake, landslide. Present your zoning as a table and justify each zone using the data sheet. Which town is most exposed, and why?
\item \textbf{Land-use plan.} Propose suitable locations for: (a) intensive agriculture; (b) a quarry for volcanic materials; (c) a geothermal prospecting borehole; (d) new residential quarters. For each choice, give one geological argument \emph{for} and one constraint to respect.
\item \textbf{Mitigation and monitoring.} Propose exactly \textbf{three} concrete measures (at least one structural/engineering and one monitoring measure) to reduce risk for the population of Buea and Limbe. For each measure, state what it protects, how it works, and who should operate it.
\item \textbf{Defence.} Present your plan in 5 minutes: one member presents the zoning, one the land-use plan, one the mitigation measures. Be ready to answer questions from the ``Governor's cabinet'' (the rest of the class).
\item \textbf{Comparison.} After the presentations, compare your proposals with the real monitoring programme run from the Ekona research centre (seismic and ground-deformation monitoring of Mount Cameroon, hazard mapping, public sensitisation). What did you anticipate well? What did you miss?
\end{enumerate}

\courssub{Model answer}

\textbf{Task 1 — Reading the map.}\par
The map shows three volcanic units and one sedimentary unit. The \emph{summit cone} (pink) is built of young scoria and recent basaltic lava from Strombolian--Hawaiian eruptions. The \emph{young flows} (orange, $<$100 years, including the 1999/2000 flows) are \cle{basalt} — a mafic lava with low silica content, hence low viscosity, which explains why the flows travel far down gentle slopes in narrow corridors. The \emph{older volcanic apron} (gold) consists of older, weathered basalt flows and scoria now covered by andosols. The \emph{sedimentary plain} (green) around Limbe and Idenau belongs to the Douala basin: sandstones, shales and limestones deposited in marine to coastal environments. Two \cle{faults} (F1, F2, dashed lines) cut the area with a NE--SW trend, parallel to the Cameroon Volcanic Line; they are extensional fractures that channel magma ascent (the volcano sits in their intersection zone) and localise earthquakes.

\textbf{Task 2 — Risk zoning.}

\begin{center}
\begin{tabularx}{\linewidth}{|l|X|X|X|}
\hline
\rowcolor{popblueL} \textbf{Hazard} & \textbf{High risk} & \textbf{Moderate risk} & \textbf{Low risk} \\
\hline
Lava flow & Summit cone and the two summit-to-coast corridors (young-flow zone, incl. 1999/2000 path) & Older volcanic apron flanking the corridors & Sedimentary plain away from flow axes \\
\hline
Earthquake & Strips along faults F1 and F2 (epicentres cluster there) & Whole volcanic edifice & Coastal plain far from faults \\
\hline
Landslide & Steep upper slopes above Buea (thin saturated soils, 2001/2003 events) & Mid-slopes with deep andosols & Flat coastal plain \\
\hline
\end{tabularx}
\end{center}

\textbf{Buea} is the most exposed town: it sits on the mid-flank, close to the young-flow corridor, near fault F1, and directly below steep landslide-prone slopes. Limbe is moderately exposed (lava can reach the coast, as in 1922, but the plain is flat and stable).

\textbf{Task 3 — Land-use plan.}
\begin{itemize}
\item \textbf{(a) Agriculture:} the mid-slope andosol belt and the ash-covered coastal plain. \emph{For:} deep, fertile volcanic soils, good rainfall. \emph{Constraint:} avoid the lava corridors and the steepest slopes (erosion and landslide risk); keep the upper slopes forested.
\item \textbf{(b) Quarry:} the older volcanic apron, away from the flow corridors and downslope of settlements. \emph{For:} massive basalt and scoria give excellent aggregates and pozzolana. \emph{Constraint:} not on fault traces (instability; blasting could reactivate fractures) and not directly above Buea (noise, dust, slope destabilisation).
\item \textbf{(c) Geothermal borehole:} near the intersection of faults F1/F2 on the lower flank, where hot springs already emerge. \emph{For:} faults provide fracture permeability and the volcano provides the heat source. \emph{Constraint:} seismic monitoring during drilling (risk of induced seismicity) and a safe distance from the lava corridor.
\item \textbf{(d) New settlements:} the sedimentary plain between Limbe and Idenau, off the flow axes and away from fault traces. \emph{For:} flat, stable ground with low lava and landslide risk. \emph{Constraint:} check flood and coastal risk, and keep evacuation routes to the Buea--Limbe road open.
\end{itemize}

\textbf{Task 4 — Three mitigation and monitoring measures.}
\begin{enumerate}
\item \textbf{Volcano monitoring network (monitoring):} a network of seismometers, GPS stations and gas sensors on the flanks, operated by the Ekona research centre, to detect magma ascent weeks to days before an eruption and trigger evacuation of the corridor zones. Protects Buea, Limbe and the road.
\item \textbf{Reforestation and drainage of the steep slopes above Buea (structural):} tree cover plus engineered drainage channels to reduce soil saturation and stabilise slopes against rain-triggered landslides; combined with a ban on construction on the steepest sectors. Protects Buea.
\item \textbf{Lava-flow diversion barriers and land-use regulation (structural + regulatory):} earthen barriers along the historic flow corridors to deflect future flows away from the road and villages, together with legal zoning forbidding new building inside the high-risk corridors and an evacuation plan with marked routes and sensitisation campaigns. Protects all settlements.
\end{enumerate}

\textbf{Task 6 — Comparison with the real programme.}\par
The real monitoring programme run from the Ekona research centre indeed operates seismic and ground-deformation monitoring of Mount Cameroon, produces hazard maps and runs public sensitisation — our measures 1 and 3 anticipated this well. What student plans often miss: the importance of \emph{gas} monitoring (a lesson of the Lake Nyos disaster of 1986, which occurred on the same volcanic line), the social dimension (reluctance of populations to abandon fertile land), and the need for continuous funding and data transmission. A good plan is geological \emph{and} human.

\begin{attention}[Common mistakes in integration answers]
\begin{itemize}
\item Proposing a land use (e.g. ``build houses near Limbe'') with \emph{no geological justification} — every decision must cite a rock, a structure, a soil or a hazard fact.
\item Contradicting yourself: zoning the flow corridor ``high risk'' and then placing a quarry or a quarter inside it.
\item Forgetting that faults are \emph{both} a hazard (earthquakes, instability) \emph{and} a resource (geothermal permeability, groundwater) — a good answer exploits this duality.
\item Listing mitigation measures without saying \emph{what they protect} and \emph{who operates them}.
\end{itemize}
\end{attention}

\begin{aretenir}[What the examiner expects]
A strong integration answer always follows the chain: \textbf{observation on the map} $\rightarrow$ \textbf{geological interpretation} $\rightarrow$ \textbf{consequence for land use or risk} $\rightarrow$ \textbf{justified decision}. Never propose a land use without a geological argument, and never describe geology without drawing the practical consequence.
\end{aretenir}

\newpage

\coursec{Methods and worked examples}

\courssub{Methods and Worked Examples}

\begin{methode}[Identifying the Main Igneous Provinces of Cameroon]
\textbf{Objective:} Distinguish the three major igneous provinces (Congo Craton basement, Cameroon Volcanic Line, and the Benue Trough-related magmatism) on a geological map and in hand specimen.
\begin{enumerate}
\item \textbf{Observe the regional setting:} Locate the study area on a 1:1\,000\,000 geological map of Cameroon. Note the position relative to the \cle{Congo Craton} (south‑east), the \cle{Cameroon Volcanic Line (CVL)} (SW–NE trend), and the \cle{Benue Trough} (north).
\item \textbf{Recognise rock types:} 
\begin{itemize}
\item Craton: high‑grade gneisses, granitoids, charnockites (Archean–Palaeoproterozoic).
\item CVL: alkaline basalts, basanites, phonolites, trachytes, and associated pyroclastics (Cenozoic to Recent).
\item Benue Trough: tholeiitic basalts, dolerites, gabbros (Cretaceous), often as dykes and sills.
\end{itemize}
\item \textbf{Use geochemical discriminants (if data available):} Plot \ce{Na2O+K2O} vs \ce{SiO2} (TAS diagram) or \ce{Zr/Nb} vs \ce{Nb/Y} to separate alkaline (CVL) from sub‑alkaline (Benue) suites.
\item \textbf{Check age relationships:} Radiometric ages (K–Ar, Ar–Ar, U–Pb) confirm: Craton > 2\,Ga; Benue magmatism ~130–90\,Ma; CVL activity from ~66\,Ma to present.
\item \textbf{Synthesize:} Assign each outcrop to a province and record the dominant lithology, age, and tectonic setting.
\end{enumerate}
\textbf{Key reflex:} Always correlate field observation (texture, mineralogy) with regional map position before invoking geochemistry.
\end{methode}

\begin{exoresolu}[Mapping Igneous Provinces on a Simplified Geological Map]
\textbf{Statement:} A simplified geological map of Cameroon (scale 1:5\,000\,000) shows three coloured zones: Zone A (pink) in the south‑east, Zone B (green) along a SW–NE line from Mount Cameroon to Lake Chad, and Zone C (yellow) in the north. Hand specimens from three localities are provided:
\begin{itemize}
\item Locality 1 (Zone A): coarse‑grained biotite granite with zircon U–Pb age 2\,150\,Ma.
\item Locality 2 (Zone B): vesicular basanite with olivine phenocrysts, K–Ar age 12\,Ma.
\item Locality 3 (Zone C): fine‑grained dolerite dyke cutting Cretaceous sandstones, Ar–Ar age 95\,Ma.
\end{itemize}
Using the method above, assign each locality to its igneous province, justify with two independent criteria, and sketch a labelled cross‑section (SW–NE) showing the three provinces.
\tcblower
\textbf{Detailed Solution:}
\begin{enumerate}
\item \textbf{Locality 1 (Zone A – pink)}: 
\begin{itemize}
\item \textit{Map position}: South‑east Cameroon, within the Congo Craton.
\item \textit{Lithology \& age}: Coarse‑grained biotite granite, U–Pb zircon age 2\,150\,Ma (Palaeoproterozoic). Typical of cratonic granitoids.
\item \textit{Province}: \cle{Congo Craton basement}.
\end{itemize}
\item \textbf{Locality 2 (Zone B – green)}: 
\begin{itemize}
\item \textit{Map position}: On the SW–NE trend of the Cameroon Volcanic Line (Mount Cameroon to Lake Chad).
\item \textit{Lithology \& age}: Vesicular basanite with olivine phenocrysts, K–Ar age 12\,Ma (Miocene). Alkaline composition, young age, typical of CVL volcanism.
\item \textit{Province}: \cle{Cameroon Volcanic Line (CVL)}.
\end{itemize}
\item \textbf{Locality 3 (Zone C – yellow)}: 
\begin{itemize}
\item \textit{Map position}: Northern Cameroon, within the Benue Trough.
\item \textit{Lithology \& age}: Fine‑grained dolerite dyke cutting Cretaceous sandstones, Ar–Ar age 95\,Ma (Cretaceous). Tholeiitic affinity, associated with rift‑related magmatism.
\item \textit{Province}: \cle{Benue Trough magmatism}.
\end{itemize}
\item \textbf{Cross‑section (SW–NE)}: 
\begin{popfigure}
\begin{tikzpicture}[scale=0.8]
% Craton
\fill[poppink!30] (0,0) rectangle (4,2);
\node[align=center] at (2,1) {Congo Craton\\Granitoids, Gneisses};
% CVL
\fill[popgreen!30] (4,0) rectangle (10,3);
\draw[popgreen, thick] (4,3) -- (5,4) -- (6,3.5) -- (7,4.2) -- (8,3.8) -- (9,4) -- (10,3);
\node[align=center] at (7,3.5) {CVL\\Alkaline Volcanoes};
% Benue Trough
\fill[poporange!30] (10,0) rectangle (14,2);
\node[align=center] at (12,1) {Benue Trough\\Cretaceous Sediments\\+ Dolerites};
% Basement line
\draw[thick] (0,0) -- (14,0);
% Labels
\node[below] at (2,-0.5) {SW};
\node[below] at (12,-0.5) {NE};
\end{tikzpicture}
\legende{Schematic SW–NE cross‑section across Cameroon showing the three igneous provinces.}
\end{popfigure}
\end{enumerate}
\textbf{Examiner's tip:} Always quote the \emph{map position}, \emph{lithology}, and \emph{geochronology} as three independent lines of evidence.
\end{exoresolu}

\begin{methode}[Constructing a Stratigraphic Column for a Sedimentary Basin]
\textbf{Objective:} Build an accurate stratigraphic column (lithostratigraphy) for the \cle{Logone‑Birni Basin} (or any Cameroonian basin) from borehole logs and outcrop data.
\begin{enumerate}
\item \textbf{Gather data:} Collect borehole logs (depth, lithology, gamma ray), outcrop descriptions, and any biostratigraphic ages (palynology, foraminifera).
\item \textbf{Establish a vertical scale:} Choose a scale (e.g., 1\,cm = 100\,m) and draw a vertical line representing total depth.
\item \textbf{Identify formation boundaries:} Use lithological changes (sandstone $\leftrightarrow$ shale, carbonate $\leftrightarrow$ clastic) and wireline log markers (gamma ray spikes, resistivity shifts).
\item \textbf{Assign formation names and ages:} Correlate with the standard Cameroonian stratigraphic chart (e.g., \cle{Bima Sandstone} – Early Cretaceous; \cle{Yola Formation} – Albian; \cle{Pindiga Formation} – Cenomanian–Turonian; \cle{Gombe Sandstone} – Santonian–Campanian; \cle{Kerri‑Kerri Formation} – Paleocene).
\item \textbf{Add key surfaces:} Mark unconformities (angular, disconformity), maximum flooding surfaces (MFS), and sequence boundaries.
\item \textbf{Annotate:} Include thickness, dominant lithology, sedimentary structures, fossil content, and depositional environment (fluvial, deltaic, shallow marine).
\item \textbf{Review:} Check that the column respects the law of superposition and matches regional correlation panels.
\end{enumerate}
\textbf{Key formula:} Thickness of a unit = (bottom depth – top depth) corrected for dip if the borehole is deviated.
\end{methode}

\begin{exoresolu}[Drawing a Stratigraphic Column from Borehole Data]
\textbf{Statement:} A borehole (BH‑01) in the Logone‑Birni Basin penetrated the following succession (depth in metres, lithology, gamma ray API):
\begin{center}
\begin{tabular}{|c|c|c|}\hline
Depth (m) & Lithology & Gamma Ray (API)\\\hline
0–120 & Fine‑grained sandstone, cross‑bedded & 45\\\hline
120–250 & Dark grey shale, pyritic & 120\\\hline
250–340 & Bioturbated limestone, fossiliferous & 30\\\hline
340–410 & Coarse sandstone, conglomeratic base & 50\\\hline
410–500 & Red mudstone, desiccation cracks & 80\\\hline
\end{tabular}
\end{center}
Palynology from the shale (120–250\,m) yields \emph{Afropollis} spp. (Albian). Foraminifera from the limestone (250–340\,m) indicate Cenomanian–Turonian. Using the method, construct a stratigraphic column at 1\,cm = 50\,m, label formations using the Cameroonian nomenclature, and indicate the depositional environments.
\tcblower
\textbf{Detailed Solution:}
\begin{enumerate}
\item \textbf{Scale:} 1\,cm = 50\,m $\Rightarrow$ total column height = 500\,m / 50 = 10\,cm.
\item \textbf{Unit boundaries and thicknesses:}
\begin{itemize}
\item 0–120\,m: 120\,m thick (2.4\,cm) – Fine sandstone.
\item 120–250\,m: 130\,m thick (2.6\,cm) – Dark shale.
\item 250–340\,m: 90\,m thick (1.8\,cm) – Limestone.
\item 340–410\,m: 70\,m thick (1.4\,cm) – Coarse sandstone.
\item 410–500\,m: 90\,m thick (1.8\,cm) – Red mudstone.
\end{itemize}
\item \textbf{Formation assignment (Cameroon chart):}
\begin{itemize}
\item 0–120\,m: \cle{Bima Sandstone} (Early Cretaceous, fluvial–deltaic). Cross‑bedding, low GR.
\item 120–250\,m: \cle{Yola Formation} (Albian, marine shale). High GR, pyrite, \emph{Afropollis}.
\item 250–340\,m: \cle{Pindiga Formation} (Cenomanian–Turonian, shallow marine limestone). Low GR, bioturbation, foraminifera.
\item 340–410\,m: \cle{Gombe Sandstone} (Santonian–Campanian, fluvial channel). Conglomeratic base, moderate GR.
\item 410–500\,m: \cle{Kerri‑Kerri Formation} (Paleocene, continental red beds). Desiccation cracks, moderate GR.
\end{itemize}
\item \textbf{Key surfaces:}
\begin{itemize}
\item Base of Yola (120\,m) = Transgressive surface (TS).
\item Top of Pindiga (340\,m) = Maximum Flooding Surface (MFS).
\item Base of Gombe (340\,m) = Sequence Boundary (SB) – incised valley fill.
\end{itemize}
\item \textbf{Stratigraphic column (sketch):}
\begin{popfigure}
\begin{tikzpicture}[scale=0.6]
% Column frame
\draw[thick] (0,0) rectangle (3,10);
% Units
\fill[poporange!40] (0,0) rectangle (3,2.4);
\node[align=center] at (1.5,1.2) {Bima Sandstone\\120\,m\\Fluvial–deltaic};
\fill[popblue!30] (0,2.4) rectangle (3,5.0);
\node[align=center] at (1.5,3.7) {Yola Formation\\130\,m\\Marine shale (Albian)};
\fill[popgreen!30] (0,5.0) rectangle (3,6.8);
\node[align=center] at (1.5,5.9) {Pindiga Formation\\90\,m\\Shallow marine limestone};
\fill[poporange!40] (0,6.8) rectangle (3,8.2);
\node[align=center] at (1.5,7.5) {Gombe Sandstone\\70\,m\\Fluvial channel};
\fill[poppink!30] (0,8.2) rectangle (3,10);
\node[align=center] at (1.5,9.1) {Kerri‑Kerri Formation\\90\,m\\Continental red beds};
% Surfaces
\draw[dashed, thick, popdark] (0,2.4) -- (3,2.4) node[right] {TS};
\draw[dashed, thick, popdark] (0,6.8) -- (3,6.8) node[right] {MFS / SB};
% Scale
\draw[<->, thick] (3.5,0) -- (3.5,10) node[midway, right] {Depth (m)};
\foreach \y/\lab in {0/0, 2.4/120, 5.0/250, 6.8/340, 8.2/410, 10/500}
  \draw (3.3,\y) -- (3.7,\y) node[right] {\lab};
\end{tikzpicture}
\legende{Stratigraphic column for BH‑01, Logone‑Birni Basin (1\,cm = 50\,m).}
\end{popfigure}
\end{enumerate}
\textbf{Examiner's tip:} Always quote the biostratigraphic age \emph{and} the lithology to justify formation assignment; a single criterion is insufficient for full marks.
\end{exoresolu}

\begin{methode}[Interpreting Lake Nyos Gas‑Rich Water Column Stability]
\textbf{Objective:} Assess the risk of a limnic eruption at Lake Nyos using temperature, dissolved gas (CO$_2$), and density profiles.
\begin{enumerate}
\item \textbf{Obtain vertical profiles:} Measure temperature ($T$), conductivity, pH, and dissolved CO$_2$ concentration ($C_{\ce{CO2}}$) at 1\,m intervals from surface to bottom.
\item \textbf{Calculate water density ($\rho$):} Use the UNESCO equation of state for freshwater, incorporating $T$, pressure ($P$), and dissolved solids (approximated by conductivity). For CO$_2$-rich water, add the partial molar volume effect: $\rho = \rho_{\text{water}}(T,P) + \alpha C_{\ce{CO2}}$ where $\alpha \approx 0.001\,\mathrm{kg\,m^{-3}\,ppm^{-1}}$ (or per mmol\,kg$^{-1}$ with appropriate conversion).
\item \textbf{Plot $\rho(z)$:} Density must increase monotonically with depth for stable stratification. Any inversion (density decreasing with depth) signals potential overturn.
\item \textbf{Compute the Schmidt stability ($S$):} $S = \frac{1}{A_0}\int_0^{z_{\max}} (\rho(z)-\bar{\rho})\,g\,(z-\bar{z})\,dz$, where $A_0$ is surface area, $\bar{\rho}$ mean density, $\bar{z}$ mean depth. High $S$ ($>10^4\,\mathrm{J\,m^{-2}}$) indicates strong stability.
\item \textbf{Evaluate gas saturation:} Compare measured $C_{\ce{CO2}}$ with solubility at in‑situ $T$ and $P$ (Henry's law). Supersaturation ($C_{\ce{CO2}} > C_{\text{sat}}$) means free gas phase can nucleate.
\item \textbf{Risk classification:} 
\begin{itemize}
\item Stable density gradient + undersaturated $\rightarrow$ Low risk.
\item Stable gradient + supersaturated deep water $\rightarrow$ Moderate risk (degassing pipes recommended).
\item Density inversion + supersaturation $\rightarrow$ High risk (immediate mitigation).
\end{itemize}
\end{enumerate}
\textbf{Key reflex:} Always check for a \emph{chemocline} (sharp density jump due to dissolved CO$_2$) around 50–100\,m in Lake Nyos.
\end{methode}

\begin{exoresolu}[Stability Assessment of Lake Nyos Water Column]
\textbf{Statement:} A 2023 survey of Lake Nyos gave the following simplified data (depth $z$, temperature $T$, dissolved CO$_2$ concentration $C$):
\begin{center}
\begin{tabular}{|c|c|c|}\hline
$z$ (m) & $T$ ($^\circ$C) & $C$ (mmol\,kg$^{-1}$)\\\hline
0 & 24.5 & 0.5\\\hline
20 & 24.0 & 0.8\\\hline
40 & 23.5 & 2.5\\\hline
60 & 22.8 & 15.0\\\hline
80 & 22.0 & 45.0\\\hline
100 & 21.5 & 80.0\\\hline
120 & 21.0 & 110.0\\\hline
140 & 20.8 & 130.0\\\hline
160 & 20.5 & 140.0\\\hline
180 & 20.3 & 145.0\\\hline
200 & 20.2 & 148.0\\\hline
\end{tabular}
\end{center}
Assume freshwater density $\rho_w(T) \approx 1000 - 0.2(T-4)$ kg\,m$^{-3}$ (valid for 20–25$^\circ$C) and $\alpha = 0.001$ kg\,m$^{-3}$ per mmol\,kg$^{-1}$. 
\begin{enumerate}
\item Calculate $\rho(z)$ at each depth.
\item Determine whether the density profile is statically stable.
\item Estimate the CO$_2$ saturation concentration at 200\,m using Henry's law: $C_{\text{sat}} = k_H P_{\ce{CO2}}$ with $k_H = 0.034$ mol\,kg$^{-1}$\,MPa$^{-1}$ and $P_{\ce{CO2}} \approx 0.1$ MPa (partial pressure in pore water). Comment on supersaturation.
\item Classify the eruption risk.
\end{enumerate}
\tcblower
\textbf{Detailed Solution:}
\begin{enumerate}
\item \textbf{Density calculation:} $\rho(z) = \rho_w(T) + \alpha C$.
\begin{center}
\begin{tabular}{|c|c|c|c|}\hline
$z$ (m) & $\rho_w$ (kg\,m$^{-3}$) & $\alpha C$ (kg\,m$^{-3}$) & $\rho$ (kg\,m$^{-3}$)\\\hline
0 & $1000-0.2(20.5)=995.9$ & 0.0005 & 995.9005\\\hline
20 & $1000-0.2(20.0)=996.0$ & 0.0008 & 996.0008\\\hline
40 & $1000-0.2(19.5)=996.1$ & 0.0025 & 996.1025\\\hline
60 & $1000-0.2(18.8)=996.24$ & 0.0150 & 996.2550\\\hline
80 & $1000-0.2(18.0)=996.4$ & 0.0450 & 996.4450\\\hline
100 & $1000-0.2(17.5)=996.5$ & 0.0800 & 996.5800\\\hline
120 & $1000-0.2(17.0)=996.6$ & 0.1100 & 996.7100\\\hline
140 & $1000-0.2(16.8)=996.64$ & 0.1300 & 996.7700\\\hline
160 & $1000-0.2(16.5)=996.7$ & 0.1400 & 996.8400\\\hline
180 & $1000-0.2(16.3)=996.74$ & 0.1450 & 996.8850\\\hline
200 & $1000-0.2(16.2)=996.76$ & 0.1480 & 996.9080\\\hline
\end{tabular}
\end{center}
\item \textbf{Static stability:} $\rho$ increases monotonically from 995.9 to 996.9 kg\,m$^{-3}$ with depth. No inversion $\rightarrow$ \textbf{stably stratified}.
\item \textbf{CO$_2$ saturation at 200\,m:} $C_{\text{sat}} = k_H P_{\ce{CO2}} = 0.034 \times 0.1 = 0.0034$ mol\,kg$^{-1}$ = 3.4 mmol\,kg$^{-1}$. Measured $C = 148$ mmol\,kg$^{-1}$ $\gg$ $C_{\text{sat}}$. \textbf{Highly supersaturated} (factor ~43).
\item \textbf{Risk classification:} Stable density gradient + extreme supersaturation in deep water $\rightarrow$ \textbf{Moderate to High risk}. The lake is currently stable but stores enormous gas quantities; a trigger (landslide, earthquake) could initiate overturn. Degassing pipes (installed 2001) reduce risk but monitoring remains essential.
\end{enumerate}
\textbf{Examiner's tip:} Show the density calculation step‑by‑step; the monotonic increase is the key observation for stability. Always convert units consistently (mmol vs mol).
\end{exoresolu}

\begin{methode}[Recognising Major Structural Domains in Cameroon]
\textbf{Objective:} Identify and describe the three principal structural domains (Pan‑African fold belts, CVL rift structures, Benue Trough extensional faults) on a structural map and in cross‑section.
\begin{enumerate}
\item \textbf{Locate the domain boundaries:} 
\begin{itemize}
\item \cle{Pan‑African Belt (Central African Fold Belt)}: NE–SW trending folds and thrusts in the Adamawa–Yaoundé region (Neoproterozoic).
\item \cle{Cameroon Volcanic Line (CVL)}: NW–SE to SW–NE aligned volcanic centres, associated with extensional fractures and normal faults.
\item \cle{Benue Trough}: NW–SE trending rift basin with listric normal faults, half‑graben geometry, and Cretaceous sedimentary fill.
\end{itemize}
\item \textbf{Map symbols:} Use standard symbols: anticline/syncline axes, thrust teeth (upthrown side), normal fault ticks (downthrown side), strike‑slip arrows.
\item \textbf{Cross‑section construction:} 
\begin{enumerate}
\item Draw a topographic profile along a chosen traverse (e.g., SW–NE across the CVL).
\item Project surface geology and dip data onto the profile.
\item Restore stratigraphic thicknesses using known dips (apparent dip correction if section not perpendicular to strike).
\item Insert major faults with correct throw sense.
\end{enumerate}
\item \textbf{Kinematic indicators:} In the field, look for slickenlines, drag folds, Riedel shears, and fault gouge to confirm sense of movement.
\item \textbf{Regional synthesis:} Relate each domain to its geodynamic driver: Pan‑African collision (600–500\,Ma), Cretaceous Atlantic opening (Benue Trough), and Cenozoic mantle plume / lithospheric tearing (CVL).
\end{enumerate}
\textbf{Key formulae:} True dip $\delta$ from apparent dip $\alpha$ on a section at angle $\theta$ to strike: $\tan\delta = \tan\alpha / \cos\theta$. Apparent dip $\alpha$ from true dip $\delta$: $\tan\alpha = \tan\delta \cdot \cos\theta$.
\end{methode}

\begin{exoresolu}[Constructing a Structural Cross‑Section Across the CVL]
\textbf{Statement:} A SW–NE traverse across the Cameroon Volcanic Line (from the coast near Limbe to the Adamawa Plateau) cuts the following surface data (distance from coast, formation, dip direction/dip):
\begin{center}
\begin{tabular}{|c|c|c|}\hline
Distance (km) & Formation & Dip (dir/$^\circ$)\\\hline
0 & Basaltic lava flows (Mt Cameroon) & –/0 (horizontal)\\\hline
30 & Basaltic plateau (Buea) & NE/15\\\hline
60 & Granitic gneiss (Pan‑African) & NW/30\\\hline
90 & Sandstone (Cretaceous, Benue Trough edge) & SE/10\\\hline
120 & Basaltic cones (Ngaoundéré) & –/0\\\hline
\end{tabular}
\end{center}
The section line is oriented N45$^\circ$E (i.e., 45$^\circ$ from North). The strike of the Pan‑African foliation is N30$^\circ$E. 
\begin{enumerate}
\item Calculate the apparent dip of the Pan‑African foliation on the section plane.
\item Sketch a balanced cross‑section (vertical exaggeration 2:1) showing the major structures: Mt Cameroon edifice, the Pan‑African thrust belt, and the Benue Trough margin.
\item Indicate the sense of movement on the thrust fault that places gneiss over Cretaceous sandstone.
\end{enumerate}
\tcblower
\textbf{Detailed Solution:}
\begin{enumerate}
\item \textbf{Apparent dip calculation:} 
Strike of foliation = N30$^\circ$E. Section azimuth = N45$^\circ$E. Angle between section and strike $\theta = 45^\circ - 30^\circ = 15^\circ$.
True dip $\delta = 30^\circ$ (given as NW/30, dip direction NW = 315$^\circ$, but dip magnitude is 30$^\circ$).
Apparent dip $\alpha$: $\tan\alpha = \tan\delta \cdot \cos\theta = \tan30^\circ \cdot \cos15^\circ = 0.5774 \times 0.9659 = 0.5577$.
$\alpha = \arctan(0.5577) \approx 29.1^\circ$.
Since the section cuts the strike at a small angle, the apparent dip is close to true dip. The dip direction on the section will be towards the NW (i.e., towards the left on a SW–NE section).
\item \textbf{Cross‑section sketch (description for drawing):} 
\begin{popfigure}
\begin{tikzpicture}[scale=0.7]
% Topography
\draw[thick, popdark] (0,0) -- (2,0.5) -- (4,0.3) -- (6,0.8) -- (8,0.6) -- (10,0.9) -- (12,0.7) -- (14,1.0);
% Mt Cameroon edifice
\fill[poporange!40] (0,0) -- (2,0.5) -- (2,2) -- (0,2) -- cycle;
\node[align=center] at (1,1) {Mt Cameroon\\Basalts};
% Buea plateau
\fill[poporange!30] (2,0.5) -- (4,0.3) -- (4,1.5) -- (2,2) -- cycle;
\node[align=center] at (3,1) {Buea Plateau\\Basalts, dip NE 15$^\circ$};
% Pan-African thrust
\draw[thick, popblue, ->] (4,0.3) -- (6,0.8) node[midway, above, sloped] {Thrust};
\fill[popblue!20] (4,0.3) -- (6,0.8) -- (6,2.5) -- (4,2) -- cycle;
\node[align=center] at (5,1.5) {Pan‑African\\Gneiss, dip NW 30$^\circ$};
% Benue Trough margin
\draw[thick, popgreen, ->] (6,0.8) -- (8,0.6) node[midway, above, sloped] {Normal fault};
\fill[popgreen!20] (6,0.8) -- (8,0.6) -- (8,2) -- (6,2.5) -- cycle;
\node[align=center] at (7,1.5) {Cretaceous\\Sandstone, dip SE 10$^\circ$};
% Ngaoundere cones
\fill[poporange!40] (10,0.9) -- (12,0.7) -- (12,2.2) -- (10,2.5) -- cycle;
\node[align=center] at (11,1.5) {Ngaoundéré\\Basaltic cones};
% Scale
\draw[<->] (0,-0.8) -- (14,-0.8) node[midway, below] {Distance (km)};
\foreach \x/\lab in {0/0, 2/30, 4/60, 6/90, 8/120, 10/150, 12/180, 14/210}
  \draw (\x,-0.7) -- (\x,-0.9) node[below] {\lab};
% Vertical exaggeration note
\node[align=center, popmuted] at (7,-1.5) {Vertical exaggeration 2:1};
\end{tikzpicture}
\legende{Schematic balanced cross‑section SW–NE across the CVL (not to scale).}
\end{popfigure}
\item \textbf{Thrust sense:} The Pan‑African gneiss (older, metamorphic) overrides the younger Cretaceous sandstone. On the cross‑section, the thrust fault dips NW (towards the left) and the hanging wall (gneiss) moves up and over the footwall (sandstone). This is a \textbf{thrust (reverse) fault} with top‑to‑the‑NW movement.
\end{enumerate}
\textbf{Examiner's tip:} Always state the \emph{angle between section line and strike} before computing apparent dip. In balanced sections, ensure area (or volume) conservation for each stratigraphic unit.
\end{exoresolu}

\begin{methode}[Evaluating Economic Mineral Potential of a Lateritic Profile]
\textbf{Objective:} Assess the lateritic bauxite/iron ore potential of a weathering profile on the Adamawa Plateau using geochemical and mineralogical criteria.
\begin{enumerate}
\item \textbf{Describe the profile:} Identify horizons from top to bottom: \cle{Ferricrete} (duricrust), \cle{Mottled zone}, \cle{Pallid zone}, \cle{Saprolite}, \cle{Fresh bedrock}. Measure thickness of each.
\item \textbf{Collect samples:} Channel samples every 0.5–1\,m. Analyse for major oxides (XRF) and trace elements (ICP‑MS). Determine mineralogy by XRD (gibbsite, boehmite, hematite, goethite, kaolinite).
\item \textbf{Calculate key indices:}
\begin{itemize}
\item \textbf{Al$_2$O$_3$ grade:} Weight \% Al$_2$O$_3$ in the ore zone (usually mottled + pallid zones).
\item \textbf{SiO$_2$ / Al$_2$O$_3$ ratio:} Low ratio (< 0.5) favours bauxite; high ratio indicates kaolinitic laterite.
\item \textbf{Fe$_2$O$_3$ / Al$_2$O$_3$ ratio:} Distinguishes ferruginous bauxite (> 0.5) from lateritic iron ore (> 2).
\item \textbf{Weathering index (Ki):} $Ki = \frac{\text{SiO}_2/\text{Al}_2\text{O}_3}{\text{Fe}_2\text{O}_3/\text{Al}_2\text{O}_3}$ (molar ratios). Ki > 10 $\rightarrow$ bauxite; 1–10 $\rightarrow$ lateritic iron; < 1 $\rightarrow$ ferricrete.
\end{itemize}
\item \textbf{Estimate tonnage:} Volume = area $\times$ average ore thickness $\times$ (1 – porosity). Tonnage = volume $\times$ bulk density (typically 2.2–2.5\,t\,m$^{-3}$ for laterite).
\item \textbf{Check infrastructure \& environmental constraints:} Proximity to rail/port (Douala, Kribi), power, water, and protected areas.
\item \textbf{Classification (UNFC/JORC):} Assign resource category (Inferred, Indicated, Measured) based on drill spacing and data quality.
\end{enumerate}
\textbf{Key reflex:} The \emph{pallid zone} (white, gibbsite‑rich) is the prime bauxite target; the \emph{mottled zone} (red‑white) often hosts iron ore.
\end{methode}

\begin{exoresolu}[Resource Estimation for a Bauxite Laterite on the Adamawa Plateau]
\textbf{Statement:} A lateritic profile on the Adamawa Plateau (area 12\,km$^2$) has been drilled on a 200\,m grid (36 holes). The average thicknesses and grades of the ore zone (mottled + pallid zones combined) are:
\begin{itemize}
\item Average ore thickness = 8.5\,m (standard deviation 1.2\,m).
\item Mean Al$_2$O$_3$ = 42.3\,\%, SiO$_2$ = 3.8\,\%, Fe$_2$O$_3$ = 18.5\,\%, TiO$_2$ = 2.1\,\%.
\item Bulk density = 2.35\,t\,m$^{-3}$.
\item Porosity = 12\,\%.
\end{itemize}
Drill spacing qualifies for \textbf{Indicated} resources. 
\begin{enumerate}
\item Compute the SiO$_2$/Al$_2$O$_3$ and Fe$_2$O$_3$/Al$_2$O$_3$ ratios (weight basis) and the molar Ki index. Classify the laterite type.
\item Estimate the Indicated tonnage (in million tonnes) and contained Al$_2$O$_3$ metal (in million tonnes).
\item Discuss two major non‑geological factors that could affect project viability in Cameroon.
\end{enumerate}
\tcblower
\textbf{Detailed Solution:}
\begin{enumerate}
\item \textbf{Ratios and Ki index:}
\begin{align*}
\text{SiO}_2/\text{Al}_2\text{O}_3 &= 3.8 / 42.3 = 0.0898 \quad (\text{weight})\\
\text{Fe}_2\text{O}_3/\text{Al}_2\text{O}_3 &= 18.5 / 42.3 = 0.437 \quad (\text{weight})
\end{align*}
Molar masses: SiO$_2$=60.08, Al$_2$O$_3$=101.96, Fe$_2$O$_3$=159.69 g\,mol$^{-1}$.
Molar ratios:
\begin{align*}
\text{SiO}_2/\text{Al}_2\text{O}_3 (\text{mol}) &= \frac{3.8/60.08}{42.3/101.96} = \frac{0.06326}{0.4149} = 0.1525\\
\text{Fe}_2\text{O}_3/\text{Al}_2\text{O}_3 (\text{mol}) &= \frac{18.5/159.69}{42.3/101.96} = \frac{0.1158}{0.4149} = 0.2792\\
Ki &= \frac{0.1525}{0.2792} = 0.546
\end{align*}
Since $Ki < 1$, the profile is classified as \textbf{ferricrete / lateritic iron ore} rather than high‑grade bauxite. However, the low SiO$_2$/Al$_2$O$_3$ (0.09) suggests a gibbsite‑rich component; the high Fe$_2$O$_3$ lowers Ki.
\item \textbf{Tonnage calculation:}
\begin{align*}
\text{Area} &= 12\ \text{km}^2 = 12 \times 10^6\ \text{m}^2\\
\text{Volume (solid)} &= \text{Area} \times \text{thickness} \times (1 - \text{porosity})\\
&= 12 \times 10^6 \times 8.5 \times (1 - 0.12)\\
&= 12 \times 10^6 \times 8.5 \times 0.88\\
&= 89.76 \times 10^6\ \text{m}^3\\
\text{Tonnage} &= \text{Volume} \times \text{bulk density}\\
&= 89.76 \times 10^6 \times 2.35\\
&= 210.9 \times 10^6\ \text{t} \approx \mathbf{211\ \text{million tonnes (Mt)}}
\end{align*}
Contained Al$_2$O$_3$ metal:
$211\ \text{Mt} \times 0.423 = \mathbf{89.3\ \text{Mt Al}_2\text{O}_3}$.
\item \textbf{Non‑geological factors (Cameroon context):}
\begin{enumerate}
\item \textbf{Transport infrastructure:} The Adamawa Plateau is landlocked; the nearest deep‑water port is Kribi (~600\,km) or Douala (~550\,km). The existing railway (Transcam) ends at Ngaoundéré; extension or heavy‑haul road upgrades are capital‑intensive.
\item \textbf{Energy supply:} Bauxite refining (Bayer process) requires abundant, low‑cost electricity. Cameroon's grid (mainly hydro) faces dry‑season deficits; a dedicated power plant (gas or hydro) would be needed, increasing CAPEX.
\end{enumerate}
Other factors: fiscal regime (mining code 2001, revised 2023), community land rights, and environmental permitting (protected areas nearby).
\end{enumerate}
\textbf{Examiner's tip:} Always convert weight \% to molar ratios for Ki; using weight ratios gives a different (and incorrect) index. Show all unit conversions.
\end{exoresolu}

\begin{methode}[Integrating Multi‑Disciplinary Data for Geohazard Zonation (Landslide Susceptibility)]
\textbf{Objective:} Produce a qualitative landslide susceptibility map for the Mount Cameroon region using geological, geomorphological, and climatic parameters.
\begin{enumerate}
\item \textbf{Define the study area:} Delineate the watershed(s) on a DEM (30\,m SRTM or 12.5\,m ALOS).
\item \textbf{Prepare thematic layers (GIS):}
\begin{itemize}
\item \textit{Lithology:} Competent (basalt, trachyte) vs weak (pyroclastics, weathered tuff, laterite).
\item \textit{Slope angle:} Derive from DEM; classify: $<15^\circ$ (low), $15–30^\circ$ (moderate), $>30^\circ$ (high).
\item \textit{Slope aspect:} South‑facing slopes receive more rainfall (orographic effect).
\item \textit{Drainage density:} High density $\rightarrow$ higher erosion potential.
\item \textit{Lineament density:} From satellite imagery (faults, fractures) – structural control.
\item \textit{Rainfall:} Mean annual precipitation (MAP) from meteorological stations; > 3000\,mm/yr = very high.
\item \textit{Land use/cover:} Forest (stabilising) vs agriculture/bare soil (destabilising).
\end{itemize}
\item \textbf{Assign weights (Analytic Hierarchy Process or expert judgment):} Typical weights: Lithology 25\%, Slope 20\%, Rainfall 15\%, Lineaments 15\%, Land use 15\%, Drainage 10\%.
\item \textbf{Calculate susceptibility index (SI) per pixel:} $SI = \sum w_i \times r_i$, where $r_i$ is the class rating (1–5) for each layer.
\item \textbf{Classify SI:} Very Low, Low, Moderate, High, Very High (natural breaks or equal interval).
\item \textbf{Validate:} Overlay known landslide inventory (historical events, field mapping). Compute success rate curve (AUC).
\item \textbf{Produce map:} Layout with legend, scale, north arrow, coordinate system (WGS84 / UTM Zone 32N).
\end{enumerate}
\textbf{Key reflex:} In Cameroon, the \emph{pyroclastic deposits} on Mount Cameroon's flanks are the weakest lithological unit and coincide with highest rainfall – they dominate the susceptibility model.
\end{methode}

\begin{exoresolu}[Landslide Susceptibility Index Calculation for a Mount Cameroon Slope]
\textbf{Statement:} For a 1\,km$^2$ pixel on the southwestern flank of Mount Cameroon, the following thematic class ratings (1 = low susceptibility, 5 = high) and weights are given:
\begin{center}
\begin{tabular}{|l|c|c|}\hline
Parameter & Weight ($w_i$) & Class Rating ($r_i$)\\\hline
Lithology (weathered pyroclastics) & 0.25 & 5\\\hline
Slope angle (35$^\circ$) & 0.20 & 5\\\hline
Rainfall (3500\,mm/yr) & 0.15 & 5\\\hline
Lineament density (high) & 0.15 & 4\\\hline
Land use (subsistence farming) & 0.15 & 4\\\hline
Drainage density (moderate) & 0.10 & 3\\\hline
\end{tabular}
\end{center}
\begin{enumerate}
\item Compute the Susceptibility Index (SI) for this pixel.
\item Classify the pixel using the scheme: SI $<$ 2.0 = Very Low; 2.0–2.5 = Low; 2.5–3.0 = Moderate; 3.0–3.5 = High; $>$ 3.5 = Very High.
\item If the pixel area is 1\,km$^2$ and the average regolith thickness is 12\,m, estimate the potential landslide volume (in Mm$^3$) assuming a 30\% mobilisation ratio.
\item Suggest two mitigation measures appropriate for the local context.
\end{enumerate}
\tcblower
\textbf{Detailed Solution:}
\begin{enumerate}
\item \textbf{SI calculation:}
\begin{align*}
SI &= \sum w_i r_i \\
&= (0.25 \times 5) + (0.20 \times 5) + (0.15 \times 5) + (0.15 \times 4) + (0.15 \times 4) + (0.10 \times 3)\\
&= 1.25 + 1.00 + 0.75 + 0.60 + 0.60 + 0.30\\
&= \mathbf{4.50}
\end{align*}
\item \textbf{Classification:} SI = 4.50 $>$ 3.5 $\rightarrow$ \textbf{Very High susceptibility}.
\item \textbf{Potential landslide volume:}
\begin{align*}
\text{Regolith volume} &= \text{Area} \times \text{thickness} = 1\ \text{km}^2 \times 12\ \text{m}\\
&= 1 \times 10^6\ \text{m}^2 \times 12\ \text{m} = 12 \times 10^6\ \text{m}^3 = 12\ \text{Mm}^3\\
\text{Mobilised volume} &= 0.30 \times 12\ \text{Mm}^3 = \mathbf{3.6\ \text{Mm}^3}
\end{align*}
\item \textbf{Mitigation measures (Cameroon context):}
\begin{enumerate}
\item \textbf{Bio‑engineering:} Plant deep‑rooted vetiver grass (\emph{Chrysopogon zizanioides}) and native trees (e.g., \emph{Prunus africana}) on contour lines to reinforce soil cohesion. Low cost, community‑driven, proven on Mount Cameroon slopes.
\item \textbf{Surface drainage control:} Construct lined diversion ditches and check‑dams using local basalt blocks to intercept runoff and reduce pore‑water pressure in the pyroclastic regolith. Integrate with existing agricultural terraces.
\end{enumerate}
\end{enumerate}
\textbf{Examiner's tip:} Show the weighted sum explicitly; examiners check each multiplication. For volume, remember 1\,km$^2$ = $10^6$\,m$^2$. Mobilisation ratio is a judgement – state your assumption.
\end{exoresolu}

\begin{methode}[Correlating Stratigraphic Units Across Cameroonian Basins Using Sequence Stratigraphy]
\textbf{Objective:} Correlate the Cretaceous–Cenozoic successions of the \cle{Logone‑Birni}, \cle{Douala}, and \cle{Rio del Rey} basins using sequence stratigraphic surfaces.
\begin{enumerate}
\item \textbf{Identify key surfaces in each basin:} 
\begin{itemize}
\item \cle{Base Cretaceous Unconformity (BCU)} – on basement.
\item \cle{Aptian–Albian Transgressive Surface (TS)} – marine incursion.
\item \cle{Cenomanian–Turonian Maximum Flooding Surface (MFS)} – condensed shales/limestones.
\item \cle{Santonian–Campanian Sequence Boundary (SB)} – forced regression.
\item \cle{Paleocene–Eocene MFS} – marine transgression in coastal basins.
\end{itemize}
\item \textbf{Use biostratigraphy:} Palynomorphs (e.g., \emph{Afropollis}, \emph{Classopollis}), foraminifera (e.g., \emph{Rotalipora}, \emph{Globotruncana}), nannofossils to date surfaces.
\item \textbf{Wireline log correlation:} Gamma ray (shale) and resistivity (sand/carbonate) patterns allow subsurface correlation where fossils are absent.
\item \textbf{Construct a correlation panel:} Align basins along a SW–NE transect (Douala – Rio del Rey – Logone‑Birni). Plot chronostratigraphic chart with systems tracts (LST, TST, HST).
\item \textbf{Interpret controls:} Eustasy (global sea‑level curves), tectonics (rift subsidence, CVL uplift), sediment supply (Congo River, Niger Delta).
\item \textbf{Document lateral facies changes:} Coastal basins (Douala, Rio del Rey) show thick marine shales and turbidites; inland basin (Logone‑Birni) shows fluvial–deltaic sandstones and lacustrine shales.
\end{enumerate}
\textbf{Key reflex:} The \emph{Cenomanian–Turonian MFS} is the most regionally extensive correlatable surface in West‑Central Africa; use it as the primary datum.
\end{methode}

\begin{exoresolu}[Sequence Stratigraphic Correlation Across Three Cameroonian Basins]
\textbf{Statement:} You are given simplified stratigraphic logs for three wells: Well A (Douala Basin, offshore), Well B (Rio del Rey Basin, onshore), Well C (Logone‑Birni Basin, inland). Each log shows gamma ray (GR) and lithology. Key biostratigraphic markers are:
\begin{itemize}
\item Well A: \emph{Rotalipora cushmani} (Cenomanian) at 2450\,m; \emph{Globotruncana} spp. (Turonian) at 2300\,m.
\item Well B: \emph{Afropollis jardinus} (Albian) at 1800\,m; \emph{Classopollis} (Cenomanian) at 1650\,m.
\item Well C: \emph{Afropollis} spp. (Albian) at 250\,m; \emph{Classopollis} (Cenomanian) at 200\,m.
\end{itemize}
The GR logs show a prominent low‑GR (clean sand) package at the top of the Albian in all wells, overlain by a high‑GR (shale) package in Wells A and B, but a low‑GR (sandstone) package in Well C.
\begin{enumerate}
\item Identify the sequence stratigraphic surfaces (TS, MFS, SB) in each well using the biostratigraphy and GR patterns.
\item Draw a correlation panel (schematic) aligning the three wells on the Cenomanian–Turonian MFS.
\item Explain the facies difference above the MFS between the coastal basins (A, B) and the inland basin (C).
\end{enumerate}
\tcblower
\textbf{Detailed Solution:}
\begin{enumerate}
\item \textbf{Surface identification:}
\begin{itemize}
\item \textbf{Albian Transgressive Surface (TS):} Marked by the base of the clean sand package (low GR) in all wells. In Well A at ~2500\,m, Well B at ~1850\,m, Well C at ~260\,m. This is the \emph{Aptian–Albian marine incursion}.
\item \textbf{Cenomanian–Turonian MFS:} In Wells A and B, the high‑GR shale package (condensed section) with \emph{Rotalipora}/\emph{Classopollis} represents the MFS (maximum marine flooding). In Well C, the MFS is represented by a thin, fossiliferous shale interbed within the sandstone at ~200\,m (condensed but sandy due to proximity to shore).
\item \textbf{Santonian SB:} Not directly dated here, but would be a basinward shift of facies (incised valleys) above the Turonian.
\end{itemize}
\item \textbf{Correlation panel (schematic):}
\begin{popfigure}
\begin{tikzpicture}[scale=0.65]
% Well A
\draw[thick] (0,0) -- (0,8);
\fill[poporange!30] (0,0) rectangle (1.5,2); \node[align=center] at (0.75,1) {Basement};
\fill[popblue!30] (0,2) rectangle (1.5,4); \node[align=center] at (0.75,3) {Albian Sand (LST)};
\fill[popgreen!30] (0,4) rectangle (1.5,5.5); \node[align=center] at (0.75,4.75) {Cenom–Turon MFS\\Shale};
\fill[poporange!30] (0,5.5) rectangle (1.5,7); \node[align=center] at (0.75,6.25) {Turon–Sant HST};
\fill[poppink!30] (0,7) rectangle (1.5,8); \node[align=center] at (0.75,7.5) {Paleogene};
\node[align=center] at (0.75,-0.5) {Well A (Douala)};
% Well B
\draw[thick] (3,0) -- (3,8);
\fill[poporange!30] (3,0) rectangle (4.5,1.5); \node[align=center] at (3.75,0.75) {Basement};
\fill[popblue!30] (3,1.5) rectangle (4.5,3.5); \node[align=center] at (3.75,2.5) {Albian Sand (LST)};
\fill[popgreen!30] (3,3.5) rectangle (4.5,5); \node[align=center] at (3.75,4.25) {Cenom–Turon MFS\\Shale};
\fill[poporange!30] (3,5) rectangle (4.5,6.5); \node[align=center] at (3.75,5.75) {Turon–Sant HST};
\fill[poppink!30] (3,6.5) rectangle (4.5,8); \node[align=center] at (3.75,7.25) {Paleogene};
\node[align=center] at (3.75,-0.5) {Well B (Rio del Rey)};
% Well C
\draw[thick] (6,0) -- (6,8);
\fill[poporange!30] (6,0) rectangle (7.5,0.8); \node[align=center] at (6.75,0.4) {Basement};
\fill[popblue!30] (6,0.8) rectangle (7.5,2.5); \node[align=center] at (6.75,1.65) {Albian Sand (LST)};
\fill[popgreen!30] (6,2.5) rectangle (7.5,3.2); \node[align=center] at (6.75,2.85) {Cenom–Turon MFS\\Thin shale};
\fill[popblue!30] (6,3.2) rectangle (7.5,5.5); \node[align=center] at (6.75,4.35) {Cenom–Turon HST\\Sandstone};
\fill[poporange!30] (6,5.5) rectangle (7.5,7); \node[align=center] at (6.75,6.25) {Santonian–Camp\\Fluvial};
\fill[poppink!30] (6,7) rectangle (7.5,8); \node[align=center] at (6.75,7.5) {Paleogene};
\node[align=center] at (6.75,-0.5) {Well C (Logone‑Birni)};
% Correlation lines
\draw[dashed, thick, popdark] (0,4) -- (7.5,2.85) node[midway, above, sloped] {MFS};
\draw[dashed, thick, popdark] (0,2) -- (7.5,1.65) node[midway, below, sloped] {TS};
% Scale
\draw[<->] (0,-1.2) -- (7.5,-1.2) node[midway, below] {Distance (approx. 300 km)};
\end{tikzpicture}
\legende{Schematic correlation panel aligned on the Cenomanian–Turonian MFS.}
\end{popfigure}
\item \textbf{Facies difference explanation:}
\begin{itemize}
\item \textbf{Coastal basins (A, B):} During the Cenomanian–Turonian highstand, the shoreline was far inland; the basins received distal, fine‑grained pelagic/hemipelagic sediments (shales, marls) forming a thick condensed section at the MFS.
\item \textbf{Inland basin (C):} The Logone‑Birni Basin was a shallow, restricted epicontinental sea or large lake. The MFS corresponds to a brief deepening but still within the reach of fluvial sediment supply; thus the condensed section is thin and sandy (shoreface/delta front). The overlying HST is dominantly fluvial–deltaic sandstone, not marine shale.
\end{itemize}
\end{enumerate}
\textbf{Examiner's tip:} Always align correlation panels on a \emph{chronostratigraphic surface} (MFS), not on lithological tops. The MFS is time‑synchronous; lithological tops are diachronous.
\end{exoresolu}

\newpage

\coursec{Exercises}

\courssub{I check my knowledge}

\begin{exercice}[MCQ: Cameroon Volcanic Line]
Which of the following statements about the Cameroon Volcanic Line (CVL) is \textbf{incorrect}?
\begin{enumerate}
\item The CVL extends from the Atlantic Ocean (Pagalu/Annobón island) to Lake Chad (Biu Plateau).
\item The oceanic sector includes the islands of Bioko, São Tomé, Príncipe and Pagalu.
\item The continental sector includes Mount Cameroon, Manengouba, Bamboutos, Oku and Mandara Mountains.
\item The dominant rock types are basalts, trachytes, phonolites and rhyolites.
\item The CVL is a classic example of a hotspot track formed by the African plate moving over a fixed mantle plume.
\end{enumerate}
\end{exercice}

\begin{exercice}[True/False with Justification]
State whether each statement is True or False. Justify your answer in one sentence.
\begin{enumerate}
\item The Ntem Complex in southern Cameroon consists of Archaean granitoids and greenstone belts similar to the Kaapvaal Craton.
\item The Douala Basin (Kribi-Campo sub-basin) contains Cretaceous to Recent sediments deposited during the opening of the South Atlantic.
\item The Pan-African orogeny (600–500 Ma) affected only the northern part of Cameroon (Adamawa-Yadé domain).
\item Lake Nyos is a maar lake formed by a phreatomagmatic explosion; its 1986 disaster was caused by a sudden overturn releasing dissolved CO\textsubscript{2}.
\item Lateritic bauxite deposits in Cameroon (e.g. Minim-Martap, Ngaoundal) formed by intense tropical weathering of Precambrian basement rocks.
\end{enumerate}
\end{exercice}

\begin{exercice}[Definitions]
Define the following geological terms and give a specific Cameroonian example for each:
\begin{enumerate}
\item Limnic eruption
\item Shear zone (or fault zone)
\item Lateritic bauxite
\item Angular unconformity
\item Pan-African granitoids
\end{enumerate}
\end{exercice}

\begin{exercice}[Direct Calculation: Lake Nyos CO\textsubscript{2} Profile]
The table below shows dissolved CO\textsubscript{2} concentration (in mmol/kg) measured at different depths in Lake Nyos in March 2023.

\begin{center}
\begin{tabular}{|c|c|c|c|c|c|c|}
\hline
Depth (m) & 0 & 50 & 100 & 150 & 200 & 207 (bottom) \\
\hline
CO\textsubscript{2} (mmol/kg) & 0.2 & 0.5 & 2.1 & 15.8 & 45.3 & 52.7 \\
\hline
\end{tabular}
\end{center}
\begin{enumerate}
\item Plot the CO\textsubscript{2} concentration vs. depth profile on graph paper (depth increasing downwards).
\item Calculate the approximate total mass of CO\textsubscript{2} (in tonnes) stored in the bottom 50 m layer (150–207 m), assuming the lake has a uniform cross-sectional area of $1.2 \times 10^5\ \mathrm{m}^2$ in this depth range, water density $1000\ \mathrm{kg\,m^{-3}}$, and molar mass of CO\textsubscript{2} = $44\ \mathrm{g\,mol^{-1}}$. Use the average concentration of the 150 m and 207 m values.
\end{enumerate}
\end{exercice}

\courssub{I practise}

\begin{exercice}[Structured Question: Cameroon Volcanic Line]
\begin{enumerate}
\item Describe the geographical extent of the Cameroon Volcanic Line (CVL), naming its oceanic and continental sectors and at least three major volcanoes in each sector. [4 marks]
\item Draw a labelled sketch cross-section showing the typical structure of a CVL stratovolcano (e.g. Mount Cameroon) including lava flows, pyroclastic layers, a central vent, and a parasitic cone. [4 marks]
\item Discuss \textbf{two} hypotheses for the origin of the CVL (e.g. mantle plume/hotspot vs. lithospheric fracture/structural control). For each hypothesis, give one supporting argument and one counter-argument based on geological/geophysical evidence. [6 marks]
\item Explain why Mount Cameroon is considered the most active volcano in the CVL, citing its eruption history and lava composition. [3 marks]
\end{enumerate}
\end{exercice}

\begin{exercice}[Map Exercise: Geological Features of Cameroon]
On the outline map of Cameroon provided below (or draw a sketch map), locate and label the following features:
\begin{enumerate}
\item Ntem Complex (Archaean basement)
\item Douala Basin (Kribi-Campo sub-basin)
\item Mamfe Basin (Cretaceous rift basin)
\item Sanaga Fault (major shear zone)
\item Lake Nyos (crater lake, Northwest Region)
\item Mbalam iron ore deposit (East Region)
\item Mount Cameroon (active stratovolcano)
\item Bamboutos Mountains (caldera complex, West Region)
\end{enumerate}
Choose \textbf{two} of the features above and justify their economic importance to Cameroon in 3–4 sentences each. [12 marks total]
\end{exercice}

\begin{exercice}[Data Analysis: Lake Nyos Degassing]
The degassing pipe installed in Lake Nyos (2001) and the additional pipes (2011) aim to prevent limnic eruptions.
\begin{enumerate}
\item Explain the physical principle behind the self-driven degassing pipe (no external energy required). Use the concept of gas lift and density difference. [4 marks]
\item Given that the pipe has an inner diameter of $0.14\ \mathrm{m}$ and the two-phase mixture (water + CO\textsubscript{2} bubbles) rises at an average velocity of $2.5\ \mathrm{m\,s^{-1}}$, calculate the volumetric flow rate of the mixture in $\mathrm{m}^3\,\mathrm{s}^{-1}$ and the mass flow rate of water in $\mathrm{kg\,s^{-1}}$ (assume mixture density $\approx 900\ \mathrm{kg\,m^{-3}}$). [3 marks]
\item If the deep water contains $50\ \mathrm{mmol\,kg^{-1}}$ of dissolved CO\textsubscript{2}, estimate the mass of CO\textsubscript{2} removed per day (in tonnes/day) by one pipe operating continuously. [3 marks]
\item Why is it critical to maintain the pipe inlet at a specific depth (e.g. 60–70 m below surface) rather than at the very bottom? [2 marks]
\end{enumerate}
\end{exercice}

\begin{exercice}[Stratigraphic Column Construction]
Arrange the following Cameroonian rock units into a correct stratigraphic column from oldest (bottom) to youngest (top). For each unit, state the geological Era/Period and the main rock types.
\begin{itemize}
\item Quaternary basalts of Mount Cameroon
\item Cretaceous sandstones and shales of the Mamfe Basin
\item Pan-African granites (syn- to post-tectonic, ~550 Ma)
\item Nyong Group (metasediments: schists, quartzites, banded iron formations)
\item Ntem Complex gneisses and migmatites (Archaean basement)
\item Cenozoic volcanic rocks of the Bamboutos caldera (Miocene–Pliocene)
\item Jurassic–Cretaceous basalts of the Douala Basin (rift volcanics)
\end{itemize}
Present your answer as a table with columns: Stratigraphic Position, Unit Name, Age (Era/Period), Main Rock Types, Tectonic Setting. [14 marks]
\end{exercice}

\courssub{I prepare for the GCE}

\begin{exercice}[Essay: Geology of Cameroon — Blessing and Threat]
\textbf{``The geology of Cameroon is both a blessing and a threat.''} Discuss this statement with reference to:
\begin{itemize}
\item Mineral resources (at least three types, e.g. bauxite, iron ore, gold, diamonds, hydrocarbons, limestone, cobalt/nickel) — their geological setting, location and economic potential. [10 marks]
\item Geologic hazards (at least three types, e.g. volcanic eruptions, limnic eruptions, landslides, earthquakes, flooding/erosion) — their causes, recent examples and mitigation measures. [10 marks]
\item Sustainable management: how can geological knowledge (mapping, monitoring, land-use planning) reduce risks while optimising resource exploitation? [5 marks]
\end{itemize}
Total: [25 marks]
\end{exercice}

\begin{exercice}[Problem Solving: Village Relocation near Bamboutos Mountains]
A village of 500 inhabitants is to be relocated due to high landslide risk on the steep slopes of the Bamboutos Mountains (West Region). Three potential sites (A, B, C) have been identified. Geological and climatic data are summarised below.

\begin{center}
\begin{tabular}{|l|c|c|c|}
\hline
\textbf{Parameter} & \textbf{Site A} & \textbf{Site B} & \textbf{Site C} \\
\hline
Underlying geology & Weathered basalt (laterite >15 m) & Fractured trachyte with clay gouge & Alluvial terrace (sand/gravel, 8 m) over impermeable shale \\
Slope angle & $12^\circ$ & $28^\circ$ & $3^\circ$ \\
Annual rainfall & $2200\ \mathrm{mm}$ & $2200\ \mathrm{mm}$ & $2200\ \mathrm{mm}$ \\
Rainy season intensity & High (convective storms) & High & High \\
Distance to river & $800\ \mathrm{m}$ & $150\ \mathrm{m}$ & $50\ \mathrm{m}$ (floodplain) \\
Evidence of past instability & None observed & Tension cracks, tilted trees & Old flood deposits, waterlogging \\
\hline
\end{tabular}
\end{center}

\begin{enumerate}
\item Using the data, assess the landslide and flood risk for each site. Rank the sites from safest to most hazardous. Justify your ranking with reference to specific geological and hydrological factors. [12 marks]
\item Recommend the best site for relocation. Propose two engineering or land-use measures to further reduce risk at the chosen site. [5 marks]
\item Explain how the volcanic history of the Bamboutos (caldera collapse, pyroclastic deposits, hydrothermal alteration) influences the current stability of slopes in the area. [5 marks]
\end{enumerate}
Total: [22 marks]
\end{exercice}

\begin{exercice}[Integration: Stratigraphy, Structure and Resources]
The figure below (not shown) represents a simplified geological cross-section across southern Cameroon from the Congo Craton (south) to the Adamawa Plateau (north), passing through the Ntem Complex, the Pan-African belt, the Sanaga Fault, and the Douala Basin.

\begin{enumerate}
\item Describe the major tectonostratigraphic domains crossed by this section, from south to north, giving the age range and dominant rock types for each domain. [8 marks]
\item The Sanaga Fault is a major NE–SW trending shear zone. Explain its role in:
  \begin{enumerate}
  \item The Pan-African orogeny (600–500 Ma)
  \item The Cretaceous opening of the South Atlantic and formation of the Douala Basin
  \item The Cenozoic magmatism of the CVL (e.g. Mount Cameroon alignment)
  \end{enumerate}
  [6 marks]
\item The Lom Basin (East Cameroon) and the Mamfe Basin (Southwest) are both Cretaceous rift basins. Compare their stratigraphy, depositional environments and hydrocarbon potential. [6 marks]
\item The Mbalam iron deposit (East Region) is hosted in Banded Iron Formation (BIF) of the Nyong Group. Explain the genetic model for this BIF (Archaean/Palaeoproterozoic) and why it is now economically viable. [5 marks]
\end{enumerate}
Total: [25 marks]
\end{exercice}

\newpage

\coursec{Solutions to the exercises}

\begin{corrige}[Exercise 1 — MCQ: Cameroon Volcanic Line]
\textbf{Correct answer: (e).}

\begin{itemize}
\item (a) \textbf{Correct}: the CVL is a Y-shaped alignment stretching about \SI{1600}{km}, from Pagalu (Annobón) island in the Gulf of Guinea through the continental sector to the Biu Plateau near Lake Chad in Nigeria.
\item (b) \textbf{Correct}: the oceanic sector comprises Pagalu, São Tomé, Príncipe and Bioko.
\item (c) \textbf{Correct}: the continental sector includes Mount Cameroon, Rumpi Hills, Manengouba, Bamboutos, Oku, Ngaoundéré and the Mandara Mountains.
\item (d) \textbf{Correct}: CVL magmas are alkaline, dominated by basalts and basanites with differentiates (trachytes, phonolites, rhyolites).
\item (e) \textbf{Incorrect}: the CVL is \emph{not} a classic age-progressive hotspot track. Radiometric ages do \textbf{not} increase systematically along the line (volcanism is roughly coeval along the chain, with no clear age progression), which contradicts the simple plume/plate-motion model. This is why alternative models (lithospheric reactivation of Pan-African shear zones, edge-driven convection, decompression melting) are invoked.
\end{itemize}

\begin{attention}
A classic trap: because Hawaii is the textbook hotspot track, candidates assume every linear volcanic chain is one. The GCE examiner expects you to know that the \emph{absence of age progression} is the key argument against the simple plume model for the CVL.
\end{attention}
\end{corrige}

\begin{corrige}[Exercise 2 — True/False with Justification]
\begin{enumerate}
\item \textbf{True.} The Ntem Complex (south Cameroon, northern edge of the Congo Craton) comprises Archaean ($\sim 2.9$–$2.5$ Ga) TTG granitoids, charnockites and greenstone belts, comparable in age and lithology to other Archaean cratons such as Kaapvaal.
\item \textbf{True.} The Douala/Kribi-Campo basin is an Atlantic passive-margin basin whose sedimentary fill (Aptian Cretaceous to Recent) records rifting and drifting during the opening of the South Atlantic.
\item \textbf{False.} The Pan-African orogeny affected the \emph{whole} mobile belt between the Congo Craton and the West African Craton — i.e. the Adamawa-Yadé, Yaoundé and Poli domains across central and northern Cameroon, not only the north.
\item \textbf{True.} Lake Nyos occupies a maar (phreatomagmatic crater) on the Oku volcanic field; the August 1986 disaster was a limnic eruption: sudden overturn of the stratified lake released a dense cloud of dissolved magmatic \ce{CO2} that killed about 1700 people.
\item \textbf{True.} The Minim-Martap and Ngaoundal plateaux bauxites are residual lateritic deposits formed by intense tropical weathering (leaching of silica, concentration of \ce{Al2O3}) of aluminous Precambrian basement (granites, gneisses) under a hot humid climate.
\end{enumerate}
\end{corrige}

\begin{corrige}[Exercise 3 — Definitions]
\begin{enumerate}
\item \textbf{Limnic eruption} (lake overturn): a rare natural disaster in which dissolved gas (usually \ce{CO2}) accumulated in the deep layers of a stratified lake is suddenly exsolved and erupted, producing a dense gas cloud that flows downhill and asphyxiates life. \emph{Cameroonian example:} Lake Nyos, 21 August 1986 (and Lake Monoun, 1984).
\item \textbf{Shear zone (fault zone)}: a tabular zone of intense ductile or brittle deformation where rocks are displaced relative to one another, often marked by mylonites, cataclasites and fault breccias. \emph{Cameroonian example:} the Sanaga Fault (Central Cameroon Shear Zone), a NE–SW Pan-African shear zone reactivated in the Cretaceous.
\item \textbf{Lateritic bauxite}: a residual aluminium ore (gibbsite, boehmite) formed by prolonged tropical chemical weathering that removes silica and bases and concentrates aluminium hydroxides in a laterite profile. \emph{Cameroonian example:} Minim-Martap and Ngaoundal deposits (Adamawa Plateau).
\item \textbf{Angular unconformity}: a surface of erosion separating tilted or folded older strata below from younger, less deformed strata deposited horizontally above; it records deformation, uplift, erosion, then renewed subsidence. \emph{Cameroonian example:} the unconformity between folded Pan-African basement and the flat-lying Cretaceous sandstones of the Mamfe Basin.
\item \textbf{Pan-African granitoids}: granitic intrusions (syn-, late- and post-tectonic granites, granodiorites, diorites) emplaced during the Pan-African orogeny ($\sim 600$–$500$ Ma) into the mobile belt between the Congo and West African cratons. \emph{Cameroonian example:} the post-tectonic granites of the Adamawa-Yadé domain (e.g. around Ngaoundéré and Bertoua).
\end{enumerate}
\end{corrige}

\begin{corrige}[Exercise 4 — Direct Calculation: Lake Nyos \ce{CO2} Profile]
\textbf{1. Profile.} Depth is plotted on the vertical axis increasing downwards, concentration on the horizontal axis. The curve shows a strongly stratified lake: near-zero \ce{CO2} in the upper 50 m, then an exponential-like increase below 100 m, reaching $52.7\ \mathrm{mmol\,kg^{-1}}$ at the bottom — the signature of magmatic \ce{CO2} charging from below.

\begin{popfigure}
\begin{center}
\begin{tikzpicture}[xscale=0.09, yscale=0.028]
\draw[->, thick] (0,0) -- (60,0) node[below left, font=\small] {\ce{CO2} (mmol/kg)};
\draw[->, thick] (0,0) -- (0,215) node[above left, font=\small, align=center] {depth (m)};
\foreach \d in {0,50,100,150,200} \draw (0.6,\d) -- (-0.6,\d) node[left, font=\scriptsize] {\d};
\foreach \c in {0,10,20,30,40,50} \draw (\c,3) -- (\c,-3) node[below, font=\scriptsize] {\c};
\draw[very thick, popblue, mark=*, mark size=1.5pt] plot coordinates {(0.2,0) (0.5,50) (2.1,100) (15.8,150) (45.3,200) (52.7,207)};
\draw[dashed, popmuted] (15.8,0) -- (15.8,150);
\node[popblue, font=\scriptsize] at (35,120) {strong stratification};
\end{tikzpicture}
\end{center}
\legende{Dissolved \ce{CO2} profile of Lake Nyos (March 2023 data): concentration rises sharply below 100 m.}
\end{popfigure}

\textbf{2. Mass of \ce{CO2} in the 150–207 m layer.}

\begin{methode}[Trapezoidal integration for a stratified layer]
\begin{enumerate}
\item Identify the concentration at the top and bottom of the layer from the profile.
\item Compute the average concentration (trapezoidal rule).
\item Calculate the layer volume: $V = A \times \Delta h$.
\item Convert water volume to mass using density $\rho \approx 1000\ \mathrm{kg\,m^{-3}}$.
\item Convert average concentration from $\mathrm{mmol\,kg^{-1}}$ to $\mathrm{mol\,kg^{-1}}$ ($\times 10^{-3}$).
\item Multiply: $n = \bar{c} \times m_w$.
\item Convert moles to mass: $m = n \times M_{\ce{CO2}}$ ($M = 44\ \mathrm{g\,mol^{-1}}$).
\item Convert grams to tonnes ($1\ \mathrm{t} = 10^6\ \mathrm{g}$).
\end{enumerate}
\end{methode}

Average concentration (trapezoidal rule):
\[ \bar{c} = \frac{c_{150} + c_{207}}{2} = \frac{15.8 + 52.7}{2} = 34.25\ \mathrm{mmol\,kg^{-1}}. \]

Volume of the layer:
\[ V = A \times h = 1.2\times 10^{5}\ \mathrm{m^2} \times (207-150)\ \mathrm{m} = 1.2\times 10^{5} \times 57 = 6.84\times 10^{6}\ \mathrm{m^3}. \]

Mass of water:
\[ m_{w} = \rho V = 1000 \times 6.84\times 10^{6} = 6.84\times 10^{9}\ \mathrm{kg}. \]

Moles of \ce{CO2}:
\[ n = \bar{c}\, m_{w} = 34.25\times 10^{-3}\ \mathrm{mol\,kg^{-1}} \times 6.84\times 10^{9}\ \mathrm{kg} = 2.343\times 10^{8}\ \mathrm{mol}. \]

Mass of \ce{CO2}:
\[ m_{\ce{CO2}} = n \times M = 2.343\times 10^{8}\ \mathrm{mol} \times 44\ \mathrm{g\,mol^{-1}} = 1.031\times 10^{10}\ \mathrm{g} \approx 1.03\times 10^{4}\ \mathrm{tonnes}. \]

\begin{aretenir}
The bottom 50 m of Lake Nyos alone stores roughly \textbf{10 000 tonnes of dissolved \ce{CO2}} — which is why continuous artificial degassing is indispensable.
\end{aretenir}

\begin{attention}
Common mistakes: forgetting to convert mmol to mol ($\times 10^{-3}$), or grams to tonnes ($1\ \mathrm{t} = 10^{6}\ \mathrm{g}$). Keep track of units at every step.
\end{attention}
\end{corrige}

\begin{corrige}[Exercise 5 — Structured Question: Cameroon Volcanic Line]
\textbf{(a) Geographical extent [4 marks].} The CVL is a $\sim$\SI{1600}{km} Y-shaped alignment of Cenozoic volcanoes straddling the continent–ocean boundary. \emph{Oceanic sector} (Gulf of Guinea): Pagalu (Annobón), São Tomé, Príncipe, Bioko. \emph{Continental sector}: Mount Cameroon, Rumpi Hills, Manengouba, Bamboutos, Oku, Ngaoundéré, Mandara Mountains, extending to the Biu Plateau (Nigeria). [1 mark for extent, 1 for sectors, 2 for six named volcanoes.]

\textbf{(b) Cross-section [4 marks].}

\begin{popfigure}
\begin{center}
\begin{tikzpicture}[scale=1.0]
\fill[poporange!70] (-4,0) -- (0,3.4) -- (4,0) -- cycle;
\fill[popgold!80] (-4,0) -- (-1.6,1.35) -- (-0.9,0) -- cycle;
\fill[popgold!80] (4,0) -- (1.7,1.3) -- (0.8,0) -- cycle;
\draw[very thick, popdark] (-4,0) -- (0,3.4) -- (4,0);
\draw[very thick, popdark] (-4,0) -- (4,0);
\draw[thick, poppurple, dashed] (0,3.4) -- (0,-1.2);
\draw[thick, poppurple, dashed] (2.6,0.6) -- (1.4,-1.2);
\fill[popgreen!60] (2.2,0) -- (2.9,0.85) -- (3.6,0) -- cycle;
\draw[thick, popdark] (2.2,0) -- (2.9,0.85) -- (3.6,0);
\draw[popink, ->] (0,3.4) -- (0,4.1);
\node[font=\scriptsize, align=center] at (0,4.5) {central vent\\(crater)};
\node[font=\scriptsize, align=center] at (-2.6,2.2) {alternating lava flows\\and pyroclastic layers};
\node[font=\scriptsize, align=center] at (4.6,1.5) {parasitic cone\\(flank vent)};
\node[font=\scriptsize] at (0,-0.7) {magma chamber / conduit};
\node[font=\scriptsize] at (-2.2,-0.4) {basement};
\end{tikzpicture}
\end{center}
\legende{Schematic cross-section of a CVL stratovolcano (e.g. Mount Cameroon).}
\end{popfigure}

[1 mark per correctly drawn and labelled element: layered cone, central vent, parasitic cone, magma conduit.]

\textbf{(c) Two hypotheses for the origin of the CVL [6 marks].}

\begin{center}
\begin{tabularx}{\linewidth}{|l|X|X|}
\hline
\rowcolor{popblueL} \textbf{Hypothesis} & \textbf{Supporting argument} & \textbf{Counter-argument} \\
\hline
Mantle plume / hotspot & Linear ocean-island chain; OIB-type alkaline basalts with deep-mantle geochemical signature; uplift and high heat flow. & No systematic age progression along the line (volcanism is coeval), unlike Hawaii; no clear plume conduit imaged beneath the whole line. \\
\hline
Lithospheric fracture / structural control & The line follows reactivated Pan-African shear zones (e.g. Sanaga, Foumban); volcanoes align on faults; decompression melting during rifting. & Does not easily explain the oceanic islands far from the faults, nor the large magma volumes and long-lived activity. \\
\hline
\end{tabularx}
\end{center}

[3 marks per hypothesis: 1 statement, 1 supporting argument, 1 counter-argument.]

\textbf{(d) Why Mount Cameroon is the most active [3 marks].} It is the only CVL volcano with frequent historical eruptions (about eight recorded since the 19th century, e.g. 1922, 1959, 1982, 1999, 2000), fed by a shallow, well-connected plumbing system producing fluid basaltic (basanite–hawaiite) lavas of low viscosity that erupt easily along fissures on the flanks. Its position at the continent–ocean transition, where fractures are actively reactivated, maintains magma supply. [1 mark eruption history, 1 mark composition/viscosity, 1 mark structural setting.]
\end{corrige}

\begin{corrige}[Exercise 6 — Map Exercise: Geological Features of Cameroon]
\textbf{Locations (8 marks, 1 each):}
\begin{itemize}
\item \textbf{Ntem Complex}: far south, around Ebolowa–Sangmélima–Ambam, bordering Gabon/Equatorial Guinea (Congo Craton margin).
\item \textbf{Douala Basin (Kribi-Campo)}: coastal strip from Douala southwards through Kribi to Campo.
\item \textbf{Mamfe Basin}: Southwest Region, around Mamfe, along the Cross River near the Nigerian border.
\item \textbf{Sanaga Fault}: NE–SW lineament crossing central Cameroon, roughly along the Sanaga River between the Yaoundé region and the coast.
\item \textbf{Lake Nyos}: Northwest Region, Oku volcanic field, NE of Wum.
\item \textbf{Mbalam iron deposit}: East Region, near the Congo border, SE of Abong-Mbang.
\item \textbf{Mount Cameroon}: Southwest coast, near Buea/Limbe — highest peak in West and Central Africa ($\sim$\SI{4095}{m}).
\item \textbf{Bamboutos Mountains}: West Region, east of Mbouda, on the CVL.
\end{itemize}

\textbf{Economic importance (2 features, 2 marks each):}

\emph{Mbalam iron deposit:} hosted in Archaean–Palaeoproterozoic Banded Iron Formation of the Nyong Group, it is one of the largest undeveloped high-grade iron ore resources in Central Africa. Its exploitation (with the associated deep-sea port at Kribi and a trans-Cameroon railway) would diversify the economy, create thousands of jobs and generate major export revenue, reducing dependence on oil.

\emph{Douala Basin:} it is Cameroon's main petroleum-bearing basin (oil and gas fields such as those of the Rio del Rey/Douala offshore), supplying the refinery and the energy sector. Its Cretaceous–Cenozoic source and reservoir rocks, formed during the South Atlantic opening, also yield construction materials (sand, limestone) and support coastal agriculture and port activities.

[Examiner expectation: accurate placement on the map (SW, S, Centre, NW, E regions), and justification linking \emph{geological setting} to \emph{economic value}, not just ``it brings money''.]
\end{corrige}

\begin{corrige}[Exercise 7 — Data Analysis: Lake Nyos Degassing]
\textbf{(a) Principle of the self-driven pipe [4 marks].} The pipe inlet sits in deep, \ce{CO2}-rich water. As the water rises inside the pipe, hydrostatic pressure decreases; dissolved \ce{CO2} exsolves and forms bubbles. The water–gas mixture inside the pipe is therefore \emph{less dense} than the surrounding lake water at the same depth. The density difference creates buoyancy (gas-lift effect): the column of light mixture is pushed upward by the heavier surrounding water, producing a self-sustaining fountain at the surface — no pump or external energy is needed once the flow is primed. Degassed water falls back at the surface, harmless.

\textbf{(b) Flow rates [3 marks].}

\begin{methode}[Pipe flow calculation]
\begin{enumerate}
\item Cross-sectional area: $A = \pi r^2$.
\item Volumetric flow rate: $Q = A v$.
\item Mass flow rate: $\dot{m} = \rho Q$ (use mixture density).
\end{enumerate}
\end{methode}

Cross-sectional area of the pipe:
\[ A = \pi r^2 = \pi (0.07)^2 = \pi \times 4.9\times 10^{-3} \approx 1.539\times 10^{-2}\ \mathrm{m^2}. \]

Volumetric flow rate:
\[ Q = A\,v = 1.539\times 10^{-2} \times 2.5 \approx 3.85\times 10^{-2}\ \mathrm{m^3\,s^{-1}} \;(\approx 0.038\ \mathrm{m^3\,s^{-1}}). \]

Mass flow rate:
\[ \dot{m} = \rho Q = 900 \times 3.85\times 10^{-2} \approx 34.6\ \mathrm{kg\,s^{-1}}. \]

\textbf{(c) Mass of \ce{CO2} removed per day [3 marks].}

Mass of water lifted per day:
\[ m_{w} = 34.6\ \mathrm{kg\,s^{-1}} \times 86400\ \mathrm{s} \approx 2.99\times 10^{6}\ \mathrm{kg\,day^{-1}}. \]

Moles of \ce{CO2} per day:
\[ n = 50\times 10^{-3}\ \mathrm{mol\,kg^{-1}} \times 2.99\times 10^{6}\ \mathrm{kg} \approx 1.49\times 10^{5}\ \mathrm{mol\,day^{-1}}. \]

Mass of \ce{CO2}:
\[ m = n \times 44\ \mathrm{g\,mol^{-1}} \approx 6.57\times 10^{6}\ \mathrm{g} \approx 6.6\ \mathrm{tonnes\,day^{-1}}. \]

\textbf{(d) Why not at the very bottom [2 marks].} (i) Drawing from the very bottom would disturb and erode the deepest, most gas-charged layer too fast, risking triggering the very overturn the pipe is meant to prevent; (ii) sediment and corrosive, gas-saturated bottom water could clog or damage the inlet. Pumping from 60–70 m removes gas-rich water progressively while preserving the stability of the stratification below.
\end{corrige}

\begin{corrige}[Exercise 8 — Stratigraphic Column Construction]
\begin{center}
\begin{tabularx}{\linewidth}{|c|X|l|X|X|}
\hline
\rowcolor{popblueL} \textbf{Position} & \textbf{Unit} & \textbf{Age} & \textbf{Main rock types} & \textbf{Tectonic setting} \\
\hline
7 (youngest) & Mount Cameroon volcanics & Quaternary & Basalts, basanites, pyroclastics & CVL intraplate rift magmatism \\
\hline
6 & Bamboutos caldera volcanics & Cenozoic (Miocene–Pliocene) & Trachytes, phonolites, ignimbrites, basalts & CVL caldera volcanism \\
\hline
5 & Douala Basin rift volcanics & Jurassic–Cretaceous & Tholeiitic basalts & South Atlantic rifting \\
\hline
4 & Mamfe Basin fill & Cretaceous & Sandstones, shales, conglomerates & Intracontinental rift (Benue trough arm) \\
\hline
3 & Pan-African granites & Neoproterozoic ($\sim$550 Ma) & Granites, granodiorites, diorites & Pan-African collisional orogeny \\
\hline
2 & Nyong Group & Palaeoproterozoic & Schists, quartzites, BIF & Eburnean metasedimentary basin \\
\hline
1 (oldest) & Ntem Complex & Archaean & Gneisses, migmatites, charnockites & Congo Craton basement \\
\hline
\end{tabularx}
\end{center}

\textbf{Marking guide [14 marks]:} 7 marks for the correct order (1 per correctly placed unit, deduct for inversions); 3 marks for ages; 2 marks for rock types; 2 marks for tectonic settings.

\begin{attention}
Frequent error: placing the Pan-African granites \emph{below} the Nyong Group. Remember: the Nyong metasediments (Palaeoproterozoic) were deposited on and around the Archaean craton and were then intruded and reworked by the much younger Pan-African granitoids ($\sim$550 Ma).
\end{attention}
\end{corrige}

\begin{corrige}[Exercise 9 — Essay: Geology of Cameroon — Blessing and Threat]
\textbf{Indicative mark scheme and model content [25 marks].}

\textbf{Introduction (2 marks):} define the dual nature — Cameroon's varied geology (craton, Pan-African belt, basins, CVL) endows it with resources but also exposes it to hazards.

\textbf{Mineral resources (10 marks) — expect at least three, with setting, location, potential:}
\begin{itemize}
\item \emph{Bauxite} (Minim-Martap, Ngaoundal, Adamawa): lateritic weathering of Precambrian basement; among the largest untapped reserves in Africa.
\item \emph{Iron ore} (Mbalam–Nabeba, East Region): BIF of the Nyong Group; high-grade, world-class deposit awaiting infrastructure.
\item \emph{Hydrocarbons} (Douala/Rio del Rey and Logone Birni basins): Cretaceous–Cenozoic passive-margin and rift fills; backbone of export revenue.
\item \emph{Gold and diamonds} (East and Adamaoua, e.g. Bétaré Oya): Pan-African shear-zone-hosted and alluvial; artisanal and industrial potential.
\item \emph{Limestone} (Figuil, Douala Basin): cement industry; \emph{cobalt/nickel} (Lomié laterites).
\end{itemize}
[Marking: 1 mark per resource correctly located + setting; up to 6; 4 marks for discussion of economic potential and limits (infrastructure, artisanal vs industrial, price volatility).]

\textbf{Geologic hazards (10 marks) — expect at least three, with cause, example, mitigation:}
\begin{itemize}
\item \emph{Volcanic eruptions}: Mount Cameroon (1922, 1959, 1982, 1999, 2000) — fluid basaltic lava flows threatening plantations, roads and Buea/Limbe; mitigation: monitoring (seismic, gas), land-use zoning, evacuation plans.
\item \emph{Limnic eruptions}: Lake Nyos (1986, $\sim$1700 deaths) and Lake Monoun (1984) — magmatic \ce{CO2} accumulation in stratified maar lakes; mitigation: degassing pipes (2001, 2011), relocation of shoreline villages, lake-level control.
\item \emph{Landslides}: Bamboutos, Bafoussam (2019 disaster), Bamenda escarpment — steep volcanic slopes, deep lateritic weathering, intense rainfall, deforestation; mitigation: reforestation, drainage, slope engineering, building codes.
\item \emph{Earthquakes}: moderate seismicity along the Sanaga Fault and CVL; \emph{flooding/erosion}: Logone plain, Douala.
\end{itemize}
[Marking: 2 marks per hazard (cause 1, example/mitigation 1) up to 6; 4 marks for quality of discussion.]

\textbf{Sustainable management (5 marks):} geological mapping (BRGM-style 1:200 000 sheets) to delineate resources and hazard zones; continuous monitoring networks (volcano observatories, lake sensors, seismic stations); land-use planning that excludes construction on unstable slopes and floodplains; environmental impact assessment for mines; valorisation of artisanal mining into regulated channels; education of populations.

\textbf{Conclusion (integrated in marks above):} geology is a blessing only if knowledge guides exploitation and risk reduction.

\begin{attention}
Examiner expectations: precise Cameroonian localities, correct genetic settings, and balanced treatment of both sides. A purely descriptive list without linking geology to economy/risk caps the mark at about 60\%.
\end{attention}
\end{corrige}

\begin{corrige}[Exercise 10 — Problem Solving: Village Relocation near Bamboutos]
\textbf{(a) Risk assessment and ranking [12 marks].}

\emph{Site A — weathered basalt, slope $12^\circ$:} thick laterite ($>15$ m) is permeable and can lose strength when saturated, but the gentle slope ($12^\circ$) is below the typical failure threshold, there is no evidence of past instability, and the site is \SI{800}{m} from the river (no flood risk). \textbf{Landslide risk: low–moderate; flood risk: negligible.}

\emph{Site B — fractured trachyte with clay gouge, slope $28^\circ$:} steep slope close to the angle of repose for weathered volcanic material; clay gouge along fractures acts as a lubricated slip surface; tension cracks and tilted trees prove \emph{active} movement; high convective rainfall will infiltrate fractures and raise pore pressure. Proximity to the river (150 m) adds undercutting/flood risk. \textbf{Landslide risk: very high — the most hazardous site.}

\emph{Site C — alluvial terrace over impermeable shale, slope $3^\circ$:} flat, so landslide risk is negligible; however it lies on a floodplain \SI{50}{m} from the river, with old flood deposits and waterlogging — clear evidence of recurrent flooding; the impermeable shale prevents drainage, aggravating waterlogging and foundation problems. \textbf{Landslide risk: very low; flood risk: high.}

\textbf{Ranking, safest to most hazardous:} Site A (safest) $>$ Site C (flood-prone but manageable) $>$ Site B (actively failing — unacceptable).

\textbf{(b) Recommendation [5 marks].} Choose \textbf{Site A}. Mitigation measures (any two):
\begin{enumerate}
\item Install a surface and subsurface drainage network (cut-off drains upslope, lined channels) to prevent saturation of the laterite during convective storms.
\item Reforest the upper slopes with deep-rooted species and ban cultivation on the steepest sectors to increase cohesion and reduce runoff.
\item (Alternative) Terrace construction platforms and enforce building codes (light structures, foundations anchored below the weathered layer).
\end{enumerate}

\textbf{(c) Influence of volcanic history on slope stability [5 marks].} The Bamboutos are a caldera complex: (i) caldera collapse produced thick, poorly consolidated pyroclastic and ignimbrite deposits that are weak and easily eroded; (ii) alternating lava and ash layers create permeability contrasts — water perches on impermeable horizons, building pore pressure and promoting sliding along them; (iii) hydrothermal alteration around the former magma system converted feldspars and glass to soft clays, weakening the rock mass; (iv) steep caldera walls and fault scarps inherited from collapse provide ready-made steep slopes and fracture planes. Combined with deep tropical weathering (laterite) and intense rainfall, this legacy explains the chronic landslide susceptibility of the massif.
\end{corrige}

\begin{corrige}[Exercise 11 — Integration: Stratigraphy, Structure and Resources]
\textbf{(a) Tectonostratigraphic domains, south to north [8 marks].}
\begin{enumerate}
\item \textbf{Congo Craton / Ntem Complex (south)}: Archaean ($\sim 2.9$–$2.5$ Ga) gneisses, migmatites, charnockites, TTG granitoids and greenstone relics — stable cratonic basement.
\item \textbf{Nyong Group (craton margin)}: Palaeoproterozoic ($\sim 2.1$ Ga) metasediments — schists, quartzites, banded iron formations — deposited on the craton edge and metamorphosed during the Eburnean orogeny.
\item \textbf{Pan-African belt (Yaoundé and Adamawa-Yadé domains, centre)}: Neoproterozoic ($\sim 800$–$500$ Ma) metasediments and metavolcanics (schists, gneisses, amphibolites) thrust southwards onto the craton, intruded by syn- to post-tectonic granitoids ($\sim 600$–$500$ Ma).
\item \textbf{Sanaga Fault zone}: NE–SW Pan-African ductile shear zone (mylonites) reactivated later.
\item \textbf{Douala Basin (west/coast)}: Cretaceous to Recent passive-margin sediments (sandstones, shales, limestones, evaporites) with Jurassic–Cretaceous rift basalts.
\end{enumerate}
[2 marks per well-described domain, any four.]

\textbf{(b) Role of the Sanaga Fault [6 marks].}
\begin{enumerate}
\item \emph{Pan-African orogeny}: it was a major transcurrent (strike-slip) shear zone accommodating collision between the Congo Craton and the northern blocks, localising deformation, metamorphism and granite emplacement ($\sim 600$–$500$ Ma).
\item \emph{Cretaceous Atlantic opening}: the old zone of weakness was \emph{reactivated} as a transfer/rift fault during South Atlantic rifting, controlling the geometry, subsidence and sedimentation of the Douala Basin (and the Benue trough system inland).
\item \emph{Cenozoic CVL magmatism}: reactivated fractures provided pathways for magma ascent; the alignment of volcanic centres (e.g. Mount Cameroon at the intersection of the CVL with coastal fracture systems) reflects this structural inheritance.
\end{enumerate}
[2 marks each: 1 for the mechanism, 1 for the consequence.]

\textbf{(c) Lom Basin vs Mamfe Basin [6 marks].}

\begin{center}
\begin{tabularx}{\linewidth}{|l|X|X|}
\hline
\rowcolor{popblueL} \textbf{Criterion} & \textbf{Mamfe Basin (SW)} & \textbf{Lom Basin (East)} \\
\hline
Setting & Cretaceous rift, arm of the Benue trough & Cretaceous intracratonic rift (Central African rift system) \\
\hline
Stratigraphy & Basal conglomerates, fluvial sandstones, shales; marine influence at times & Fluvio-lacustrine sandstones, siltstones, shales \\
\hline
Environment & Fluvial to deltaic, locally marginal marine & Continental fluvial and lacustrine \\
\hline
Hydrocarbon potential & Proven shows; source shales, reservoir sandstones; explored for oil/gas & Speculative; potential source/reservoir pairs but poorly explored \\
\hline
\end{tabularx}
\end{center}
[6 marks: 2 for stratigraphy comparison, 2 for environments, 2 for hydrocarbon potential.]

\textbf{(d) Genetic model of the Mbalam BIF and economic viability [5 marks].} \emph{Genesis:} during the Archaean–Palaeoproterozoic, the oceans were rich in dissolved ferrous iron (\ce{Fe^{2+}}) from hydrothermal and weathering inputs. In a shallow marine basin on the Congo Craton margin (Nyong Group), periodic oxidation — linked to the rise of photosynthetic oxygen (Great Oxidation Event context) — precipitated iron oxides (magnetite, hematite) in alternating bands with silica (chert), forming the Banded Iron Formation. Subsequent metamorphism recrystallised and locally upgraded the ore. \emph{Economic viability:} (i) very large tonnage with high Fe grade (locally $> 50\%$ Fe, enriched itabirite); (ii) rising global iron demand and prices; (iii) planned infrastructure — a railway to the deep-water port of Kribi — solves the historic transport bottleneck; (iv) possibility of direct-shipping ore with limited processing. These factors transform a long-known geological resource into a bankable reserve.
\end{corrige}

\newpage

\coursec{Summary sheet}

\begin{popfigure}
\begin{tikzpicture}[
  mindmap,
  grow cyclic,
  every node/.style={concept, font=\small},
  root concept/.append style={concept color=popink, font=\small\bfseries, text=white, minimum size=2.6cm},
  level 1 concept/.append style={level distance=3.6cm, sibling angle=51.4, font=\scriptsize, minimum size=1.9cm},
  level 2 concept/.append style={level distance=2.2cm, font=\tiny, minimum size=1.3cm},
]
\node[root concept]{Geology of Cameroon}
  child[concept color=poporange]{ node {Igneous \& CVL}
    child { node {Mt Cameroon} }
    child { node {Ring complexes} }
  }
  child[concept color=popgreen]{ node {Metamorphics}
    child { node {Craton gneiss} }
    child { node {Pan-African belt} }
  }
  child[concept color=popblue]{ node {Sedimentary basins}
    child { node {Douala} }
    child { node {Benue trough} }
    child { node {Chad basin} }
  }
  child[concept color=poppurple]{ node {Limnology}
    child { node {Lake Nyos} }
    child { node {Crater lakes} }
  }
  child[concept color=poppink]{ node {Structures}
    child { node {Faults, folds} }
    child { node {CC shear zone} }
  }
  child[concept color=popgold]{ node {Resources}
    child { node {Bauxite, oil} }
    child { node {Gold, iron} }
  }
  child[concept color=popmuted]{ node {Stratigraphy \& hazards}
    child { node {Column Precambrian--Recent} }
    child { node {Eruptions, limnic} }
  };
\end{tikzpicture}
\legende{Mind map of the chapter ``The Geology of Cameroon'': seven branches linking rock units, basins, lakes, structures, resources, stratigraphy and hazards.}
\end{popfigure}

\begin{bilan}
\textbf{1. Key definitions to master}
\begin{itemize}
\item \cle{Cameroon Volcanic Line (CVL)}: a $\sim\SI{1600}{km}$ Y-shaped alignment of Cenozoic volcanoes and plutonic ring complexes running from the Gulf of Guinea islands (Annob\'on, S\~ao Tom\'e, Pr\'incipe, Bioko) through Mount Cameroon, the Bambouto--Bamenda highlands to Lake Chad; both oceanic and continental segments.
\item \cle{Ring complex}: anorogenic plutonic association (granites, syenites, gabbros) emplaced along ring faults, typical of the continental CVL (e.g.\ the granitic massifs of the Adamaoua--Bamenda axis).
\item \cle{Craton}: old, stable, rigid part of continental crust; in Cameroon the northern edge of the \cle{Congo craton} (Archean--Palaeoproterozoic gneisses, charnockites, greenstone belts of the Ntem complex) occupies the south.
\item \cle{Pan-African orogeny}: Neoproterozoic ($\sim$750--550~Ma) mountain-building event that produced the metamorphic belt (schists, gneisses, migmatites, granitoids) of central and northern Cameroon.
\item \cle{Sedimentary basin}: subsiding region filled with sediments; Cameroon has coastal basins (Douala/Kribi-Campo, Rio del Rey), intracratonic basins (Chad, Mb\'er\'e--Djerem) and the Garoua basin linked to the \cle{Benue trough}.
\item \cle{Limnology}: study of inland waters; \cle{meromictic lake}: permanently stratified lake whose deep water (monimolimnion) does not mix with surface water -- the dangerous case of Lakes Nyos and Monoun.
\item \cle{Limnic eruption}: sudden overturn of a CO$_2$-charged deep lake releasing a dense gas cloud (Nyos 1986, Monoun 1984).
\end{itemize}

\textbf{2. Facts, ages and ``laws'' to know by heart}
\begin{itemize}
\item Basement: Archean Ntem complex (south) $\rightarrow$ Palaeoproterozoic (Eburnean) $\rightarrow$ Neoproterozoic Pan-African belt (centre/north) $\rightarrow$ Cretaceous--Recent sedimentary cover $\rightarrow$ Cenozoic CVL volcanism.
\item Mount Cameroon ($\sim\SI{4095}{m}$): active basaltic stratovolcano; recent eruptions 1922, 1959, 1982, 1999, 2000 -- the most active volcano of the CVL.
\item CVL magmas: alkaline, intraplate (basanites, basalts $\rightarrow$ trachytes, phonolites); origin debated (hot line / reactivated lithospheric fractures), \emph{not} a classic subduction zone.
\item Stratigraphic order of basins: Benue--Garoua (Albian--Cenomanian marine transgression), Douala basin (Cretaceous--Tertiary, oil-bearing), Chad basin (Cretaceous--Recent, ``Continental Terminal'').
\item Supergene enrichment: tropical weathering concentrates \cle{bauxite} (Minim-Martap, Ngaoundal) and lateritic iron; residual vs.\ transported deposits.
\end{itemize}

\textbf{3. Methods to apply}
\begin{itemize}
\item Reading a geological map of Cameroon: identify units by colour/age legend, locate craton, Pan-African belt, basins, CVL; draw a labelled cross-section.
\item Building a stratigraphic column: oldest at bottom, indicate lithology, thickness, age, unconformities.
\item Interpreting lake stability: thermal stratification + dissolved gas $\Rightarrow$ hazard assessment; degassing columns as mitigation.
\item Linking structure to resources: faults control ring complexes, mineralisation and basin geometry.
\end{itemize}

\textbf{4. Common mistakes to avoid}
\begin{itemize}
\item Confusing the CVL (Cenozoic, alkaline, intraplate) with a subduction volcanic arc.
\item Placing the Benue trough entirely in Cameroon: only the Garoua--Yola arm concerns Cameroon.
\item Saying Lake Nyos ``exploded'': it was a gas \emph{release} (overturn), not a volcanic eruption.
\item Mixing up metamorphic grades or attributing Pan-African ages to the Congo craton.
\item Forgetting that bauxite is a \emph{weathering} product, not a magmatic ore.
\end{itemize}

\textbf{5. GCE tips}
\begin{itemize}
\item Always support answers with a \textbf{labelled sketch} (map, cross-section, column) -- marks are reserved for diagrams.
\item Quote at least one \textbf{Cameroonian example} per theme (Mt Cameroon, Nyos, Minim-Martap, Douala basin).
\item For hazard questions, structure: cause $\rightarrow$ mechanism $\rightarrow$ effects $\rightarrow$ prevention (e.g.\ Nyos degassing pipes).
\item Learn the chronological table craton $\rightarrow$ Pan-African $\rightarrow$ basins $\rightarrow$ CVL; it organises the whole chapter.
\end{itemize}
\end{bilan}

\findecours

\end{document}
